% FUNCTIONAL ANALYSIS AND DIFFERENTIAL EQUATIONS                               %
%------------------------------------------------------------------------------%
% File  : Functional_Analysis_and_Differential_Equations.tex                   %
% Author: Klaus Hoeflinger, Beethovenstrasse 11, D-73117 Wangen                %
% Email : khoeflinger@t-online.de                                              %
%------------------------------------------------------------------------------%

%------------------------------------------------------------------------------%
%                       DOCUMENT CLASS and INCLUDE FILES                       %
%------------------------------------------------------------------------------%
\documentclass[11pt,twoside]{article}
\usepackage{geometry}
\geometry{paperheight=241mm,
          paperwidth=152mm,
          top=22mm,
          bottom=17mm,
          left=20mm,
          right=20mm,
          textwidth=112mm,
          textheight=202mm,
          headheight=10mm,
          headsep=3mm,
          footskip=9mm,
          includehead,
          includefoot,
          marginparsep=0mm,
          marginparwidth=0mm}

\renewcommand{\baselinestretch}{0.945}\normalsize

\AtBeginDocument{\setlength\abovedisplayskip{2mm}}
\AtBeginDocument{\setlength\belowdisplayskip{2mm}}

%------------------------------------------------------------------------------%
% define title formats                                                         %
%------------------------------------------------------------------------------%
\usepackage{titlesec}

\renewcommand{\thesection}{\arabic{section}}
\renewcommand{\sfdefault}{FiraSans}
\renewcommand{\ttdefault}{pcr}
\renewcommand{\rmdefault}{antiqua}

\titleformat{\part}[display]
{\normalfont \large \rmfamily \bfseries \centering}
{\small{\thepart.}}{2pt}{}

\titleformat{\section}[runin]
{\normalfont \rmfamily \bfseries}
{\thesection}{1pt}{.\;}[.]

%\titleformat{\section}
%{\normalfont \bfseries \small \filcenter}
%{\thesection.\;}{0mm}{}

\titleformat{\subsection}[runin]
{\normalfont \itshape}
{}{0mm}{}

\titleformat{\subsubsection}[runin]
{\normalfont \itshape}
{}{}{}{}

\titleformat{\paragraph}[runin]
{\normalfont}
{}{}{}{}

\titleformat{\subparagraph}[runin]
{\normalfont}
{}{}{}{}

\titlespacing*{\part}          { 0pt}{ 0pt}{ 8pt}
\titlespacing*{\section}       { 0pt}{ 7pt}{ 4pt}
\titlespacing*{\subsection}    { 0pt}{14pt}{ 6pt}
\titlespacing*{\subsubsection} { 0pt}{ 6pt}{12pt}
\titlespacing*{\paragraph}     { 0pt}{ 6pt}{12pt}
\titlespacing*{\subparagraph}  { 0pt}{ 0pt}{12pt}

%------------------------------------------------------------------------------%
% define enumerate parameters                                                  %
%------------------------------------------------------------------------------%
\usepackage{enumitem}
\setlength{\parindent}{3pt}
\setlength{\parskip}  {0pt}

%\setlist{noitemsep}

\setlist[enumerate]{
         leftmargin=12pt,
         rightmargin=0pt,
         itemsep=1pt,
         itemindent=3pt,
         topsep=3pt,
         labelsep=4pt,
         partopsep=0pt,
         labelwidth=0pt,
         listparindent=0pt,
         parsep=0pt}

\setlist[itemize]{
         leftmargin=12pt,
         rightmargin=0pt,
         itemsep=0pt,
         itemindent=1pt,
         topsep=1pt,
         labelsep=4pt,
         partopsep=1pt,
         labelwidth=0pt,
         listparindent=0pt,
         parsep=0pt}

%------------------------------------------------------------------------------%
% define packages                                                              %
%------------------------------------------------------------------------------%
\usepackage{upref}
\usepackage{amssymb}
\usepackage{slantsc}
\usepackage{amsmath}
\usepackage{amsthm}
\usepackage{thmtools}
\usepackage{amscd}
\usepackage{verbatim}
\usepackage{latexsym}
\usepackage{multicol}
\usepackage{graphicx}
\usepackage{setspace}
\usepackage[ngerman,english]{babel}
\usepackage[protrusion=true,expansion=true]{microtype}
\usepackage{tikz}
\usetikzlibrary{arrows,arrows.meta,decorations.markings}
\usepackage{mathrsfs}
\usepackage{amsfonts}
\usepackage{datetime2}
\usepackage{textgreek}
\usepackage{hyperref}
\usepackage{pifont}
\usepackage{relsize}
%\usepackage{biblatex}
\relsize{-0.05}
\usepackage[skip=0mm]{parskip}
\usetikzlibrary{decorations.pathmorphing}

\usepackage{mlmodern}
\usepackage[T5,T1]{fontenc}
\usepackage{ragged2e}

%\usepackage{mathabx}
%\usepackage{times}

%\usepackage{fontspec}
%\setmainfont{Nimbus Roman No9 L}
%\usepackage[T1]{fontenc}

%------------------------------------------------------------------------------%
% define page layout                                                           %
%------------------------------------------------------------------------------%
\usepackage{fancyhdr}
\pagestyle{fancy}
\setlength{\headheight}{30.0pt}
\renewcommand{\headrulewidth}{0pt}

\AtBeginDocument{
  \addtolength\abovedisplayskip{-0.0\baselineskip}
  \addtolength\belowdisplayskip{-0.0\baselineskip}
}

\fancyhead{}
\fancyfoot{}
\addtolength{\headwidth}{0.02\marginparwidth}

\makeatletter
\let\ps@plain\ps@fancy
\makeatother

\renewcommand\sectionmark[1]
{
 \def\sectionname{#1}
 \markright{\thesection #1}
}

\makeatletter
\@addtoreset{section}{part}
\makeatother

%------------------------------------------------------------------------------%
% define footnote without number                                               %
%------------------------------------------------------------------------------%
\newcommand{\myfootnote}[1]{%
\renewcommand{\thefootnote}{}%
\footnotetext{#1}%
\renewcommand{\thefootnote}{\arabic{footnote}}%
}

%------------------------------------------------------------------------------%
% define numbering in equations                                                %
%------------------------------------------------------------------------------%
\numberwithin{equation}{section}

%------------------------------------------------------------------------------%
% define newtheorem                                                            %
%------------------------------------------------------------------------------%
\declaretheoremstyle[
headfont=\scshape, 
headindent = 0mm, %\parindent,
%postheadhook = {\hspace{0mm}\newline},
%postheadspace = \newline, 
spaceabove = 3mm, 
spacebelow = 3mm,
numberwithin=section,
name=Theorem]{khMATH}
\declaretheorem[style=khMATH]{theorem}

\declaretheoremstyle[
headfont=\bfseries, 
headindent = 0mm, %\parindent,
%postheadhook = {\hspace{0mm}\newline},
%postheadspace = \newline, 
spaceabove = 3mm, 
spacebelow = 3mm,
numberwithin=part,
name=Theorem]{khMATH1}
\declaretheorem[style=khMATH1]{theorem1}

\declaretheoremstyle[
headfont=\scshape, 
headindent = 0mm, %\parindent,
%postheadhook = {\hspace{0mm}\newline},
%postheadspace = \newline, 
spaceabove = 3mm, 
spacebelow = 3mm,
numberwithin=section,
name=Corollary]{khMATH1}
\declaretheorem[style=khMATH1]{corollary}

%            name         counter     label              numeration-mode
\theoremstyle{plain}
\newtheorem*{introduction}            {\sc Introduction}
\newtheorem {standard}                {}                 [section]
\newtheorem*{definition}              {\sc Definition}
\newtheorem{satz1}                    {\sc Satz}
\newtheorem{lemma}                    {\sc Lemma}
\newtheorem*{proposition}             {\sc Proposition}
%\newtheorem{corollary}               {\sc Corollary}
\newtheorem*{corollar}                {\sc Korollar}
\newtheorem*{remark}                  {\sc Remark}
\newtheorem*{remarks}                 {\sc Remarks}
\newtheorem*{example}                 {Example}
\newtheorem*{examples}                {Examples}
\newtheorem*{theo}                    {\sc Theorem}

%------------------------------------------------------------------------------%
% define tabular                                                               %
%------------------------------------------------------------------------------%
\usepackage{tabularx}
\newcolumntype{L}[1]{>{\raggedright\arraybackslash}p{#1}}
\newcolumntype{C}[1]{>{\centering\arraybackslash}  p{#1}} 
\newcolumntype{R}[1]{>{\raggedleft\arraybackslash} p{#1}} 

%------------------------------------------------------------------------------%
% define \sh, \zB                                                              %
%------------------------------------------------------------------------------%
\def\dsh{{d.\,h.}}
\def\ie{{i.\,e.}}
\def\zB{z.\,B.}
\renewcommand{\abstractname}{ }

%------------------------------------------------------------------------------%
% change default spacing for cases                                             %
%------------------------------------------------------------------------------%
\makeatletter
\renewcommand*\env@cases[1][1.6]{%
  \let\@ifnextchar\new@ifnextchar
  \left\lbrace
  \def\arraystretch{#1}%
  \array{@{}l@{\quad}l@{}}%
}
\makeatother

\newenvironment{rcases}
  {\left.\begin{aligned}}
  {\end{aligned}\right\rbrace}

%------------------------------------------------------------------------------%
% change default for \predisplaypenalty                                        %
%------------------------------------------------------------------------------%
\predisplaypenalty=990
\allowdisplaybreaks \numberwithin{equation}{section}

%------------------------------------------------------------------------------%
% general notations                                                            %
%------------------------------------------------------------------------------%
\def\*{\star}
\def\={\,=\,}
\def\+{\,+\,}
\def\xsubset{{\,\subset\,}}
\def\xtimes{{\,\times\,}}
\def\xeof{{\,\in\,}}
\def\eq{{\,=\,}}
\def\xto{{\,\to\,}}
\def\eqdef{\,:=\,}
\def\:{{\,:\,}}
\def\ele{{\,\in\,}}
\def\exists{\mbox{ exists }}
\def\and{\mbox{ and }}
\def\for{\mbox{ for }}
\def\fa{\mbox{ for all }}
\def\fe{\mbox{ for each }}
\def\qfa{\quad \mbox{ for all }}
\def\df{\,\dotfill}
\def\iff{if and only if }
\def\wrt{with respect to}

\makeatletter
\newcommand \Df {\leavevmode \cleaders \hb@xt@ .70em{\hss .\hss }\hfill \kern \z@}
\makeatother

\def\sql{\mathfrak{a}}               % sesquilinear form a
\def\DF{B}                           % Dirichlet form a
%------------------------------------------------------------------------------%
% number sets and special elements                                             %
%------------------------------------------------------------------------------%
\def\P{\mathbb{H}}                   % half plane of $\C$
\def\J{I}                            % interval I
\def\N{{\mathbb{N}}}                 % unsigned integers
\def\Q{{\mathbb{Q}}}                 % rational integers
\def\Z{{\mathbb{Z}}}                 % integers
\def\R{{\mathbb{R}}}                 % real numbers
\def\Rn{{\mathbb{R}}^n}              % real numbers in Rn
\def\Rd{{\mathbb{R}}^d}              % real numbers in Rn
\def\C{{\mathbb{C}}}                 % complex numbers
\def\F{{\mathbb{K}}}                 % arbitrary field
\def\Torus{{\mathbb{T}}}             % complex circle line
\def\T{\tau}                         % upper bound of interval (0,t)
\def\sign#1{\operatorname{sign}(#1)} % sign of a number

%------------------------------------------------------------------------------%
% set theory                                                                   %
%------------------------------------------------------------------------------%
\def\G{G}                            % domain $G$
\def\clG{\overline{\G}}              % closure of $G$
\def\bndG{\partial \G}               % boundary of $G$
\def\setA{\mathcal{A}}               % set A
\def\setB{\mathcal{B}}               % set B
\def\setC{\mathcal{C}}               % set C
\def\setM{M}                         % set M
\def\setN{N}                         % set N
\def\setS{\mathcal{S}}               % set S
\def\famF{\mathfrak{F}}              % family of sets
\def\indI{\mathcal{I}}               % index set I
\def\pot#1{\mathfrak{P}(#1)}         % power set

\def\relR{\mathbf{R}}                % relation R
\def\mapf{f}                         % mapping f
\def\mapg{g}                         % mapping g

\def\set#1#2{\{ \, #1 \,: \, #2 \, \}}
\def\tensor#1#2{#1 \otimes #2}
\def\Tens{\oplus}

\def\cal#1{{\mathcal{#1}}}
\def\aut#1{\operatorname{Aut}(#1)}

\def\cross#1#2{#1 {\,\times\,} #2}
\def\multicross#1#2{#1 \times  \ldots \times  #2}
\def\multitens#1#2 {#1 \otimes \ldots \otimes #2}

%------------------------------------------------------------------------------%
% group theory                                                                 %
%------------------------------------------------------------------------------%
\def\1{{\mathsf{1}}}                % unit element of a monoid
\def\0{{\mathsf{0}}}                % unit element of an Abelian monoid
\def\kernel#1{\mathfrak{N}(#1)}     % kernel of a homomorphism
\def\cokernel#1{\mathfrak{C}{(#1)}} % cokernel of a homomorphism
\def\Hom#1{\operatorname{Hom}(#1)}

%------------------------------------------------------------------------------%
% linear algebra                                                               %
%------------------------------------------------------------------------------%
\def\vecV{\mathit{V}}               % vector space V
\def\vecH{\mathit{H}}               % vector space H
\def\vecW{\mathit{W}}               % vector space W
\def\vecU{\mathit{U}}               % subspace U
\def\convC{\mathcal{C}}             % convex set C
\def\basH{\mathfrak{B}}             % Hamel base B

\def\LO#1{\mathscr{L}(#1) }         % algebra of linear transformations

\def\scalar#1#2{ \langle #1,#2 \rangle }
\def\trace#1{\operatorname{tr}(#1)}
\def\dim#1{{\operatorname{dim} \, #1 }}
\def\codim#1{{\operatorname{codim} \, #1 }}
\def\dim{\operatorname{dim}}
\def\codim{\operatorname{codim}}
\def\ind{\operatorname{ind}}

\def\genG{\cal{G}}

\def\algA{\cal{A}}
\def\algB{\cal{B}}
\def\algU{\cal{U}}

\def\idI{\cal{I}}

%------------------------------------------------------------------------------%
% topology                                                                     %
%------------------------------------------------------------------------------%
\def\compK{{K}}          % compact set C
\def\openO{\mathcal{O}}  % open set O
\def\openU{\mathcal{U}}  % open set U
\def\U{\mathcal{U}}      % neighbourhood U
\def\V{\mathcal{V}}      % neighbourhood V
\def\d{\mathit{d}}       % metric function d

\def\interior#1{\mathring{#1}}
\def\closure#1{{\overline{#1}}}

\def\topT{\mathcal{T}}   % topology T
\def\topX{X}             % topological space X
\def\topY{Y}             % topological space Y
\def\topZ{Z}             % topological space Z
\def\metrS{M}            % metric space M

\def\grpG{G}             % topological group G
\def\grpA{A}             % topological group A
\def\grpB{B}             % topological group B
\def\grpC{C}             % topological group C
\def\grpH{H}             % topological group H
\def\grpU{U}             % topological group U
\def\grpV{V}             % topological group V
\def\topG{G}             % topological group G
\def\topH{H}             % topological group H
\def\topU{U}             % topological subgroup U

\def\subS{\cal{S}}
\def\riR{R}              % ring R

%------------------------------------------------------------------------------%
% calculus                                                                     %
%------------------------------------------------------------------------------%
\def\supp#1{\operatorname{supp}(#1)}
\def\singsupp#1{\operatorname{sing \, supp}(#1)}

\def\re{\operatorname{Re}}
\def\im{\operatorname{Im}}
\def\grad{\operatorname{grad}}

%\def\abs#1 {\left|#1\right|}
\def\abs#1 {\lvert #1 \rvert}
%\def\abs#1 {\left|\,#1\,\right|}

%------------------------------------------------------------------------------%
% measure theory                                                               %
%------------------------------------------------------------------------------%
\def\salgA{\Sigma}               % sigma-algebra S
\def\salgB{\cal{B}}              % sigma-algebra B
\def\Borel#1{\cal{B}_{#1}}       % Borel-algebra 
\def\LB{\lambda}                 % Lebesgue-Borel measure

%------------------------------------------------------------------------------%
% function spaces                                                              %
%------------------------------------------------------------------------------%
\def\Lp{L^p(\G)}                 % spaces L_p
\def\Lq{L^q(\G)}                 % spaces L_q
\def\Lh{L^2(\G)}                 % spaces L_2
\def\Li{L^{\infty}(\G)}          % spaces L_infty
\def\Lc{L_{loc}^1(\G)}           % spaces L_1^{loc}
\def\Ck{C^k(\G)}                 % spaces C_k

\def\BV#1{BV(#1) }               % set of functions of bounded variation
\def\AC#1{AC(#1) }               % set of absolutely continuous functions

\def\MP#1{\mathscr{M}^+(#1) }    % set of positive measures
\def\MS#1{\mathscr{M}^-(#1) }    % vector space of finite signed measures
\def\MC#1{\mathscr{M}(#1) }      % vector space of complex measures
\def\MCR#1{\mathscr{M}_{r}(#1) } % vector space of regular complex measures

\def\SobW{W^{k,p}(\G)}                % Sobolev space W
\def\SobO{\mathring{W}^{k,p}(\G)}     % Sobolev space W
%\def\WN#1#2{\mathring{W}^{#1}(#2) }   % Sobolev space W
\def\WN#1#2{W_0^{#1}(#2) }   % Sobolev space W
%\def\HN#1#2{\mathring{H}^{#1}(#2) }   % Sobolev space H
\def\HN#1#2{H_0^{#1}(#2) }   % Sobolev space H
%------------------------------------------------------------------------------%
% functional analysis                                                          %
%------------------------------------------------------------------------------%
\def\snorm#1{\lVert #1 \rVert}
\newcommand{\norm}[1]{\left\lVert #1 \right\rVert}

\def\topV{V}                       % topological vector space V
\def\topW{W}                       % topological vector space W
\def\normS{X}                      % normed space

\def\banX{\mathit{X}}              % Banach space X
\def\banY{\mathit{Y}}              % Banach space Y
\def\banZ{\mathit{Z}}              % Banach space Z
\def\dbanX{\mathit{X^*}}           % dual space of Banach space X
\def\dbanY{\mathit{Y^*}}           % dual space of Banach space Y
\def\dbanZ{\mathit{Z^*}}           % dual space of Banach space Z

\def\banA{\mathit{A}}              % Banach algebra A
\def\banB{\mathit{B}}              % Banach algebra B
\def\banC{\mathit{C}}              % C*-algebra C


\def\domain#1{\mathfrak{D}(#1)}    % domain
\def\range#1{\mathfrak{R}(#1)}     % range

\def\isoJ{\mathfrak{J}}            % isometry J
\def\repA{{\cal{A}_{\sql}}}        % representation operator A
\def\linS{{S}}                     % linear operator S
\def\linG{{G}}                     % linear operator G
\def\linL{{L}}                     % linear operator L
\def\linT{{T}}                     % linear operator A
\def\linA{\mathit{A}}              % linear operator A
\def\linB{{B}}                     % linear operator B
\def\linM{{M}}                     % linear operator M
\def\linU{{U}}                     % linear operator U
\def\linI{{I}}                     % linear operator I
\def\A{\cal{A}}                    % linear differential operator A
\def\BDO{\cal{B}}                  % linear boundary differential operator B
\def\linU{{U}}                     % unitary operator U
\def\linP{{P}}                     % projection P
\def\compA{k}                      % compact operator k
\def\linFR{f}                      % finite rank operator f
\def\fredA{A}                      % Fredholm operator F

\def\HAM{\cal{H}}                  % Hamilton operator
\def\PA{\partial^{\alpha}}         % elementary differential operator
\def\DO{\cal{A}}                   % linear differential operator
\def\BC{\cal{B}}                   % boundary operator
\def\SLO{\cal{L}_{p,q}}            % Sturm-Liouville operator
\def\MO{\cal{L}_{M}}               % momentum operator

\def\HS#1{S(#1)           }        % algebra of Hilbert-Schmidt operators
\def\FR#1{F(#1)           }        % algebra of finite rank operators
\def\TC#1{N(#1)           }        % algebra of trace class operators
\def\BU#1{\mathscr{U}(#1) }        % algebra of unitary linear operators
\def\BO#1{\mathscr{L}(#1) }        % algebra of bounded linear operators
\def\CO#1{\mathscr{C}(#1) }        % algebra of compact linear operators
\def\CL#1{\mathscr{C}(#1) }        % algebra of compact linear operators
\def\FO#1{\Phi(#1) }               % set of Fredholm operators
\def\FK#1#2{\Phi_{#2}(#1) }        % set of Fredholm operators with index k

\def\hilH{\mathit{H}}              % Hilbert space H
\def\clU{\mathit{U}}               % closed subspace U

\def\orthO{\mathfrak{O}}           % orthonormal system O
\def\orthS{\mathfrak{S}}           % orthonormal system S
\def\basB{\mathfrak{B}}            % basis B of a Hilbert space

\def\GL#1 {\mathbf{GL}(#1)}        % linear group
\def\OG#1 {\mathbf{O}(#1)}         % linear group
\def\SOG#1 {\mathbf{SO}(#1)}       % linear group

%------------------------------------------------------------------------------%
% distribution theory                                                          %
%------------------------------------------------------------------------------%
\def\disF{F}                  % distribution F
\def\disG{G}                  % distribution G
\def\disH{H}                  % distribution H
\def\disU{U}                  % distribution U
\def\disT{T}                  % distribution T
\def\SF{C^{\infty}(\G)}       % space of smooth functions
\def\TF{C_c^{\infty}(\G)}     % space of compactly supported smooth functions
\def\TFRd{C_c^{\infty}(\R^d)} % space of test functions
\def\SRd{\mathscr{S}(\R^d)}   % Schwartz space
\def\SR1{\mathscr{S}(\R)}     % Schwartz space
\def\B1#1{C^1(#1) }           % space of continuously differentiable functions
\def\Bm#1{C^m(#1) }           % space of continuously differentiable functions
\def\Ck{C^k(\G)}
\def\TD{\mathscr{S}'(\R^d)}   % space of temperated distributions
\def\E{\cal{E}}               % space of distributions with compact support
\def\D{\mathscr{D}}           % D
\def\FT{\mathcal{F}}          % Fourier transformation F
\def\FT{\mathscr{F}}          % Fourier transformation F

%------------------------------------------------------------------------------%
% other definitions                                                            %
%------------------------------------------------------------------------------%
\def\Proof{\textsc{Proof}}
\def\FRECHET{Fr\'{e}chet}
\def\ARZELA{Arzel\'{a}}
\def\GARDING{G\r{a}rding}
\def\HOELDER{H\"older}
\def\SCHROEDINGER{Schr\"odinger}
\def\JOERGENS{J\"orgens}
\def\WUEST{W\"uest}
\def\KAEHLER{K\"aehler}
\def\HOEFLINGER{H\"oflinger}
\def\POINCARE{Poincar\'{e}}
\def\FA{Functional Analysis}

\def\Author   {K.\,{\HOEFLINGER}}
\def\AuthorUC {K.\,HOEFLINGER}
\def\Version  {1.0}
\def\Email    {khoeflinger@t-online.de}


%\renewcommand\thesection{\Roman{section}}

\def\Title {Functional Analysis and Differential Equations}

\titleformat{\section}
{\normalfont \bfseries \filcenter}
{\thesection.\;}{0mm}{}

\titleformat{\subsection}%[runin]
{\normalfont \itshape}
{}{}{}

\titlespacing*{\section}   {0pt}{12pt}{4pt}
\titlespacing*{\subsection}{0pt}{6pt}{0pt}

\def\G{G}

%------------------------------------------------------------------------------%
%                                BEGIN OF DOCUMENT                             %
%------------------------------------------------------------------------------%
\begin{document}
\thispagestyle{empty}

\centerline{\large{\textbf{\Title}}}

\vspace{4pt}
\centerline{\footnotesize{{\textsc{\Author}}}}
\vspace{6pt}

\fancyhead[C] {\small{\textsc{\Title}}}
\fancyhead[OL]{\small{\thepage}}
\fancyhead[ER]{\small{\thepage}}

\myfootnote
{
  Version {\Version} from \today;
  Email:  \textit{khoeflinger@\,t-online.de}
}

%------------------------------------------------------------------------------%
%                                 INTRODUCTION                                 %
%------------------------------------------------------------------------------%
The mathematical area of \textit{functional analysis} has been developed
in the first half of the 20$^{\textrm{th}}$ century
by a multiplicity of pioneers,
where we single out
{D.\,Hilbert},
{S.\,Banach},
{F.\,Riesz},
{J.\,v.\,Neumann},
{M.\,Stone} and
{K.\,O.\,Friedrichs}.
Nowadays it forms a well-elaborated subject in modern mathematics.

\subparagraph{}
Rather functional analysis constitutes a closed subject in its own,
it also provides a couple of powerful tools
that can be applied to the boundary and initial value problems
in mathematical physics and engineering science.
These problems are connected to partial differential equations,
usually governed by elliptic operators.
Moreover,
possible solutions of the problems have to fulfill prescribed values
on the boundary of the underlying domain
and in the evolutionary case also prescribed values at some initial time $t_0$.

\paragraph{}
This exposition consists of overall ten sections,
where the first eight ones deal with abstract functional analytic topics,
all of them sharing the notion of \textit{linear operator}
as one of their central concepts.
The nineth section outlines applications of the presented theory
to linear partial differential operators and the boundary value problems
governed by them,
while the last one provide a glimpse to applications of distribution theory
to linear partial differential equations.

\subparagraph{} %- honoration to Prof. Dr. Werner
I wrote this treatise in honour of Prof.\;Dr.\;P.\;Werner (1932-2018),
who conveyed my fascination of these beautiful mathematical topics
by his lessons hold in 1984 and 1985
at Mathematical Institute A of University Stuttgart.

%------------------------------------------------------------------------------%
\section{Fundamentals of Operator Theory}                                      %
%------------------------------------------------------------------------------%
Throughout this text let $\F$ be one of the fields $\R$ and $\C$.
By a \textit{normed space} $\banX$ we mean a $\F$-linear space
that differs from other abstract linear spaces
by the presence of a particular mapping $\norm{\,\cdot\,}\: \banX \xto \R$,
denoted the \textit{norm} on $\banX$,
which fulfills for all $x,y \in \banX$ and $\lambda \in \F$

\begin{itemize}[leftmargin=22mm,itemsep=2pt]
\item[(N-1)]
$\norm{x} \, \geq \, 0$,

\item[(N-2)]
$\norm{\lambda \, x} \= \abs{\lambda} \cdot \norm{x}$,

\item[(N-3)]
$\norm{x+y} \, \leq \, \norm{x} + \norm{y}$
\; \textit{(triangle inequality)},

\item[(N-4)]
$\norm{x} \= 0$ \iff $x \eq \0$.
\end{itemize}

\subparagraph{}
Linear spaces can of course be furnished with different norms.
Arbitrary norms $\norm{\,\cdot\,}_{\alpha}$ and $\norm{\,\cdot\,}_{\beta}$
defined on a linear space $\banX$
are called \textit{equivalent} to each another,
if there are constants $c_1,c_2 > 0$ such that
$
c_1 \norm{x}_{\alpha} \leq
    \norm{x}_{\beta}  \leq 
c_2 \norm{x}_{\alpha}$
for all $x \in \banX$.

\subparagraph{}
Each normed linear space $\banX$ may be equipped
with a very natural topology $\topT$,
called the \textit{norm topology} on $\banX$.
The definition of $\topT$ is done either by the assertion
that $\topT := \topT_{\norm{\cdot}}$ is the smallest topology on $\banX$
such that the norm mapping remains continuous,
or by the metric topology $\topT_d$,
when $\banX$ is considered as a metric space with metric function
$d(x,y) := \norm{x-y} \fa x,y \in \banX$.
In fact,
both definitions lead to the same collection of so-called
"open" subsets in $\banX$,
so we have $\topT := \topT_{\norm{\cdot}} \eq \topT_d$,
which justifies to speak of \textit{the} topology of a normed linear space.
Hence a multiplicity of topological concepts,
like the \textit{closure} of a set,
the \textit{neighbourhood} of a point,
the \textit{convergence} of a series and many others are meaningful
in the setting of normed linear spaces.
Notice
that equivalent norms on $\banX$ are creating one and the same topology.
For example,
if $\banX$ is a finite-dimensional linear space,
then all norms are equivalent
and consequently there is exactly one norm topology on $\banX$.

\subparagraph{} %- sequentially complete normed spaces, Banach spaces
Recall that a sequence $(x_n)$ in $\banX$ is called \textit{Cauchy},
if to each $\epsilon > 0$ there is $N_{\epsilon} \in \N$
such that $m,n > N_{\epsilon}$ implies $\norm{x_m-x_n} < \epsilon$.
Any normed linear space $\banX$ is called \textit{sequentially complete}
(or shortly \textit{complete}),
if for any Cauchy sequence $(x_n)$ in $\banX$
there is $x \in \banX$
such that $\lim_{\,n \xto \infty} \norm{x_n-x} \eq 0$.
Equivalently,
$\banX$ is complete
iff
for every absolutely convergent series
$\sum_{n=1}^{\infty} \norm{x_n} < \infty$
there is $x \in \banX$ such that $\norm{x-x_n} \xto 0$ for $n \xto \infty$.
A complete normed space is called a \textit{Banach space}.

\begin{comment}
\paragraph{} %- honoration to Prof. Dr. Werner
I wrote this treatise in honour of Prof.\;Dr.\;P.\;Werner (1932-2018),
who conveyed my fascination of these beautiful mathematical topics
by his lessons hold in 1984 and 1985
at Mathematical Institute A of University Stuttgart.
\end{comment}

\paragraph{} %-- closed linear operators
%------------------------------------------------------------------------------%
The notation of \textit{densely defined} and \textit{closed linear operator}
will be the key concept in the existence and uniqueness theory
of the mentioned problems and hence plays a dominant role in our considerations.
In fact,
there are various examples of mappings in functional analysis,
which prove to be linear, densely defined and closed.
Especially a lot of partial differential expressions
own this remarkable property,
supposed they are acting between suitable Banach spaces of functions.

\subparagraph{} %-- densely defined and closed operators
%------------------------------------------------------------------------------%
Let us turn to the definition of our key notation.
First remember that by a \textit{topological linear space}
is meant an arbitrary set $\banX$,
which first 
is a $\F$-linear space
($\F$ being the field of real or complex numbers)
and a topological space fulfilling \textsc{Hausdorff}'s separation axiom
at the same time,
and second the basic operations
$(x,y) \mapsto x+y$ and $(\lambda,x) \mapsto \lambda x$ in the space $\banX$
are continuous.
Then a mapping $\linA \: \domain{\linA} \xsubset \banX \xto \banY$
is denoted a \textit{densely defined closed linear operator},
if the following conditions are fulfilled:

\begin{itemize}[leftmargin=12mm]
\item[(C-1)]
$\banX$ and $\banY$ are topological linear spaces;
\item[(C-2)]
$\domain{\linA}$,
the \textit{domain of definition} of the mapping $\linA$,
is a linear subspace of $\banX$,
which lies dense in $\banX$;
\item[(C-3)]
$\linA$ is \textit{linear},
that is,
$\linA(x+y) \eq \linA x + \linA y$ and $\linA(\lambda x) \eq \lambda \linA x$
for $x,y \xeof \domain{\linA}$ and $\lambda \xeof \F$;
\item[(C-4)]
$\linA$ is \textit{graph-closed} or briefly \textit{closed},
that is,
the set
\[
\Gamma_{\linA} \eqdef
\{\,(x,\linA x): x \xeof \domain{\linA}\,\} \; \xsubset \; \banX \xtimes \banY,
\]
denoted the \textit{graph} of $\linA$,
is closed in $\banX \xtimes \banY$
{\wrt} the product topology $\topT_{\,\banX \xtimes \banY}$,
which especially implies that
the kernel $\kernel{\linA} \eqdef \{ x \xeof \domain{\linA}: \linA x \eq \0 \}$
of $\linA$ forms a closed set in $\banX$.
\end{itemize}

\subparagraph{} %- domain and range space
The linear spaces $\banX$ and $\banY$ are said
to be the \textit{domain space} and the \textit{range space}
of $\linA$.
In case of $\banY \eq \banX$,
however,
it is usual to say that $\linA$ is an operator acting in $\banX$.

\paragraph{} %- sequentially closed operators
More can be said if $\banX$ and $\banY$ constitute {\FRECHET}, Banach
or Hilbert spaces:
in these cases the operator $\linA$ is closed \iff
it is \textit{sequentially closed},
that is,
$x_n \xto x$ and $\linA x_n \xto y$ implies
$x \xeof \domain{\linA}$ and $\linA x \eq y$,
briefly written
\[
\begin{rcases}
 x_n       \xto x \quad \\
 \linA x_n \xto y \quad \\
\end{rcases}
\quad \Longrightarrow \quad
x \xeof \domain{\linA} \, \and \, \linA x \eq y.
\]

\subparagraph{} %- closeness of operators based on Schauder estimates
A further statement to ensure closenss of a linear operator
relies on the verification of sequentially closeness
and uses the possible validness of an abstract \textsc{Schauder} estimate:

\begin{theorem1}\label{BROWDER_THEOREM}
Let $\linA_0 \: \banX \xto \banY$ be a continuous linear operator,
where the Banach space $\banX$ is continuously and densely embedded
into the Banach space $\banY$ by a mapping $\iota \: \banX \xto \banY$.
If the estimate
\[
\norm{x}_{\banX} \, \leq \,
c \cdot
\{ \,\norm{\iota x}_{\banY} + \norm{\linA_0 \, x}_{\banY} \, \}
\]
is fulfilled for some $c {\,>\,} 0$ and all $x \xeof \banX$,
then the linear operator $\linA$ acting in $\banY$ with
$\domain{\linA} := \iota(\banX)$ and
$\linA\, \iota x := \linA_0 x$ for $x \xeof \banX$ is closed.
\end{theorem1}

\paragraph{} %-- pre-closed operators
%------------------------------------------------------------------------------%
Many linear operators acting between Banach or Hilbert spaces are not closed,
but can be extended to a closed operator.
We shall call such operators \textit{closable} or \textit{pre-closed}.
If a linear operator $\linA$ is closable,
then there is a \textit{minimal closed extension} $\linA_c$ of $\linA$,
called the \textit{closure} of $\linA$.
This mapping $\linA_c$ is given by the fact
that each other closed extension of $\linA$ extends $\linA_c$,
and $\linA_c$ is nothing but the intersection
of all closed extensions of $\linA$.
There is a well-known criterion for closability:
the  operator $\linA$ is closable,
if each sequence $\{(x_n,\linA x_n)\}_{n \xeof \N}$ converging to $(\0,y)$
implies $y \eq \0$,
briefly written
\[
\begin{rcases}
 x_n       \xto \0 \quad \\
 \linA x_n \xto  y \quad \\
\end{rcases}
\quad \Longrightarrow \quad
y \= \0.
\]

\paragraph{} %-- continuous operators
%------------------------------------------------------------------------------%
On the other hand,
a linear operator $\linA \: \domain{\linA} \xsubset \banX \xto \banY$
is called \textit{continuous},
if the preimage of each open set in $\banY$ forms an open set in $\banX$.
This is the well-known definition of continuity for mappings
between arbitrary topological spaces.

\subparagraph{} %- continuity in normed linear spaces
For linear mappings between normed linear spaces,
however,
we may replace continuity by \textit{sequential continuity},
which means that a linear operator $\linA$ is continuous in $x \xeof \banX$
if $x_n \xto x$ for $n \xto \infty$ implies $\linA x_n \xto \linA x$
for each convergent series with limit element $x$.
It is also known that $\linA$ is continuous iff
it is so in at least one element $x \xeof \banX$
(for example in $x \eq \0$).

\subparagraph{} %- closed graph theorem
Let us emphasize that in general 
the concepts of closeness and continuity of a linear operator $\linA$
are not comparable.
However,
if $\banX$ and $\banY$ are {\FRECHET}, Banach or Hilbert spaces
and $\domain{\linA}$ is a \textit{closed} set in $\banX$,
then closeness of $\linA$ is an easy consequence of continuity of $\linA$.
Reversely,
the \textit{closed graph theorem} claims that closeness implies continuity,
so both concepts are equivalent to each other
under this restrictive assumption.
This is especially true in the following situation:
if $\linA$ is densely defined and continuous,
then it can be extended to a continuous linear operator
$\linA_{ext} \: \banX \xto \banY$,
and since $\banX$ forms a closed set,
the extension\footnote{
\,The operator $\linB$ is called an \textit{extension} of the operator $\linA$
(notation $\linA \prec \linB$),
if $\domain{\linA} \xsubset \domain{\linB}$ and the restriction of $\linB$
to $\domain{\linA}$ coincides with $\linA$.}
$\linA_{ext}$ is closed,
too.
Hence continuous linear operators of the form $\linA \: \banX \xto \banY$
fit in a natural way
into our class of densely defined and closed linear operators.

\subparagraph{} %- bounded inverse theorem
Another application of the closed graph theorem is of outstanding meaning:
if a linear operator $\linA$ is closed and bijectiv,
then its inverse $\linA^{-1}$ is also linear and closed,
and since the domain of definition $\domain{\linA^{-1}}$
is the entire space $\banY$,
the closed graph theorem implies that $\linA^{-1} \: \banY \xto \banX$
is even continuous and bounded,
respectively.
This fact is known as the \textit{bounded inverse theorem} for closed operators.
But notice that it is also valid for continuous linear operators
$\linA \: \banX \xto \banY$,
which is a consequence of \textsc{Banach}'s \textit{open mapping theorem}.

\paragraph{} %-- spectral theory
%------------------------------------------------------------------------------%
\textit{Spectral theory} is a subject within functional analysis
that proves to be helpful in the study of abstract linear equations
of type $\linA x \eq y$
as well as of eigen\-value equations of type $\linA x {-} \lambda x \eq y$.

\subparagraph{} %- resolvent set and spectrum
Given a complex Banach space $(\banX,\norm{\,\cdot\,})$
and a densely defined linear operator
$\linA \: \domain{\linA} \xsubset \banX \xto \banX$,
the basic idea in spectral theory is to split the complex plane $\C$
into two disjoint sets $\rho(\linA)$ and $\sigma(\linA)$,
denoted the \textit{resolvent set} and the \textit{spectrum} of $\linA$,
respectively,
where by definition $\zeta \xeof \rho(\linA)$ \iff

\begin{itemize}[leftmargin=12mm]
\item[(R-1)]
$\linA-\zeta := \linA - \zeta \, Id$\footnote{
\,$Id \: \banX \xto \banX$ is the identity mapping
in the linear space $\banX$.}
$\: \domain{\linA} \xto \banX$ is bjective,

\item[(R-2)]
the inverse mapping $(\linA - \zeta)^{-1} \: \banX \xto \domain{\linA}$
is continuous,
\end{itemize}

otherwise $\zeta$ is said to belong to $\sigma(\linA)$.
The numbers in the sets $\rho(\linA)$ and $\sigma(\linA)$
are called \textit{regular values} and \textit{spectral values},
respectively.
Notice that if $\linA$ is closed,
then property (R-2) is a consequence of the closed graph theorem
and does not have to be postulated.

\subparagraph{} %- properties of the resolvent set and the spectrum
The set $\rho(\linA)$ is open,
but may also be the empty set, a proper subset of the complex plane
or the set $\C$ itself.
One often denotes a linear operator $\linA$ to be \textit{regular},
if $\rho(\linA) \neq \emptyset$,
and in this case $\linA$ must necessarily be closed.
If $\linA$ is continuous and everywhere defined,
then we actually have $\emptyset \neq \rho(\linA) \neq \C$,
so $\sigma(\linA) \neq \emptyset$ and
$\abs{\lambda} \leq \norm{\linA}$ for all $\lambda \xeof \sigma(\linA)$.

\subparagraph{} %-- resolvent functions
%------------------------------------------------------------------------------%
Moreover,
the operator-valued mapping
$R(\linA,\cdot) \: \rho(\linA) \xto \BO{\banX}$
defined by $\zeta \mapsto R(\linA,\zeta) := (\linA-\zeta)^{-1}$ is called
the \textit{resolvent function} or briefly the \textit{resolvent} of $\linA$,
where $\BO{\banX}$ denotes the unital Banach algebra
of bounded linear operators acting in $\banX$.
The resolvent is holomorphic in every $\zeta_0 \xeof \rho(\linA)$,
that is,
$R(\linA,\zeta)$ can be represented by a power series
$
\sum_{n \xeof \N} \, b_n(\zeta_0) \, (\zeta-\zeta_0)^n
$
for $\abs{\zeta-\zeta_0} $ sufficiently small
and with coefficents $b_n(\zeta_0) \xeof \BO{\banX}$.

\subparagraph{} %- resolvent equation
All the mappings $R(\linA,\zeta)$ are coupled by the \textit{resolvent equation}
\begin{equation}\label{RESOLVENT_EQUATION}
R(\linA,\zeta_1) - R(\linA,\zeta_2)
\=
(\zeta_2-\zeta_1) \, R(\linA,\zeta_1) \circ R(\linA,\zeta_2)
\end{equation}
for $\zeta_1,\zeta_2 \xeof \rho(\linA)$,
which especially implies that
the linear operators $R(\linA,\zeta_1)$ and $R(\linA,\zeta_2)$
commute with each another.

\paragraph{} %- norm estimate of the resolvent
Estimates for the operator norm of resolvents play an important role
in later sections.
In fact,
the estimate
\begin{equation}\label{RESOLVENT_ESTIMATE_I}
\norm{R(\linA,\zeta)} \, \geq \,
r(R(\linA,\zeta))     \, \geq \,
\mbox{dist}(\zeta,\sigma(\linA))^{-1}
\end{equation}
holds for any $\zeta \xeof \rho(\linA)$,
where $r(R(\linA,\zeta)) \eq
\sup \{ \abs{\lambda} \: \lambda \in \sigma(R(\linA,\zeta)) \}$
is the \textit{spectral radius} of $R(\linA,\zeta)$ and
$\mbox{dist}(\zeta,\sigma(\linA))$ denotes the minimal distance
between the singleton $\{\zeta\}$ and the spectrum $\sigma(\linA)$.
This is a consequence of the so-called \textit{spectral mapping theorem}
\[
\sigma(R(\linA,\zeta)) \setminus \{0\} \=
\{(\lambda-\zeta)^{-1}: \lambda \xeof \sigma(\linA)\},
\]
which establishes a connection between the spectra of the resolvents
$R(\linA,\zeta)$ and the operator $\linA$ itself.

\paragraph{} %-- various kinds of spectra
%------------------------------------------------------------------------------%
We next will have a look at the spectral values
of a closed operator $\linA$ and study the possible reasons
that prevent the operator $\linA - \lambda$ to be bijective
for $\lambda \xeof \sigma(\linA)$.
Preliminary,
let the \textit{range} $\range{\linA}$ of $\linA$ be the linear subspace
$\{y \in \banX: y \eq \linA \, x \mbox{ for some } x \in \domain{\linA} \}$.

\begin{itemize}[leftmargin=4mm,itemsep=5pt]
\item[1.]
Recall from linear algebra
that a complex number $\lambda$ is called an \textit{eigen\-value}
of a linear operator $\linA$,
if there is an element $x \xeof \banX$ such that
$x \neq 0$ and $\linA \,x \eq \lambda x$.
In this case the element $x$ is called an \textit{eigenvector} of $\linA$
{\wrt} $\lambda$,
and all these eigenvectors form a linear subspace of $\banX$,
called the \textit{eigenspace} of $\linA$ for $\lambda$.
The \textit{point spectrum} of $\linA$ is defined to be
the set $\sigma_p(\linA)$ of all eigen\-values of $\linA$.
Clearly,
for $\lambda \xeof \sigma_p(\linA)$ the operator
$\linA - \lambda$ is not injective,
since its kernel is non-trivial.
\end{itemize}

If $\linA - \lambda$ is injective,
but not surjective,
then $(\linA-\lambda)^{-1}$ is defined on the range $\range{\linA - \lambda}$.
We distinguish between two cases:

\begin{itemize}[leftmargin=4mm]
\item[2.]
if $\range{\linA - \lambda}$ is dense in $\banX$,
then $\lambda$ is said to be an element of the set $\sigma_c(\linA)$,
called the \textit{continuous spectrum} of $\linA$,
otherwise

\item[3.]
the number $\lambda$ belongs to the \textit{residual spectrum}
$\sigma_r(\linA)$.
\end{itemize}

\subparagraph{}
In summary,
for each closed linear operator
$\linA \: \domain{\linA} \xsubset \banX \xto \banX$
the set of complex numbers forms a disjoint union
of the resolvent set and the three different kinds of spectra, namely
\[
\C \= \rho    (\linA) \, \sqcup
        \, \sigma_p(\linA) \, \sqcup
        \, \sigma_c(\linA) \, \sqcup
        \, \sigma_r(\linA).
\]
The decomposition of the spectrum into these three parts
is the classical approach.
We shall introduce in the next section
another kind of decomposition of the set $\sigma(\linA)$,
which relies on the concept of \textit{semi-Fredholm operator}.

\paragraph{} %-- adjoint operators
%------------------------------------------------------------------------------%
In operator theory it is often advantageous to combine
the study of a densely defined linear operator
with that of its \textit{adjoint}.

\subparagraph{} %- definition of the adjoint operator
Let $\banX,\banY$ be arbitrary Banach spaces,
$\banX^*,\banY^*$ be their adjoint spaces\footnote{
\,By the \textit{adjoint space} $\banX^*$ of a normed linear space $\banX$
is meant the linear space of continuous linear functionals $f \: \banX \xto \F$,
which itself forms a Banach space with norm
$\norm{f} := \sup_{\norm{x}=1} \abs{f(x)} $.
The celebrated \textsc{Hahn-Banach} theorem claims that $\banX^* \neq \{\0\}$,
so there are non-trivial continuous linear functions on $\banX$.
}
and $\linA \: \domain{\linA} \xsubset \banX \xto \banY$
be a linear operator with domain of definition $\domain{\linA}$
lying dense in $\banX$.
Then the domain of definition $\domain{\linA^*}$
of the {adjoint operator} $\linA^*$ is defined by
\[
\domain{\linA^*} \eqdef
\{\, g \xeof \banY^*:
g \circ \linA \: \domain{\linA} \xto \F \mbox{ is continuous}\,\}.
\]
For $g \xeof \domain{\linA^*} $ let
$\linA^* g := \overline{g \circ \linA}$,
the unique extension of the continuous linear functional $g \circ \linA$
to the whole of $\banX$,
which implies $\linA^* g \xeof \banX^*$.
Thus the linear operator $\linA^*$
has domain space $\banY^*$ and range space $\banX^*$,
and the domain of definition $\domain{\linA^*}$
forms a linear subspace of $\banY^*$,
on which $\linA^*$ proves to be linear.
We emphasize that
$\linA$ is supposed to be densely defined in order to introduce its adjoint.

\subparagraph{} %- resolvent set and spectrum of adjoint operators
The resolvent set $\rho(\linA^*)$ and spectrum $\sigma(\linA^*)$
are mirrored sets of $\rho(\linA)$ and $\sigma(\linA)$
and obtained by reflection at the real line.
There is also an interesting identity between the spectrum
of a densely defined and closed linear operator $\linA$
and that of its adjoint $\linA^*$,
namely
$\sigma_r(\linA) \eq \sigma_p(\linA^*)$ (see \cite{EN}, p.\,161).

\paragraph{} %-- orthogonality relationships
The usage of the operator $\linA^*$ is founded by the validness
of two relationships between the kernels and ranges
of $\linA$ and $\linA^*$,
namely
\begin{equation}\label{ORTHOGONALITY_RELATIONS}
\overline{\range{\linA^*} } \= \kernel{\linA}^{\bot}
\quad \and \quad
\overline{\range{\linA} } \= ^{\bot}\kernel{\linA^*},
\end{equation}
where for any subset $M \xsubset \banX$ and $N \xsubset \banX^*$
the \textit{annihilator}     $M^{\bot}$ and
the \textit{pre-annihilator} $^{\bot}N$
are defined by
\[
\begin{cases}[1.2]
\quad M^{\bot} \eqdef \{\,x^*\xeof \banX^*: \scalar{x}{x^*} = 0
      \fa x   \xeof M \,\},\\
\quad ^{\bot}N \eqdef \{\,x  \xeof \banX:   \scalar{x}{x^*} = 0
      \fa x^* \xeof N \,\}\\
\end{cases}
\]
and $\scalar{x}{x^*} := x^*(x)$ denotes the pairing of elements in $\banX$
and $\banX^*$.

\subparagraph{}
In the Hilbert space setting $\banX \eq \banY \eq \hilH$,
the sets $M^{\bot}$ and $^{\bot}N$ are nothing
but the orthogonal spaces of $M$ and $N$,
where $N^{\bot} \xsubset \hilH^*$ can be seen as a linear subspace of $\hilH$,
which is a consequence of \textsc{Riesz}'s representation theorem
establishing a one-to-one relationship
between the spaces $\hilH$ and $\hilH^*$.

\begin{comment}
\paragraph{} %- formally adjoint pairs
We finally introduce the concept of \textit{adjoint pair} of linear operators.
Two densely defined linear operators
$\linA\:\domain{\linA} \xsubset \banX   \xto \banY$ and
$\linB\:\domain{\linB} \xsubset \banY^* \xto \banX^*$
are said to be \textit{formally adjoint} to each other
or to be a \textit{formally adjoint pair} $(\linA,\linB)$ if
\[
\scalar{\linA \,x}{f}_{\banY \xtimes \banY^*} \=
\scalar{x}{\linB f}_{\banX \xtimes \banX^*}
\]
holds for all $x \in \domain{\linA}$ and $f \in \domain{\linB}$.
In addition,
we define the \textit{restricted adjoint} $\linB_r$
of the linear operator $\linB$ according to
$\domain{\linB_r} :=
\{ x \in \banX: \mbox{there is } y \in \banY \mbox{ such that }
(\linB f)\,x \eq f(y) \fa f \in \domain{\linB} \}$
and $\linB_r \, x \eqdef y$.
The mapping $\linB_r$ acts between the Banach spaces $\banX$ and $\banY$
and proves to be linear.
We may reformulate the notation of {formally adjoint pair} by saying that
$\linA$ and $\linB$ are formally adjoint to one another
if $\linB_r$ is an extension of $\linA$.

\subparagraph{} %- realization of a formally adjoint pair
To the end,
by a \textit{realization} of a formally adjoint pair $(\linA,\linB)$
we mean any densely defined and closed linear operator $R$
lying between $\linA$ and $\linB_r$,
so $\linA \prec R \prec \linB_r$ holds.
\end{comment}

%------------------------------------------------------------------------------%
\section{Normally Solvable Operators}                                          %
%------------------------------------------------------------------------------%
Given a densely defined linear operator $\linA$,
the orthogonality relations \eqref{ORTHOGONALITY_RELATIONS} can be sharpened
if the ranges of $\linA$ and $\linA^*$ are \textit{closed} subspaces
of $\banY$ and $\banX^*$,
so in this case we even would have
\[
{\range{\linA} } \= ^{\bot}\kernel{\linA^*}
\quad \and \quad
{\range{\linA^*} } \= \kernel{\linA}^{\bot}.
\]
One also says that
a linear operator $\linA$ is \textit{normally solvable}
if it is densely defined and
$\range{\linA} \eq \, ^{\bot}\kernel{\linA^*}$ holds.

\subparagraph{} %-- closed range theorem
%------------------------------------------------------------------------------%
The celebrated \textit{closed range theorem} 
is one of the classical results in functional analysis 
and provides equivalent conditions
for a densely defined and closed linear operator to have closed range:

\begin{theorem1}\label{CLOSED_RANGE_THEOREM}
Let $\banX,\,\banY$ be Banach spaces
and $\linA \: \domain{\linA} \xsubset \banX \xto \banY$
be densely defined and closed.
Then the following assertions are equivalent:

\begin{itemize}[leftmargin=9mm]
\item[(a)] $\linA$ has closed range.
\item[(b)] $\linA^*$ has closed range.
\item[(c)] $\range{\linA} $ is the pre-annihilator of $\kernel{\linA^*}$,
           {\ie} $\range{\linA} \= ^{\bot}\kernel{\linA^*} $.
\item[(d)] $\range{\linA^*}$ is the annihilator of $\kernel{\linA}$,
           {\ie} $\range{\linA^*} \= \kernel{\linA} ^{\bot}$.
\end{itemize}
\end{theorem1}

\subparagraph{}
For a proof of this theorem see for example \cite{Yo}, pp. 205-207.
The statement was originally proved by \textsc{S.\,Banach}
for continuous linear operators and
have later been extended to closed linear operators,
see \textsc{Browder} \cite{Br1}, theorem\,1.1 and theorem\,1.2.

\paragraph{} %-- operators bounded below
%------------------------------------------------------------------------------%
Besides this theorem,
it is natural to search for further conditions
that imply normal solvability.
A linear operator $\linA$ is called \textit{bounded below}
if there is $\theta {\,>\,} 0$ such that
\begin{equation}\label{BOUNDED_BELOW_CONDITION}
\norm{\linA x}_Y \, \geq \, \theta \norm{x}_X \fa x \xeof \domain{\linA}.
\end{equation}

\subparagraph{} %- bounded below theorem
The \textit{bounded-below theorem} provides a fundamental assertion
in operator theory and is helpful to verify the closeness of $\range{\linA}$:

\begin{theorem1}\label{BOUNDED_BELOW_THEOREM}
If $\linA$ fulfills \eqref{BOUNDED_BELOW_CONDITION},
then $\linA$ is injective and the operator
$\linA^{-1} \: \range{\linA} \xsubset \banY \xto \banX$ is bounded
with operator norm $\norm{\linA^{-1}} \leq \theta^{-1}$.
Furthermore,
the range $\range{\linA}$ is closed \iff $\linA$ is closed.
\end{theorem1}

Indeed,
injectivity of $\linA$ can easily proved,
supposed \eqref{BOUNDED_BELOW_CONDITION} holds,
and closeness of $\range{\linA}$ is a consequence of sequentially closeness
of $\linA$.

\paragraph{} %-- semi-Fredholm operators
%------------------------------------------------------------------------------%
We want to extend our considerations to a more general class of operators.
Preliminary,
let $\alpha(\linA)$ (the \textit{nullity index} of $\linA$)
and $\beta(\linA)$ (the \textit{deficiency index} of $\linA$)
be the Hamel dimensions of the linear subspaces
$\kernel{\linA} \xsubset \banX$ and $\cokernel{\linA} \xsubset \banY$.

\subparagraph{} %- Fredholm index
Both integers together can be used
to measure the deviation of $\linA$ from bijectivity.
Moreover,
if at least one of them is finite,
then the so-called \textit{Fredholm index}
\[
\chi(\linA) \eqdef
\begin{cases}[1.2]
\quad + \infty, & \alpha(\linA) = \infty\\
\quad \alpha(\linA) \, - \, \beta(\linA), &
\alpha(\linA) < \infty, \, \beta(\linA) < \infty\\
\quad - \infty, & \beta(\linA) = \infty\\
\end{cases}
\]
is well-defined and takes values in the set
$\{- \infty\} \cup \N \cup \{+ \infty\}$.
In the special case,
where $\alpha(\linA) {\,<\,} \infty$,
$\beta(\linA) {\,<\,} \infty$
and consequently $\chi(\linA) \xeof \Z$,
the operator $\linA$ is said to have \textit{finite index}.

\paragraph{} %- definition of semi-Fredholm operator
The linear operator $\linA \: \domain{\linA} \xsubset \banX \xto \banY$
is called \textit{semi-Fredholm},
if
\begin{itemize}[leftmargin=12mm]
\item[(F-1)] the range $\range{\linA}$ is closed,
\item[(F-2)] the kernel $\kernel{\linA} $ or the cokernel $\cokernel{\linA} $
             is finite-dimensional
             (rsp. $\alpha(\linA) {\,<\,} \infty$ or
                   $\beta(\linA)  {\,<\,} \infty$).
\end{itemize}

Notice the remarkable fact
that if $\linA$ is continuous with $\domain{\linA} \eq \banX$
and $\beta(\linA) {\,<\,} \infty$,
then the range $\range{\linA}$ is necessarily closed,
since it constitutes the topological complement
of the finite-dimensional linear subspace $\cokernel{\linA} \xsubset \banY$,
that is,
$\banX \eq \range{\linA} \oplus \cokernel{\linA}$.

\subparagraph{} %- lower and upper semi-Fredholm operators
We refine the concept of semi-Fredholm operator as follows:
a linear operator $\linA \: \domain{\linA} \xsubset \banX \xto \banY$
with closed range $\range{\linA}$ is called

\begin{itemize}[leftmargin=12mm]
\item[(F-3)]
\textit{upper semi-Fredholm},
if $\alpha(\linA) {\,<\,} \infty$ rsp. $\chi(\linA) {\,<\,} \infty$,

\item[(F-4)]
\textit{lower semi-Fredholm},
if $\beta(\linA) {\,<\,} \infty$ rsp. $\chi(\linA) {\,>\,} {-}\infty$.
\end{itemize}

\subparagraph{} %- example of upper semi-Fredholm operator
Examples of upper semi-Fredholm operators are the bounded-below ones.
In fact,
theorem \ref{BOUNDED_BELOW_THEOREM} implies that
these operators are upper semi-Fredholm with vanishing nullity index.

\paragraph{} %- Schauder estimates and Peetre's lemma
If a continuous linear operator $\linA$ satisfies
an abstract \textit{Schauder estimate},
one may quote a sufficient condition for $\linA$ to be upper semi-Fredholm
that is known as \textsc{Peetre}\textit{'s lemma}:

\begin{theorem1}\label{PEETRE_THEOREM_II}
Let $\banX, \banY, \banZ$ be three Banach spaces,
where $\banX$ is compactly embedded in $\banZ$,
that is,
there is an injective and compact linear mapping \footnote{
\,
A linear operator $\compA \: \banX \xto \banY$ is called \textit{compact},
if every bounded sequence $\{x_n\}_{n \xeof \N}$ in $\banX$
owns a subsequence $\{x_{n_k}\}_{k \xeof \N}$
with $\{ \compA x_{n_k} \}_{k \xeof \N}$ convergent in $\banY$.}
$\iota \: \banX \xto \banZ$.
Then a continuous operator $\linA \: \banX \xto \banY$ is upper semi-Fredholm
\iff
there is $c {\,>\,} 0$ such that
\[
\norm{x}_{\banX} \, \leq \,
c \, (\norm{\iota x}_{\banZ} \+ \norm{\linA \,x}_{\banY})
\quad \quad (x \in \banX).
\]
\end{theorem1}

\subparagraph{}
For a proof of this theorem, see \cite{Wl}, theorem 12.12, page 180.

\paragraph{} %-- perturbation theory of semi-Fredholm operators
%------------------------------------------------------------------------------%
Nowadays there is a well-elaborated \textit{perturbation theory}
for semi-Fredholm operators available,
where conditions are studied,
under which a linear operator of the form $\linA + \linB$ is semi-Fredholm,
assumed $\linA$ is semi-Fredholm and $\linB$ fulfills certain properties.
In this case operator $\linB$ is regarded as a perturbation of $\linA$.

\paragraph{} %- Yood's theorem
We next quote two theorems {\wrt} conditions on $\linB$,
which ensure the semi-Fredholm property of the operator $\linA + \linB$.
We basically assume that $\banX$ and $\banY$ are Banach spaces.

\begin{theorem1}\label{YOOD_THEOREM}
Let $\linA \: \banX \xto \banY$ be semi-Fredholm
and $\linB \: \banX \xto \banY$ be linear and bounded
with sufficiently small norm $\norm{\linB}$.
Then the operator $\linA + \linB$ remains semi-Fredholm
with $\chi(\linA+\linB) \eq \chi(\linA)$.
\end{theorem1}

\subparagraph{} %- punctured neighbourhood theorem
In other words,
this theorem asserts
that the set of semi-Fredholm operators is open in the set $\BO{\banX,\banY}$.
Especially linear operators of the form $\linA + \lambda$ for $\abs{\lambda} $ 
small enough are semi-Fredholm again.
But more can be said on the values $\alpha(\linA+\lambda)$,
$\beta(\linA+\lambda)$ and $\chi(\linA+\lambda)$
by \textsc{Gohberg}\textit{'s punctured neighbourhood theorem}:

\begin{theorem1}\label{PUNCTURED_NEIGHBOURHOOD_THEOREM}
Let the continuous linear operator $\linA \: \banX \xto \banX$ be semi-Fredholm.
Then there is $\lambda_0 {\,>\,} 0$
such that for all $\lambda \xeof (0,\lambda_0)$ we have
$\alpha(\linA + \lambda) \leq \alpha(\linA)$,
$\beta (\linA + \lambda) \leq \beta (\linA)$
and
$\chi  (\linA + \lambda) \eq  \chi  (\linA)$.
\end{theorem1}

\paragraph{} %-- Fredholm operators
%------------------------------------------------------------------------------%
Let us refine the definition of semi-Fredholm operator to a stronger concept:
a continuous linear operator $\linA \: \banX \xto \banY$
acting between normed linear spaces $\banX$ and $\banY$ is called
\begin{itemize}[leftmargin=12mm]
\item[(F-5)]
\textit{Fredholm},
if $\chi(\linA)$ is finite,
{\ie} $\alpha(\linA) {\,<\,} \infty$ and $\beta(\linA) {\,<\,} \infty$.
\end{itemize}

\subparagraph{} %- Fredholm property of the adjoint
The relationships \eqref{ORTHOGONALITY_RELATIONS} imply
that iff $\linA$ is densely defined and Fredholm,
then the adjoint operator $\linA^*$ is also Fredholm and
$\chi(\linA^*) \eq {-} \chi(\linA)$ holds.

\paragraph{} %- Fredholm's alternative theorem
In the following we cite the \textsc{Fredholm} \textit{alternative theorem}
as the outstanding result in Fredholm theory
{\wrt} the solvability of the linear equation $\linA x \eq y$.
Remember the \textit{rank-nullity theorem} known from linear algebra,
claiming that if a linear space $\vecV$ is finite-dimensional,
then a linear mapping $\linA \: \vecV \xto \vecV$ is surjective
\iff it is injective.

\subparagraph{}
The next theorem is a considerable generalization of the rank-nullity theorem
to arbitrary Banach spaces and Fredholm operators with vanishing index:

\begin{theorem1}\label{FREDHOLM_ALTERNATIVE_THEOREM}
Let $\banX,\,\banY$ be Banach spaces and
the continuous linear operator $\linA \: \banX \xto \banY$ be Fredholm
with index $\chi(\linA) \eq 0$.

\subparagraph{}
According to the solvability of the linear equation $\linA \, x \eq y$,
exactly one of the following cases occurs:

\begin{itemize}[leftmargin=9mm]
\item[(a)]
$\linA \, x \eq y$ is uniquely solvable for all $y \xeof \banX$,
so $\linA$ is bijective,

\item[(b)]
$\linA \, x \eq y$ is not for all $y \xeof \banX$ solvable,
and $\linA$ is neither injective nor surjective in this case.
However,
if the equation is solvable for an element $y_0 \xeof \banX$,
then the set of solutions forms a finite-dimensional subspace of $\banX$.
Moreover,
the equation $\linA \, x \eq y$ owns a solution $x \xeof \banX$
\iff
$y \xeof \,^{\bot}\kernel{\linA^*}$.
\end{itemize}
\end{theorem1}

\paragraph{} %-- the Fredholm property of closed linear operators
%------------------------------------------------------------------------------%
We have introduced the Fredholm property of linear operators
under the assumption that
these are continuous and everywhere defined in their domain spaces.
If a linear operator $\linA$ is closed instead of continuous,
however,
we define the possible Fredholm property by regarding
$\linA$ as a continuous linear operator $\linA \: \domain{\linA} \xto \banY$,
where the domain of definition $\domain{\linA}$ is endowed
with the \textit{graph norm}
\[
\norm{x}_{\linA} \eqdef \norm{x}_X \+ \norm{\linA x}_Y
\]
for $x \xeof \domain{\linA}$.
The pair $(\domain{\linA},\norm{\,\cdot\,}_{\linA})$
forms a Banach space due to the fact that
the operator $\linA$ is assumed to be closed.

\paragraph{} %-- essential spectrum
%------------------------------------------------------------------------------%
The concept of semi-Fredholm operator can be used to characterize
the \textit{essential spectrum} $\sigma_{ess,1}(\linA)$
of a closed operator $\linA$ acting in a complex Banach space $\banX$:
a number $\lambda \xeof \C$ belongs to $\sigma_{ess,1}(\linA)$
\iff $\linA - \lambda$ is \textit{not} semi-Fredholm.
But there are further approaches to define the notion of essential spectrum:

\begin{itemize}[leftmargin=6mm]
\item[-] 
$\sigma_{ess,2}(\linA) := \{\lambda \xeof \C:
\linA-\lambda \mbox{ is upper semi-Fredholm} \}$;

\item[-] 
$\sigma_{ess,3}(\linA) := \{\lambda \xeof \C:
\linA-\lambda \mbox{ is Fredholm} \}$;

\item[-] 
$\sigma_{ess,4}(\linA) := \{\lambda \xeof \C:
\linA-\lambda \mbox{ is Fredholm with } \chi(\linA-\lambda) \eq 0\}$;

\item[-] 
$\sigma_{ess,5}(\linA)$ is defined to be the union of $\sigma_{ess,1}(\linA)$
and all components of $\C \setminus \sigma_{ess,1}(\linA)$,
which do not intersect with $\rho(\linA)$,
called the \textit{Weyl spectrum} of $\linA$.
\end{itemize}

\subparagraph{} %- properties of the essential spectrum
We then have the chain of inclusions 
\[
\sigma_{ess,1}(\linA) \, \xsubset \,
\sigma_{ess,2}(\linA) \, \xsubset \,
\sigma_{ess,3}(\linA) \, \xsubset \,
\sigma_{ess,4}(\linA) \, \xsubset \,
\sigma_{ess,5}(\linA).
\]
All sets $\sigma_{ess,k}(\linA)$ form closed subsets
of the spectrum $\sigma(\linA)$.
The main property of the essential spectra $\sigma_{ess,k}(\linA)$
($k \eq 1,2,3,4)$ is
that they do not change under compact perturbations of $\linA$,
hence $\sigma_{ess,k}(\linA + \compA) \eq \sigma_{ess,k}(\linA)$
for every compact linear operator $\compA \: \banX \xto \banX$.
Finally,
the set $\sigma_{ess,4}(\linA)$ is the largest subset of $\sigma(\linA)$
such that
\[
\sigma_{ess,4}(\linA + \compA)
\=
\sigma_{ess,4}(\linA)
\]
for every compact linear operator $\compA \xeof \BO{\banX}$.

\paragraph{} %-- discrete spectrum
%------------------------------------------------------------------------------%
On the other hand,
the \textit{discrete spectrum} $\sigma_d(\linA)$ is defined
to be the relative complement of $\sigma_{ess,5}(\linA)$ {\wrt} $\sigma(\linA)$,
hence we have a disjoint decomposition of the spectrum of $\linA$ according to
\[
\sigma(\linA) \eq \sigma_d(\linA) \, \sqcup \, \sigma_{ess,5}(\linA).
\]
The set $\sigma_d(\linA)$ contains those points $\lambda \xeof \sigma(\linA)$,
which are isolated in $\sigma(\linA)$ and the rank of the \textsc{Riesz}
\textit{projection} $P_{\lambda}$ is finite,
where $P_{\lambda}$ is defined by the \textsc{Dunford} integral
\[
P_{\lambda}(x) \eqdef
\frac{1}{2 \pi i} \, \int_{\,\gamma_{\lambda}} \, R(\linA,\zeta) x \, d\zeta
\quad (x \xeof \banX)
\]
and $\gamma_{\lambda}$ is a closed continuous path enclosing $\lambda$
but no other values in $\sigma(\linA)$.
Elements of the set $\sigma_d(\linA)$ are called \textit{normal eigenvalues}.
One may also introduce the set $\sigma_d(\linA)$ as the collection
of isolated points $\lambda \xeof \sigma(\linA)$,
for which $\linA-\lambda$ is Fredholm
(rsp. Fredholm with index $0$ or semi-Fredholm).

%------------------------------------------------------------------------------%
\section{Well-Posedness of Linear Equations}                                   %
%------------------------------------------------------------------------------%
The main subject of the section is the concept of \textit{well-posedness}
that may be defined for the following types of equations
related to a linear operator $\linA$
with domain space $\banX$ and range space $\banY$,
both of them assumed to be topological or normed linear spaces:

\begin{itemize}[leftmargin=12mm]
\item[(P-1)]
the stationary equation $\linA x \eq y$,
which also includes the eigen\-value equation $\linA x - \lambda x \eq y$
for $\lambda \xeof \C$,
\item[(P-2)]
the evolution equation of first order $x' + \linA x \eq y$,
\item[(P-3)]
the evolution equation of second order $x'' + \linA x \eq y$,
\end{itemize}

where in the second and third case $\banY \eq \banX$ is assumed
and $x'$ and $x''$ shall denote the first and second "time-derivative"
of a searched function $x\:[0,\T] \xto \banX$.

\paragraph{} %-- solution concepts for stationary equations
%------------------------------------------------------------------------------%
With respect to problem (P-1),
a linear equation $\linA x \eq y$ is called \textit{well-posed},
if
\textit{(i)}  $\linA$ is a bijective mapping
              between the linear spaces $\domain{\linA}
              \xsubset \banX$ and $\banY$ and
\textit{(ii)} the inverse mapping $\linA^{-1}$ is continuous,
              so small variations of $y$ cause small variations of $x$.

\subparagraph{}
By the closed graph theorem,
closeness and bijectivity of $\linA$ are already sufficient conditions
for well-posedness of (P-1).
If $\linA$ is not bijective,
however,
theorem \ref{BOUNDED_BELOW_THEOREM} has useful propositions
for the study of this equation:
in fact,
if $\linA$ is densely defined, closed and bounded below,
then injectivity of the adjoint $\linA^*$ implies well-posedness,
since in this case $\kernel{\linA^*} \eq \, ^{\bot}\range{\linA}$ holds
and consequently $\range{\linA}$ must be the whole of $\banY$.

\paragraph{} %- solvability of linear equations
We want to complete our considerations on linear equations
with a complete collection of concepts
regarding the solvability of $\linA \, x \eq y$
where again $\linA$ is assumed to be closed and
acting between Banach spaces $\banX$ and $\banY$.
The equation $\linA \, x \eq y$ is called

\begin{itemize}[leftmargin=9mm]
\item[(a)]
\textit{uniquely solvable} (\cite{Kr}),
if $\linA$ is injective;

\item[(b)]
\textit{everywhere solvable} (\cite{Kr}),
if $\linA$ is surjective;

\item[(c)]
\textit{solvable} (\cite{Br1}),
if it is uniquely and everywhere solvable,
or in other words, if $\linA$ is bijective;

\item[(d)]
\textit{densely solvable} (\cite{Kr}),
if $\overline{\range{\linA}} \eq \banY$;

\item[(e)]
\textit{normally solvable} (\cite{Kr}),
if $\overline{\range{\linA}} \eq \range{\linA}$;

\item[(f)]
\textit{correctly solvable} (\cite{Kr}),
if $\linA$ is bounded below;

\item[(g)]
\textit{semi-solvable} (\cite{Br1}), if $\linA$ is Fredholm;

\item[(h)]
\textit{almost solvable} or \textit{regularly solvable} (\cite{Br1}),
if $\linA$ is semi-solvable and $\chi(\linA) \eq 0$;

\item[(i)]
\textit{well-posed},
if it is solvable and $\linA^{-1}$ is continuous,
or equivalently,
if $0 \xeof \rho(\linA)$;

\item[(j)]
\textit{completely solvable} (\cite{Br1}) or \textit{compactly solvable}
if it is solvable and $\linA^{-1}$ is compact.
\end{itemize}

\paragraph{} %-- solution concepts for evolutional equations
%------------------------------------------------------------------------------%
We turn to equations of type (P-2) and (P-3).
By an \textit{evolution equation} is meant the differential equation
describing a dynamical system.
It describes the propagation of one or more physical magnitudes along the time,
where the initial values $x(t_0)$ of them are given
in a start point $t_0 \xeof \R$.

\subparagraph{} %- example of evolution equation
An elementary example of an evolution equation
is given by the differential equation
$x' + \lambda x \eq 0$ for $x(0) \eq x_0 \xeof \R$ and $\lambda {\,>\,} 0$.
It is well known that
this problem is solved by the exponential function
$x(t,x_0) \eq x_0 \, e^{-\lambda t}$ for $t {\,>\,} 0$.

\paragraph{} %- functional analysis and evolution equations
Functional analysis provides a rich theory
on \textit{homogeneous} linear evolution equations of the form
\begin{equation}\label{HOMOGENEOUS_CAUCHY_PROBLEM_I}
\begin{cases}[1.2]
\quad x'(t) \+ \linA \, x(t) \= 0,\\ 
\quad x(0) \= x_0 \xeof \banX,\\
\end{cases}
\end{equation}
where $\linA$ is a densely defined and closed linear operator
acting in a Banach space $\banX$
and with domain of definition $\domain{\linA}$.
Our task is to find \textit{vector-valued} solutions $x \: [0,\T) \xto \banX$
of the equation $x' + \linA x \eq 0$ 
satisfying the initial value condition $x(0) \eq x_0$.

\subparagraph{} %- second order equations
The homogeneous Cauchy problem of \textit{second order} consists of solving
the initial value problem
\begin{equation}\label{HOMOGENEOUS_CAUCHY_PROBLEM_II}
\begin{cases}[1.2]
\quad x''(t) \+ \linA \, x(t) \= \0,\\
\quad x(0)  \= x_0, \\
\quad x'(0) \= x_1,\\
\end{cases}
\end{equation}
where $\linA$ is densely defined and closed.
We also suppose $x_0 \xeof \domain{\linA}$,
but the element $x_1$ may belong to any other linear subspace of $\banX$.

\paragraph{} %- spaces of classical differentiable functions
Let us shortly review a couple of fundamental Banach spaces,
in which the possible solution of a Cauchy problem resides.
If $\banX$ is an arbitrary Banach space and $(a,b) \xsubset \R$,
then a vector-valued function $x \: [a,b] \xto \banX$
is called \textit{classical differentiable} in $t_0 \xeof (a,b)$,
if there is an element $x'(t_0) \xeof \banX$ such that
\[
\lim_{\,t \xto t_0} \,
\norm{ \, \frac{x(t)-x(t_0)}{t-t_0}\, - x'(t_0) \, } \= 0.
\]
Moreover,
the function $x$ is called {classical differentiable in} $(a,b)$,
if $x$ is so in each $t \xeof (a,b)$.
The linear space consisting of functions having \textit{continuous}
classical derivatives in $(a,b)$ is denoted $C^1(a,b;\banX)$.
The spaces $C^k(a,b;\banX)$ ($k \xeof \N$) are defined recursively as usual,
and the space $C^{\infty}(a,b;\banX)$ is defined
by taking the intersection of all spaces $C^k(a,b;\banX)$,
$k \xeof \N_0$.

\subparagraph{} %- homgeneous Cauchy problems
As already mentioned,
the \textit{homogeneous Cauchy problem} of first order associated with
a densely defined and closed linear operator $\linA$ acting in $\banX$
is to find a solution
of the initial value problem \eqref{HOMOGENEOUS_CAUCHY_PROBLEM_I}.
But we first have to define concepts of solution for this:

\begin{itemize}[leftmargin=4mm]
\item[1.]
A vector-valued function
$x \xeof C([0,\T),\domain{\linA}) \, \cap \, C^1((0,\T),\banX)$
is called a \textit{classical solution} of \eqref{HOMOGENEOUS_CAUCHY_PROBLEM_I},
if it satisfies the initial condition $x(0) \eq x_0$
and the differential equation in each point $t \xeof (0,\T)$.

\item[2.]
a continuous function $x \: [0,\T) \xto \banX$ is said to be
a \textit{mild solution} of \eqref{HOMOGENEOUS_CAUCHY_PROBLEM_I} if
\[
\mbox{(1) } \, \tilde{x}(t) \eqdef
\int_{\,0}^t x(s) \, ds \, \xeof \domain{\linA}
\quad \and \quad
\mbox{(2) } \, \linA \, \tilde{x}(t) \eq x(t)-x_0
\]
holds for all $0 {\,<\,} t {\,<\,} \T$,
where the integral on the right-hand side of (1) 
is meant in the Riemannian sense.
One is led to the notion of mild solution in a natural way
by integrating both sides of the equation $x' + \linA \, x \eq \0$,
which then becomes into
\[
\int_{\,0}^t x' ds \=
x(t)-x(0) \= -
\int_{\,0}^t \linA \, x(s) \, ds.
\]
If $\linA$ is closed,
then the action of $\linA$ and performing the integration commutes
for continuous functions $x \: [0,t] \xto \domain{\linA}$,
that is,
\[
\int_{\,0}^t \linA \, x(s) \, ds  \= \linA \, \int_{\,0}^t x(s) \, ds,
\]
and to the end we obtain the equation
\[
x(t) \= x(0) - \linA \, \int_{\,0}^t x(s) \,ds
\quad \quad (t \xeof (0,\T)).
\]
\end{itemize}

\subparagraph{} %- comparing classical and mild solutions
Notice that each classical solution is a mild solution,
and reversely a mild solution is classical
iff it is continuously differentiable in $(0,\T)$.
Thus classical and mild solutions of homogeneous Cauchy problems
only differ {\wrt} this regularity condition.

\paragraph{} %-- correctly and well-posed Cauchy problems
%------------------------------------------------------------------------------%
After having the notions of classical and mild solution for the homogeneous
Cauchy problem available,
we want to introduce the concepts of \textit{correct posedness} and
\textit{well-posedness}.

\subparagraph{} %- definitions of correctly and well-posed problems
First,
the homogeneous Cauchy problem \eqref{HOMOGENEOUS_CAUCHY_PROBLEM_I}
is called \textit{correctly posed},
if there is a unique classical solution $x$
for any initial value $x_0 \xeof \domain{\linA}$.
Second,
it is denoted \textit{well-posed},
if it is correctly posed and
the solution $x$ depends continuously on the initial value $x_0$,
or in other words,
if for each $\T {\,>\,} 0$ there is $c_{\T} {\,>\,} 0$ such that
\[
\norm{x(t)} \leq c_{\T} \norm{x_0}
\]
for all $t \xeof (0,\T)$ and all $x_0 \xeof \banX$.

\subparagraph{}
In \cite{EN}
the continuous dependency of the solutions from the initial values
is expressed as follows:
if $\{x_n\}_{n \xeof \N}$ is a converging sequence in $\domain{\linA}$
with limit element $\0 \xeof \banX$,
then the related solutions $t \mapsto x_n(t)$ are uniformly convergent
to the zero function in all compact intervals $[0,t_0]$
for $0 {\,<\,} t_0 < \T$.
Similar definitions for correct posedness and well-posedness can also be
quoted for the second-order homogeneous Cauchy problem
\eqref{HOMOGENEOUS_CAUCHY_PROBLEM_II}.

\subparagraph{}
Let us finally mention that problem \eqref{HOMOGENEOUS_CAUCHY_PROBLEM_I}
governed by a densely defined and closed linear operator $\linA$ is well-posed
iff
it is correctly posed and
in addition $\rho(\linA) \ne \emptyset$ holds.

\paragraph{} %-- the semigroup property of solution operators
%------------------------------------------------------------------------------%
It is not hard to show
that solutions $t \xeof [0,\T) \mapsto x(t,x_0)$ of the evolution equation
$x' + \linA x \eq 0$ fullfill the \textit{semi\-group property}
\[
x(t_1 + t_2,x_0) \= x(t_2, x(t_1,x_0))
\]
for all $0 \leq t_1 {\,<\,} t_2 {\,<\,} \T$.
Moreover,
if $x_1$ and $x_2$ are the solutions with initial values $x_1$ and $x_2$,
then $x_1+x_2$ is the solution with initial value $x_1+x_2$,
so the solution mapping $x(0) \mapsto x(t)$ for $t \xeof (0,\T]$
is given by a linear operator $S(t)$ acting in $\banX$,
which is actually continuous.
Hence we are faced with families $\{S(t)\}_{t \geq 0}$
of bounded linear operators $S(t) \: \banX \xto \banX$ satisfying
\[
S(t_1+t_2) = S(t_1) \circ S(t_2) \quad \quad (0 \leq t_1 < t_2 < \T),
\]
which lead us to the theory of {strongly continuous semigroups}.

\subparagraph{} %- strongly continuous semigroups
Given an arbitrary Banach space $\banX$,
a \textit{strongly continuous semigroup of bounded linear operators}
(shortly denoted a \textit{C$_0$-semigroup}) is an operator-valued mapping
\[
\linS \: [0,\infty) \xto \BO{\banX},
\]
also denoted $\{\linS(t)\}_{t \geq 0}$,
which distinguishes by the properties

\begin{itemize}[leftmargin=12mm]
\item[(G-1)]
$\linS(t+s) \eq \linS(t) \circ \linS(s)$ for all $t,s \geq 0$
(the \textit{semigroup property});

\item[(G-2)]
$\lim_{\,t \downarrow 0} \, \linS(t)\,x \eq x$ for all $x \xeof X$.
\end{itemize}

\subparagraph{} %- continuity from the right
It follows from (G-2),
which is nothing but strong continuity from the right-hand side in $t \eq 0$,
that $\linS(0)$ is the identity map in $\banX$.
The conditions (G-1) and (G-2) together
imply strong continuity from the right of the semigroup
even in every $t {\,>\,} 0$.
In case of $\banX \eq \R$,
the family $\{\cal{E}(t)\}_{\, t \geq 0}$ of exponential functions
$\cal{E}(t) \: x \xeof \R \mapsto e^{\lambda tx}$
fulfills the properties (G-1) and (G-2),
so $\{\cal{E}(t)\}_{\, t \geq 0}$ serves as an elementary example
of a C$_0$-semigroup.

\subparagraph{} %- orbit maps
By the \textit{orbit map} or \textit{trajectory} $\eta_x$
of a strongly continuous semigroup $\{\linS(t)\}_{t \geq 0}$
and associated with initial value $x \xeof \banX$
we mean the vector-valued mapping $\eta_x \: [\,0,\infty) \xto \banX$
given by $\eta_x(t) := S(t) \, x$.
The conditions (G-1) and (G-2) together
imply continuity of the function $\eta_x$.
But $\eta_x$ is actually differentiable under a very mild assumption:
if $\{\linS(t)\}_{t \geq 0}$ is a strongly continuous semigroup,
then for each $x \xeof \banX$
the mapping $\eta_x$ is differentiable in $(0,\infty)$
\iff
$\eta_x$ is right-differentiable in $t \eq 0$.

\paragraph{} %- infinitesimal generators of strongly continuous semigroups
The last statement forms the motivation for the following concept:
if $\{\linS(t)\}_{t \geq 0}$ is a strongly continuous semigroup,
its \textit{infinitesimal generator}
(briefly denoted the \textit{generator} of the semigroup)
is defined as a mapping $\linA$ in $\banX$ by means of the graph
\[
\Gamma_{\linA} \eqdef
\{\,(x,y) \xeof X \xtimes X: \;
y \= \lim_{\,h \downarrow 0} \, \frac{1}{h}(\linS(h)x-x) \mbox{ exists} \}.
\]
Hence the domain of definition $\domain{\linA}$ consists
of those elements $x \xeof \banX$,
for which the orbit map $\eta_x$ is differentiable from the right in $t \eq 0$.
But this means that
$\eta_x$ is also (left and right) differentiable in every $t>0$,
which is a consequence of property (G-1).

\subparagraph{} %- properties of infinitesimal generators
It is easy to show that
the infinitesimal generator of a strongly continuous semigroup is linear,
so it constitutes a linear operator.
Two of the basic properties of such an operator are given by the fact
that it is closed and densely defined in $\banX$,
hence infinitesimal generators of strongly continuous semigroups fit
into our base class of linear operators considered in section 1.

\paragraph{} %-- well-posedness theorem for first order Cauchy problems
%------------------------------------------------------------------------------%
The following fundamental theorem describes the well-posedness of the
homogeneous Cauchy problem in terms of properties of $\linA$:

\begin{theorem1}\label{WELL_POSEDNESS_THEOREM_I}
Let $\linA$ be a densely defined and closed linear operator
acting in a Banach space $\banX$.
Then the problem \eqref{HOMOGENEOUS_CAUCHY_PROBLEM_I} is well-posed
\iff
the operator ${-}\linA$ is the infinitesimal generator
of a strongly continuous semigroup.
\end{theorem1}

\subparagraph{}
In fact,
the family $\{S(t)\}_{t \geq 0}$ of linear operators
having ${-}\linA$ as its infinitesimal generator
describes the classical solutions $x$ with initial values
$x_0 \xeof \domain{\linA}$ of \eqref{HOMOGENEOUS_CAUCHY_PROBLEM_I},
that is,
$x(t,x_0) \eq S(t)\,x_0$ for $0 \leq t {\,<\,} \T$.
A proof of this well-posedness theorem can be found in \cite{EN},
see theorem II.6.6.

\subparagraph{}
The conditions on ${-}\linA$ to be the infinitesimal generator
of a C$_0$-semigroup,
however,
can be expressed 
\textit{(i)}  by the location of the spectrum $\sigma(\linA)$
              in the complex plane and
\textit{(ii)} by a collection of \textit{resolvent estimates}
              for the operator $\linA$.
We shall postpone a detailed review of these conditions to section 5.

\paragraph{} %-- inhomogeneous Cauchy problems
%------------------------------------------------------------------------------%
To the end of this section we want to deal with the solution theory
of the \textit{inhomogeneous} abstract Cauchy problem
\begin{equation}\label{INHOMOGENEOUS_CAUCHY_PROBLEM}
\begin{cases}[1.2]
\quad  x'(t) \+ \linA \, x(t) \= y(t)\\
\quad  x(0) \= x_0 \xeof \banX, \\
\end{cases}
\end{equation}
where $\banX$ is a Banach space,
$y \: [0,\T) \mapsto \banX$ is a given function and
the solution $u \: (0,\T) \xto \banX$ of \eqref{INHOMOGENEOUS_CAUCHY_PROBLEM}
is searched.

\subparagraph{} %- Duhamels principle
Let us first remember \textsc{Duhamel}\textit{'s principle},
also called the \textit{method of variation of constants}:
if $x$ is a \textit{classical solution} of \eqref{INHOMOGENEOUS_CAUCHY_PROBLEM}
(for the definition of this notion see below),
the right-hand side $y \: [0,\T) \xto \banX$ is continuous
and the linear operator $\linA \: \banX \xto \banX$ is bounded,
then 
\[
x(t) \= \linS(t) \, x_0 \+ \int_{\,0}^t \linA \, y(s) \, ds,
\quad (0 < t \leq \T)
\]
holds,
where $\linS(t) \xeof \BO{\banX}$ is the solution operator
of the abstract homogeneous problem \eqref{HOMOGENEOUS_CAUCHY_PROBLEM_I},
mapping an initial value $x_0 \xeof \banX$
to the state $\linS(t)x_0 \xeof \banX$ at time $t {\,>\,} 0$.

\subparagraph{} %- solvability of the inhomogeneous Cauchy problem
The main problem in solving the inhomogeneous Cauchy problem is to find
conditions on the function $y$,
which ensure the existence and uniqueness of a solution
of \eqref{INHOMOGENEOUS_CAUCHY_PROBLEM}. 
Hence we first review the notion of classical and mild solution
for inhomogeneous Cauchy problems
and assume that $y \xeof L((0,t),\banX)$ for each $0 {\,<\,} t \leq \T$,
where $L^1(\J,\banX)$ denotes the space
of \textit{Bochner-integrable} functions\footnote{
\,See also section 8,
where a short glimpse on Bochner integrals is presented.}
$f :\ J \xsubset \R \xto \banX$
for any interval $\J \xsubset \R$,
that is,
$f$ is measurable in the sense of a vector-valued mapping
and the real-valued function $t \xeof \J \mapsto \norm{f(t)}_X$
is Lebesgue-integrable.

\paragraph{} %- classical and mild solution for inhomogeneous Cauchy problems
With respect to problem \eqref{INHOMOGENEOUS_CAUCHY_PROBLEM},
a mapping $x \: [0,\T] \xto \banX$ is called

\begin{itemize}[leftmargin=4mm]
\item[1.]
\textit{classical solution},
if
\textit{(i)}   it belongs to $C([0,\T),\banX) \cap C^1((0,\T),\banX)$,
\textit{(ii)}  it satisfies the differential equation
               $x'(t) \,{+}\, \linA \, x(t) \eq y(t)$
               in every $t \xeof (0,\T)$ and
\textit{(iii)} fulfills the initial condition $x(0) \eq x_0$.

\item[2.]
\textit{mild solution},
if
\textit{(i)}  the related homogeneous problem is well-posed
              with solution operator
              \[
                (t,x_0) \xeof (0,\T) \xtimes \domain{\linA} \mapsto S(t) \, x_0
              \]
              and
\textit{(ii)} the function $x$ fulfills \textsc{Duhamel}'s formula
\begin{equation}\label{DUHAMEL_FORMULA}
x(t) \= \linS(t) \, x_0 \+ \int_{\,0}^t \linS(t-s) \, y(s) \, ds
\quad (0 < t \leq \T).
\end{equation}
\end{itemize}

\subparagraph{} %- relationship between mild and classical solution
While in the homogeneous case a mild solution is classical
iff it is continuously differentiable,
the validness of this assertion remains true for the inhomogeneous problem
only
under some additional regularity properties of the right-hand side $y$
(see for example \cite{Mi}, section 4.7).

\subparagraph{}
In general,
a mild solution may not be differentiable and it may not lie in the domain
of definition $\domain{\linA}$,
hence it may not be a classical solution.
However,
if a classical solution exists,
then the mild solution is given by formula \eqref{DUHAMEL_FORMULA}.

%------------------------------------------------------------------------------%
\section{Linear Operators in Hilbert Space}                                    %
%------------------------------------------------------------------------------%
In this section we want to review a couple of specific linear operators
that occur in the setting of Hilbert spaces.
The theory on them is extremely helpful in the analysis
of liner-elliptic differential operators acting in Hilbert space $\Lh$,
where $\G$ denotes an open and connected subset of an Euclidean spaces $\Rd$
(see section 9 for more referring to this).

\paragraph{} %-- inner products and Hilbert spaces
%------------------------------------------------------------------------------%
Recall that
a complex Hilbert space $\hilH$ differs from an arbitrary Banach space
by the presence of an \textit{inner product},
which is nothing but a symmetric and positive sesquilinear form on $\hilH$.
This form defines the norm and the metric topology,
but also determines the inner geometry to be close
to that of an Euclidean space.

\subparagraph{} %-- inner geometry of Hilbert spaces
%------------------------------------------------
Let us make these assertions precise.
Given a real or complex linear space $\hilH$,
an \textit{inner product} in $\hilH$ is defined to be a mapping
$\scalar{\,\cdot\,}{\,\cdot\,}\: \hilH \xtimes \hilH \xto \F$,
which is

\begin{itemize}[leftmargin=35pt]
\item[(I-1)]
linear in the first argument,

\item[(I-2)]
either linear or antilinear in the second argument (conditional upon,
whether $\F$ is the real or the complex number field),
\end{itemize}
together with the additional properties

\begin{itemize}[leftmargin=35pt]
\item[(I-3)]
$\scalar{y}{x} \eq           \scalar{x}{y}$  ($\F=\R$) rsp.
$\scalar{y}{x} \eq \overline{\scalar{x}{y}}$ ($\F=\C$),
\item[(I-4)]
$\scalar{x}{x} \geq 0$ for all $x \neq 0$,
\item[(I-5)]
$\scalar{x}{x} \eq 0$ implies $x \eq 0$.
\end{itemize}

\subparagraph{}
Property (I-3) is denoted the \textit{Hermitean symmetry}
of $\scalar{\,\cdot\,}{\,\cdot\,}$,
while (I-4) and (I-5) together mean \textit{positive definiteness}.
The presence of the inner product gives rise to define a norm in $\hilH$
according to
\begin{equation}\label{DERIVED_NORM}
\norm{x} \, \eqdef \, \scalar{x}{x} ^{1/2} \quad \quad (x \in \hilH),
\end{equation}
so each linear space furnished with an inner product is normed
in a natural manner.
The inner product and the norm are coupled by
the \textit{Cauchy-Schwarz} \textit{inequality}
\[
\abs{\scalar{x}{y}} \, \leq \,
\norm{x} \cdot \norm{y} \quad \quad (x,y \in \hilH),
\]
which implies that the mapping
$\scalar{\,\cdot\,}{\,\cdot\,}\:\hilH \xtimes \hilH \xto \F$ is continuous.
An inner product space $\hilH$ is called a \textit{Hilbert space}
or shortly \textit{Hilbertian},
if it is complete {\wrt} the norm defined by \eqref{DERIVED_NORM}.

\subparagraph{} %- parallelogram equation and polarization formula
The pleasant geometric properties of inner product spaces
may be attributed to the validness of the \textit{parallelogram equation}
\begin{equation}\label{PARALLELOGRAM_EQUATION}
2 \norm{x}^2 + 2 \norm{y}^2 \= \norm{x+y}^2 + \norm{x-y}^2
\qfa x,y \in \hilH.
\end{equation}
In general,
the norm $\norm{\,\cdot\,}$ on a linear space $\banX$
can be induced by an inner product
\iff \eqref{PARALLELOGRAM_EQUATION} holds.
In this case and if $\banX$ is a complex linear space,
the inner product can be recovered from the norm
by the \textit{polarization formula}
\[
\scalar{x}{y}
\=
\frac{1}{4} \,
(\,
\norm{x + y}^2 - \norm{x-y}^2 +i \norm{x+iy}^2 - i \norm{x-iy}^2
\,).
\]
If $\banX$ is a real space,
however,
the polarization formula simplifies to
\[
\scalar{x}{y} \=
\frac{1}{4} \,
(\,
\norm{x + y}^2 - \norm{x-y}^2
\,).
\]

\paragraph{} %- properties of Hilbert spaces
In the following we want to point out
that Hilbert spaces differ from general Banach spaces
by a couple of desirable properties:

\begin{itemize}[leftmargin=4mm,itemsep=6pt]
\item[1.] \textit{Orthogonality}.
Two elements $x,y \in \hilH$ are called \textit{orthogonal} to each other,
if $\scalar{x}{y} \eq 0$ (notation: $x \bot y$).
Any subset $\orthO \xsubset \hilH$ is said to be \textit{orthogonal}
or an \textit{orthogonal system},
if all elements in $\orthO$ are orthogonal to each other.
Finally,
the \textit{orthogonal space} $\setA^{\bot}$
of a subset $\setA \xsubset \hilH$ contains exactly the elements $y \in \hilH$,
which are orthogonal to all $x \in \setA$.
Such sets basically form closed linear subspaces of $\hilH$.

\item[2.] \textit{Projection Theorem}.
Let $\setA \xsubset \hilH$ be a \textit{closed} linear subspace.
Then for each $x \in \hilH$ there is exactly one $x_0 \in \setA$
such that
\[
\norm{ x - x_0 } \= \inf \, \{ \, \norm {x-y}: y \in \setA \,\}.
\]
The element $x_0$ is called the \textit{projection} of $x$ onto $\setA$
and minimizes the distance between $\setA$ and $x$.
The associated linear mapping
$\pi_{\setA}\:x \in \hilH \mapsto x_0 \in \setA$
is also denoted the \textit{projection} of $\hilH$ onto $\setA$.
In general,
by a \textit{projection} is meant
any linear mapping $\pi\:\hilH \xto \hilH$ satisfying
$\pi \circ \pi \eq \pi$ and
$\scalar{\pi x}{y} \eq \scalar{x}{\pi y}$ for all $x,y \in \hilH$.

\item[3.] \textit{Decomposition Theorem}.
Using the notation of orthogonal space $\setA^{\bot}$,
the projection of an element $x$ onto $\setA$ is characterized as follows:
$x_0 \eq \pi_{\setA} x$ iff $x-x_0 \in \setA^{\bot}$.
Hence each $x \in \hilH$ can be uniquely represented by the formula
$x \eq x_0 + x^{\bot}$,
where $x_0 \in \setA$ and $x^{\bot} \in \setA^{\bot}$,
so the entire space $\hilH$ is isometrically isomorphic
to the direct sum $\setA \, \oplus \, \setA^{\bot}$
of the closed subspaces $\setA$ and $\setA^{\bot}$.
This fact is known as \textit{decomposition theorem} in Hilbert spaces.
In other words,
each closed linear subspace $\setA$ of $\hilH$
is \textit{topologically complementable},
that is,
to each closed subspace $\setA$
there is a second closed subspace $\setB \xsubset \hilH$
such that $\setA \cap \setB \eq \{\0\}$ and $[\setA \cup \setB] \eq \hilH$.
Of course,
a closed subspace $\A$ is topologically complemented
by its orthogonal space. %$\setA^{\bot}$.

\item[4.] \textit{Existence of Orthonormal Bases}.
An orthogonal system $\orthO \xsubset \hilH$ is denoted
an \textit{orthonormal system},
if $\scalar{x}{x} \eq 1$ for all $x \in \orthO$.
However,
if $\orthO$ contains only \textit{finitely} many elements,
then the span $[\orthO]$ forms a \textit{closed} subspace in $\hilH$,
and the projection $\pi_{[\orthO]}$ of any element $x$ onto $[\orthO]$
can be expressed by the formula
\[
\pi_{[\orthO]} \, x \= \sum_{i=1}^n \, \scalar{x}{e_i} \, e_i.
\]
However,
if $\orthO$ consists of arbitrarily many elements,
the inequality $\scalar{x}{e_{\alpha}} \neq 0$ can only be valid
for mostly countable many elements $e_{\alpha} \in \orthO$,
which in conclusion enables to add the coefficients $\scalar{x}{e_{\alpha}}$
together about all elements of $\orthO$.

\subparagraph{} %- orthonormal bases
An orthonormal system $\orthO$ is called
an \textit{orthonormal base} in $\hilH$,
if its span $[ \orthO ]$ lies dense in $\hilH$.
In this case we may represent each $x \in \hilH$ according to
\begin{equation}\label{FOURIER_COEFFICIENTS}
x \= \sum_{e_{\alpha} \in \, \orthO} \, \scalar{x}{e_{\alpha}} \, e_{\alpha},
\end{equation}
which is primarily a formal sum,
but consists of mostly countable many summands.
The inner products $\scalar{x}{e_{\alpha}}$ are called
the \textit{Fourier coefficients} of $x$ {\wrt} $\orthO$.
The right hand side of \eqref{FOURIER_COEFFICIENTS}
consists of \textit{mostly countable many} non-zero summands,
and this series is invariant
by realignments of the elements $e_{\alpha} \in \orthO$.

\subparagraph{}
There is of course the need to determine,
whether a given orthonormal system $\orthO$ forms an orthonormal base.
In fact,
it can be shown that
the following statements are equivalent:

\begin{itemize}[leftmargin=9mm]
\item[(a)]
$\orthO$ is an orthonormal base in $\hilH$.

\item[(b)]
$\orthO$ is \textit{complete},
that is,
$\scalar{x}{e} \eq 0$ for all $e \in \orthO$ implies $x \eq \0$.

\item[(c)]
$\sum_{e_{\alpha} \in \,\orthO} \;
\abs{ \scalar{x}{e_{\alpha}} } ^2 \eq \norm{x}^2$ for all $x \in \hilH$
(\textit{Parseval}\textit{'s equation}).
\end{itemize}

\paragraph{} %- existence of orthonormal bases
It is also nearby to ask for the existence of orthonormal bases.
First of all,
the collection of all orthonormal systems in a Hilbert space $\hilH$
forms a partial ordered set,
and under the assumption that \textsc{Zorn}{'s lemma}\,\footnote{\,
\,The \textit{axion of choice} claims that
if $\cal{F}$ is a disjoint family $F_{\alpha}$ ($\alpha \in \J$)
of non-empty sets,
then there is a set $A$ that shares exactly one element
with each other set $F_{\alpha} \in \cal{F}$.}
is valid,
it contains a \textit{maximal} element $\orthO$
that cannot be increased,
hence each Hilbert space $\hilH \neq \{\0\}$ owns an orthonormal base.

\subparagraph{} %- Hilbert space dimension
Furthermore,
if $\orthO_1$ and $\orthO_2$ are orthonormal bases in $\hilH$,
then there is a bijective mapping between them,
thus they own identical cardinal numbers.
This legitimates us to introduce the notation
of \textit{Hilbert space dimension} $\abs{\hilH} $
being the cardinal number of any orthonormal base in $\hilH$.

\subparagraph{} %- separable Hilbert spaces
A Hilbert space $\hilH$ is called \textit{separable}
if there is a mostly countable subset lying dense in $\hilH$.
It turns out that $\hilH$ is separable
\iff
$\abs{ \hilH} \leq \abs{\N} $.
Indeed,
each orthonormal base in $\hilH$ forms a mostly countable set in this case,
and $\hilH$ proves to be norm-isomorphic to the Hilbert space $\ell^2$
of square summarizable sequences $f\:\N \xto \F$.

\item[5.] \textit{Riesz's Representation Theorem}.
It is one of the basic tasks in functional analysis to identify
the adjoint space $\banX^*$ of a given normed space $\banX$
with another known Banach space $\banY$
by finding a suitable norm-isomorphism $\Phi\:\banX^* \xto \banY$.
This can easily be done if $\banX \eq \hilH$ is a real or complex Hilbert space,
as we shall point out in what follows.

\subparagraph{} %- conjugate spaces
But when dealing with a \textit{complex} Hilbert space $\hilH$,
it is convenient to consider the linear space $\hilH_a^*$ of
\textit{continuous} and \textit{conjugate linear} functionals
$\phi\:\hilH \xto \C$,
which distinguish by
\[
\phi(\lambda \cdot x) \eq \overline{\lambda} \, \phi(x)
\; \fa x \in \hilH \and \lambda \in \C.
\]
The space $\hilH_a^*$ is furnished with the norm
\[
\norm{\phi} \, \eqdef \, \sup_{x \in \hilH,\,\norm{x} \eq 1} \, \abs{\phi(x)}
\]
and often called the \textit{conjugate space} of $\hilH$.
Due to the properties of the inner product,
each element $x \in \hilH$ gives rise to
a continuous and conjugate linear functional $\phi_x \in \hilH_a^*$ according to
\[
\phi_x(y) \, \eqdef \, \scalar{x}{y} \qfa y \in \hilH.
\]
The mapping $\Phi\:\hilH \xto \hilH_a^*$ defined by $x \mapsto \phi_x$
is linear,
called the \textsc{Riesz} \textit{map} of $\hilH$.
Since the inner product is positive definite,
$\Phi$ turns out to be injective.

\subparagraph{} %- Riesz's representation theorem
\textsc{Riesz}\textit{'s representation theorem} claims
that the mapping $\Phi$ is actually surjective,
so there is a canonical correspondence
between the elements of $\hilH$ and those of $\hilH_a^*$:

\begin{theorem1}\label{RIESZ_REPRESENTATION_THEOREM}
If $\hilH$ is a complex Hilbert space,
then $\Phi\: \hilH \xto \hilH_a^*$ is bijective and norm-preserving,
{\ie}
$\norm{\Phi(x)} \eq \norm{x}$ for $x \in \hilH$.
\end{theorem1}

\subparagraph{}
This theorem implies that
to each $\phi \in \hilH_a^*$
there is a unique element $x_{\phi} \in \hilH$,
which fulfills the equation
$\scalar{x_{\phi}}{y} \eq \phi(y)$ for every $y \in \hilH$.
Moreover,
the \textsc{Riesz} map gives rise to introduce an inner product
on the conjugate space $\hilH_a^*$ by
\[
\scalar{\phi}{\psi} \, \eqdef \, \scalar{\Phi^{-1}(\phi)}{\Phi^{-1}(\psi)}
\quad \quad (\phi,\psi \in \hilH_a^*),
\]
hence $\Phi$ is not only norm-preserving,
but also \textit{isometric},
that is,
it preserves the inner products $\scalar{x}{y}$ for $x,y \in \hilH$.

\subparagraph{} %- the Riesz map in real Hilbert spaces
On the other hand,
if $\hilH$ is a \textit{real} Hilbert space,
we may replace $\hilH_a^*$ by $\hilH^*$ and verify
that the \textsc{Riesz} map $\Phi$ is a norm-isomorphism
between $\hilH$ and $\hilH^*$ as well,
since in this case
the inner product in $\hilH$ is also linear in the second argument.
\end{itemize}

\paragraph{} %-- Hilbert adjoint
%------------------------------------------------------------------------------%
We now want to elaborate
that in the Hilbert space setting
the adjoint $\linA^*$ of a densely defined linear operator $\linA$
can be seen as a mapping acting in $\hilH$ as well.
More precisely,
if $\linA$ is linear and densely defined in $\hilH$,
then the domain of definition $\domain{\linA^*}$
of the adjoint $\linA^*$ consists of those elements $y \xeof \hilH$,
for which there is $y^* \xeof \hilH$ fulfilling
\[
\scalar{\linA \,x}{y} \eq \scalar{x}{y^*} \fa x \xeof \domain{\linA}.
\]
We then define $\linA^* y := y^*$.
Moreover,
if $\linA$ is closable,
then $\linA^*$ is also densely defined,
hence
$\linA^{**} := (\linA^*)^*$ exists and $\linA^{**} \eq \overline{\linA}$
holds.

\paragraph{} %-- numerical range and external field of regularity
%------------------------------------------------------------------------------%
Besides the spectrum $\sigma(\linA)$,
the \textit{numerical range} $\Theta(\linA)$ plays an important role
in the general theory of Hilbert space operators and is defined by
$\Theta(\linA) :=
\{ \scalar{\linA}{x}: x \xeof \domain{\linA}, \norm{x} \eq 1 \}$.
It is a well-known fact,
given by the \textit{Toeplitz-Hausdorff theorem},
that $\Theta(\linA)$ is a convex set within the complex plane.
Moreover,
let
\[
\Delta(\linA) \eqdef \C \setminus \overline{\Theta(\linA)}
\]
be the \textit{external field of regularity} of $\linA$.
If $\Delta(\linA) \neq \emptyset$,
then it consists of one or two connected components.
To be precise,
if $\Delta(\linA)$ has two connected components,
then the set $\Theta(\linA)$ is either a strip or a ray,
otherwise it forms a half-strip or a half-ray.

\subparagraph{}
For $\zeta \xeof \Delta(\linA)$ we of course have
$\scalar{\linA x}{x} \neq \zeta$
for $x \xeof \domain{\linA}$ and $\norm{x} \eq 1$,
and the following chain of inequalities
\[
0 \, < \, \mbox{dist}(\zeta,\linA) \, \leq \,
\abs{\scalar{\linA x}{x} - \zeta} \=
\abs{\scalar{(\linA-\zeta)x}{x}} \, \leq \,
\norm{(\linA -\zeta) x}
\]
is valid,
where $\mbox{dist}(\zeta,\linA)$ denotes the distance
between the singleton $\{\zeta\} \xsubset \Delta(\linA)$
and the set $\overline{\Theta(\linA)}$,
which in fact is positive,
since these sets are closed and have empty intersection.
Hence the operator $\linA -\zeta$ is bounded below,
which implies
that $\linA-\zeta$ is injective and has closed range.
In other words,
the set $\Delta(\linA)$ is a subset of the \textit{field of regularity}
$\Pi(\linA) := \{\zeta \xeof \C: \linA-\zeta \mbox{ is bounded below}\}$.

\paragraph{} %- classes of linear operators
The shapes and locations 
of the spectrum and the numerical range of
a densely defined and closed linear operator
$\linA \: \domain{\linA} \xsubset \hilH \xto \hilH$
give rise to introduce a couple of subclasses of such operators,
and the following enumeration forms a rough overview of them.
In fact,
the operator $\linA$ is called

\begin{itemize}[leftmargin=4mm,itemsep=2pt]
\item[1.]
\textit{accretive},
if $\Theta(\linA) \xsubset \C^+ := \{\zeta \xeof \C\: \re \zeta \geq 0\}$,
and $\linA$ is said to be \textit{m-accretive},
if it is accretive and does not have any proper accretive extension,
or equivalently,
$\linA$ is accretive and $(-\infty,0) \xsubset \rho(\linA)$;

\item[2.]
\textit{sectorial},
if there is $\psi \xeof [0,\pi)$ and $\omega \xeof \R$ such that

\begin{itemize}[leftmargin=9mm]
\item[(S-1)]
the spectrum $\sigma(\linA)$ is contained in the closure
$\overline{\Sigma_{\,\psi,\omega} }$ of a right-open \textit{sector}
\[
\Sigma_{\,\psi,\omega} \eqdef
\{\zeta \xeof \C: \zeta \neq \omega, arg(\zeta-\omega) < \psi\}
\]
with vertex $\omega$ and semi-angle $\psi$;
this implies that 
$\rho(\linA)$,
the resolvent set of $\linA$,
contains the complementary and left-open sector
$\Lambda_{\,\pi-\psi,\omega} :=
\C \setminus \overline{ \Sigma_{\,\psi,\omega} }$;

\item[(S-2)]
for every $\psi {\,<\,} \phi {\,<\,} \pi$
the following resolvent estimation holds:
\begin{equation}\label{RESOLVENT_ESTIMATE_II}
\sup \;
\{\, \norm{\,\zeta R(\linA,\zeta)\,}:
  \zeta \xeof \Lambda_{\,\pi-\psi,\omega}
\}
\, < \, \infty\,;
\end{equation}
\end{itemize}

\item[3.]
\textit{symmetric} if $\Theta(\linA) \xsubset \R$,
and \textit{half-bounded} if $\linA$ is symmetric
and $\Theta(\linA) \xsubset (m,\infty)$ for some $m \xeof \R$;
\item[4.]
\textit{self-adjoint},
if $\sigma(\linA) \xsubset \R$,
and $\linA$ is said to be \textit{essentially self-adjoint},
if it is closable and the closure $\overline{\linA}$ is self-adjoint;
\item[5.]
to have \textit{pure point spectrum}
if each element $\lambda \xeof \sigma(\linA)$ is an eigen\-value
of $\linA$ with finite-dimensional eigenspace $E_{\lambda}$.
\end{itemize}

\paragraph{}
For the rest of this section
we want to go into the details of these kinds of linear operators.

\begin{itemize}[leftmargin=4mm,itemsep=5pt]

\item[1.] %-- accretive operators
%------------------------------------------------------------------------------%
Let the linear operator $\linA$ acting in a Hilbert space $\hilH$ be accretive.
One can verify that
$\re \scalar{\linA x}{x} \geq 0$ for $x \xeof \domain{\linA}$ is equivalent to
\[
\norm{(\linA - \xi)x} \, \geq \, \norm{\xi x}
\]
for all $\xi {\,<\,} 0$.
But this means that the operators $\linA - \xi$ are bounded below,
that is, 
they are injective and own closed range.

\subparagraph{}
In addition,
if $\linA$ is m-accretive,
then the operators $\linA - \xi$ are also bijective for $\xi {\,<\,} 0$.
Hence we have the estimates
\[
\norm{\, \xi R(\linA,\xi)\,} \, \leq \, 1 \quad \quad (\xi < 0),
\]  
where $R(\linA,\xi)$ again denotes the resolvent
for $\xi \xeof \rho(\linA)$.
Moreover,
it turns out that if $\linA$ is m-accretive,
then it is densely defined and closed,
so the operator fits into our main class of operators considered in section 1.
The usage of m-accretive operator will be illustrated in the next section
when dealing with various types of strongly continuous semigroups of operators.

\item[2.] %-- sectorial operators
%------------------------------------------------------------------------------%
We next discuss the property of sectorialness for a linear operator.
The operator $\linA$ is called \textit{sectorial},
if (S-1) and (S-2) holds for some pair $(\psi,\omega) \xeof [0,\pi) \xtimes \R$.
Notice that
if $\linA$ is sectorial {\wrt} the angle $\psi$,
then there is $0 {\,<\,} \epsilon {\,<\,} \psi$
such that $\linA$ is also sectorial {\wrt} the smaller angle $\psi - \epsilon$.
The infimum $\psi_0$ of all these angles is called
the \textit{angle of sectoriality} of $\linA$.
It should be clear that each sectorial operator is quasi-m-accretive
if the resolvent estimate \eqref{RESOLVENT_ESTIMATE_II} holds with $M \leq 1$.

\paragraph{} %- criterion for sectorialness
The following theorem is useful to ensure sectorialness:

\begin{theorem1}\label{CRITERION_FOR_SECTORIALNESS}
Let $\linA$ be a linear operator in $\banX$ with domain of definition
$\domain{\linA} \xsubset \banX$ having the following properties:

\begin{itemize}[leftmargin=9mm]
\item[(a)]
the resolvent set $\rho(\linA)$ contains
$\C_{\omega} := \{\zeta \xeof \C: \re \zeta \leq \omega\}$
for some $\omega \xeof \R$;

\item[(b)]
there is $M {\,>\,} 0$ such that
$\norm{\zeta \, R(\linA,\zeta)} \, \leq \, M$ for all $\zeta \xeof \C_{\omega}$.
\end{itemize}
Then $\linA$ is sectorial with vertex $\omega$.
\end{theorem1}

\subparagraph{} %- proof of the criterion for sectorialness.
The proof of this theorem can be found in \cite{Lu}, p.\,43.

\paragraph{} %- properties of sectorial operators in Banach spaces
We next collect a couple of properties of a sectorial operator $\linA$
acting in a Hilbert space $\hilH$:

\begin{itemize}[leftmargin=7mm]
\item[1.]
$\linA$ is densely defined and closed.
\item[2.]
$\hilH \eq \kernel{\linA} \oplus \overline{\range{\linA}}$
(this even holds when $\linA$ is acting in a reflexive Banach space $\banX$).
\end{itemize}

\subparagraph{} %- functional calculus for sectorial operators
We also we want to sketch a functional calculus for sectorial operators $\linA$,
which especially enables us to define the family of bounded linear operators
$\{e^{-t \linA}\}_{t \geq 0}$.
In fact,
if $\linA$ is sectorial,
then one may introduce the mapping $t \mapsto e^{-t \linA}$
by means of %the contour integrals
\begin{equation}\label{CONTOUR_INTEGRAL}
e^{-t \linA} x \eqdef
\frac{1}{2 \pi i}
\int_{\gamma} \, e^{-t \zeta} \, R(\linA,\zeta)x \, d\zeta
\quad \quad (x \xeof \hilH),
\end{equation}
which generalizes the well-known \textsc{Riesz-Dunford} \textit{formula}.
Here by $\gamma$ is meant a curve in the complex plane
that drifts within the sector $S_{\phi,\omega}$
and includes the spectrum of $\linA$.

\subparagraph{}
The family of operators $\{e^{-t \linA} \}_{t \geq 0}$
given by \eqref{CONTOUR_INTEGRAL} constitutes a semigroup,
which is strongly continuous
if $\linA$ is assumed to be densely defined.
But this semigroup has some more properties than an arbitrary C$_0$-semigroup
with consequences to the solution theory of homogeneous Cauchy problems
(see section 5 for more details referring to this).

\item[3.] %-- symmetric and half-bounded operators
%------------------------------------------------------------------------------%
As already mentioned,
a linear operator $\linA$ acting in $\hilH$
and with domain of definition $\domain{\linA}$ is called \textit{symmetric},
if the numerical range $\Theta(\linA)$ is a subset of the real line,
and a symmetric operator $\linA$ is called \textit{half-bounded},
if $\Theta(\linA)$ is even a subset of a half-line.

\subparagraph{}
Symmetric operators distinguish from other ones by the fact
that all eigen\-values are real and
eigenvectors belonging to different eigen\-values
are \textit{orthogonal} to each other.
But symmetry of a linear operator $\linA$ can also be characterized
by other properties:

\begin{itemize}[leftmargin=9mm]
\item[(a)]
$\linA$ is symmetric \iff
$\scalar{\linA x}{y} \eq \scalar{x}{\linA y}$
holds for all $x,y \xeof \domain{\linA}$.
This condition can be proved by computing the values
$\scalar{\linA (x+y)}{x+y}$ and $\scalar{\linA (ix+y)}{ix+y}$
for arbitrary $x,y \xeof \hilH$ and checking the imaginary parts of them.
\item[(b)]
If $\linA$ is densely defined,
then the adjoint $\linA^*$ is well-defined,
and $\linA$ proves to be symmetric \iff $\linA \prec \linA^*$.
Since $\linA^*$ forms a \textit{closed} extension of $\linA$ in this case,
the operator $\linA$ itself is closable
and the closure $\overline{\linA}$ is symmetric,
too.
\end{itemize}

\subparagraph{} %- A+i and A-i are bounded below
A helpful property of symmetric operators is given by the fact
that the closed and symmetric operators $\linA \pm \zeta$ are bounded below
for $\im \zeta \neq 0$,
hence they are injective and own closed ranges.
The codimensions $d_{\pm}(\zeta)$ of the subspaces $\range{\linA \pm \zeta}$
are constant in the upper and lower half-plane of $\C$,
respectively,
which is a consequence of the stability theorem for semi-Fredholm operators.
The values $d_{\pm} := d_{\pm}(i) \xeof \N_0 \cup \{\infty\}$
are called the \textit{defect indices} of the operator $\linA$.

\item[4.] %-- self-adjoint operators
%------------------------------------------------------------------------------%
If a linear operator $\linA$ acting in $\hilH$ is symmetric,
then every other linear operator $\linS$ having the property
$\linA \prec \linS$ is called a \textit{symmetric extension} of $\linA$.
Notice 
that $\linA \prec \linS$ always implies $\linS^* \prec \linA^*$,
which is also true for non-symmetric operators $\linA$ and $\linS$.

\subparagraph{}
Now,
let $\linS_1,\ldots,\linS_k$ be a finite sequence
of symmetric extensions of $\linA$,
where $\linS_{i+1}$ extends $\linS_i$ for $i \eq 1,\ldots,k-1$.
Since $\linS_k \prec \linS_k^*$ and
$\linS_{i+1}^* \prec \linS_i^*$ for $i \eq 1,\ldots,k-1$,
we obtain by concatenating all these extensions the chain
\[
\linA \, \prec \, \linS_1
\, \prec \,
\linS_2
\, \prec \,
\ldots
\, \prec \,
\linS_k
\, \prec \,
\linS_k^*
\, \prec \,
\ldots
\, \prec \,
\linS_2^*
\, \prec \,
\linS_1^*
\, \prec \,
\linA^*.
\]
Thus one may hope to find a \textit{maximal} symmetric extension
$\linS_{sa}$ of $\linA$ such that
\[
\linA \, \prec \, \linS_{sa} \= \linS_{sa}^* \, \prec \, \linA^*
\]
holds,
and in this case
the operator $\linS_{sa}$ cannot be extended to a larger symmetric operator.
This expectation suggests to introduce the notion of \textit{self-adjointness}:
a symmetric linear operator $\linA$ is called \textit{self-adjoint},
if $\linA^* \prec \linA$,
or equivalently,
if $\linA \eq \linA^*$.

\subparagraph{} %- Fredholm property of self-adjoint linear operators
Let us first note
that if a self-adjoint linear operator $\linA$ is semi-Fredholm,
then $\linA$ is actually Fredholm with $\chi(\linA) \eq 0$,
which follows from the equality $\chi(\linA) \eq - \chi(\linA^*)$.
Moreover,
the various definitions of the essential spectrum of $\linA$ given in section 2
are equivalent when assuming that $\linA$ is self-adjoint.
We then have 
\[
\sigma_{ess,1}(\linA)  =
\sigma_{ess,2}(\linA)  =
\sigma_{ess,3}(\linA)  =
\sigma_{ess,4}(\linA)  =
\sigma_{ess,5}(\linA),
\]
which gives rise to define $\sigma_{ess}(\linA) := \sigma_{ess,k}$
for any $1 \leq k \leq 5$.

\subparagraph{} %- Weyl sequence
Recall that the discrete spectrum
$\sigma_d(\linA) := \sigma(\linA) \setminus \sigma_{ess}(\linA)$
is nothing but the set of isolated eigenvalues of $\linA$
having finite multiplicity,
that is,
$\kernel{\linA - \lambda}$ is finite-dimensional
for every $\lambda \xeof \sigma_d(\linA)$.
On the other hand,
any value $\lambda$ belongs to $\sigma_{ess}(\linA)$,
if there is a sequence $\{x_n\}_{n \xeof \N}$ of orthonormal elements in $\hilH$
such that $\norm{\linA x_n - \lambda x_n} \xto 0$ for $n \xto \infty$.
%which is denoted a \textsc{Weyl} \textit{sequence} {\wrt} $\lambda$.

\paragraph{} %- range condition
To ensure self-adjointess of a symmetric linear operator,
the \textit{range condition} serves as the main tool to verify this property:

\begin{theorem1}\label{RANGE_CONDITION}
A densely defined, closed and symmetric linear operator $\linA$
acting in $\hilH$ is self-adjoint iff there is $\zeta \xeof \C$ such that
\[
\range{\linA + \zeta} \= \range{\linA + \overline{\zeta}} \= \hilH,
\]
or equivalently,
if the defect indices of $\linA$ fulfill $d_+(i) \eq d_-(i) \eq 0$.
\end{theorem1}

\paragraph{} %- spectral properties of self-adjoint operators
The \textit{spectral analysis} of self-adjoint operators has applications
especially in quantum mechanics and provides a lot of useful results.
First of all,
we may define self-adjointness of a densely defined and closed linear operator
$\linA$ by means of its spectrum:
the operator $\linA$ is self-adjoint iff $\sigma(\linA) \xsubset \R$,
which is a consequence of the range condition \ref{RANGE_CONDITION}.

\subparagraph{} %- resolvent estimate for self-adjoint operators
A further result is given by the \textit{resolvent estimate}:
if $\linA$ is self-adjoint,
then we have for $\im \zeta \neq 0$ the resolvent estimation
\[
\norm{R(\linA,\zeta)} \leq \abs{\im \zeta} ^{-1},
\]
which follows from the inequality 
$\abs{\im \zeta} \leq \mbox{dist}(\zeta,\sigma(\linA))$
and the general estimate \eqref{RESOLVENT_ESTIMATE_I}.

\subparagraph{} %- residual spectrum of a self-adjoint operator is empty
Moreover,
if $\linA$ is self-adjoint and $\lambda \xeof \sigma(\linA)$,
then $\lambda$ is either an eigenvalue of $\linA$
or $\range{\linA-\lambda}$ is dense in $\hilH$,
but not the whole of $\hilH$,
hence in the first case the value $\lambda$
belongs to the {point spectrum} $\sigma_p(\linA)$ of $\linA$,
whereas the second case means that $\lambda$ belongs
to the {continuous spectrum} $\sigma_c(\linA)$.
Hence the {residual spectrum} $\sigma_r(\linA)$ proves to be empty.

\paragraph{} %- spectral decomposition of self-adjoint operators
The most important theorem {\wrt} self-adjoint linear operators
is the celebrated \textit{spectral theorem}.

\subparagraph{} %- spectral resolutions
Preliminary,
by a \textit{spectral resolution} is meant
a one-parameter family $\{\linP_{\lambda}\}_{\lambda \xeof \R}$
of linear operators $\linP_{\lambda} \: \hilH \xto \hilH$
such that

\begin{itemize}[leftmargin=12mm]
\item[(R-1)]
$\linP_{\lambda}$ is a projection for each $\lambda \xeof \R$;
\item[(R-2)]
for each $x \xeof \hilH$
the mapping $\lambda \mapsto \linP_{\lambda} x$ is continuous from the right;
\item[(R-3)]
$\lambda {\,<\,} \mu$ implies $\linP_{\lambda} \leq \linP_{\mu}$,
so $\scalar{\linP_{\lambda} x}{x} \leq \scalar{\linP_{\mu} x}{x}$
for $x \xeof \hilH$;
\item[(R-4)]
$\lim_{\lambda \xto -\infty} \linP_{\lambda} \, x \eq 0$
and
$\lim_{\lambda \xto  \infty} \linP_{\lambda} \, x \eq x$
for all $x \xeof \hilH$.
\end{itemize}

\subparagraph{}
Here each projection $\linP_{\lambda}$ is related
to a closed linear subspace $\setA_{\lambda} \xsubset \hilH$.
Moreover,
the difference $\linP_{\mu} - \linP_{\lambda}$
for $\mu \geq \lambda$ also forms a projection 
being related to the closed subspace $\setA_{\mu} \ominus \setA_{\lambda}$.

\paragraph{} %- spectral operators
A spectral resolution can be used to define an operator $\linA$
by taking the spectral integral
\[
\linA x \= \int_{\,\R} \lambda \, d\linP_{\lambda}x,
\]
where the domain of definition $\domain{\linA}$
consists of those elements $x \xeof \hilH$,
for which this improper Riemann integral converges.
The mapping $\linA$ is said to be the \textit{spectral operator}
associated with the resolution $\{\linP_{\lambda}\}_{\lambda \xeof \R}$
and proves to be linear and even self-adjoint.

\subparagraph{} %- spectral theorem for self-adjoint operators
Now the \textit{spectral theorem} for self-adjoint operators
claims the inverse statement,
namely that each self-adjoint operator $\linA$
owns a spectral resolution $\{\linP_{\lambda}\}_{\lambda \xeof \R}$
such that
\[
\domain{\linA} \=
\{\, x \xeof \hilH:
    \int_{\,\R}
\lambda^2 \, d \scalar{\linP_{\lambda} \,x}{x} < \infty\,\}
\]
and
\[
\scalar{\linA \, x}{y} \=
\int_{\,\R} \lambda \, d\scalar{\linP_{\lambda} \,x}{y}
\quad \mbox{for } x \xeof \domain{\linA} \mbox{ and } y \xeof \hilH.
\]
The latter formula is often written down in the symbolic form
\[
\linA x \= \int_{\,\R} \lambda \, d\linP_{\lambda}x.
\]
The real number $\lambda$ belongs to $\rho(\linA)$
iff
$\setA_{\lambda + \epsilon} \ominus \setA_{\lambda - \epsilon} \eq \{\0\}$
for some $\epsilon {\,>\,} 0$,
or in other words,
if the mapping $\zeta \mapsto \linP_{\zeta}$ is constant
in a neighbourhood of $\lambda$.

\subparagraph{} %- the functional calculus for self-adjoint operators
The spectral theorem is the base
for the so-called \textit{functional calculus} for self-adjoint operators:
if a given function $f \: \R \xto \C$ is sufficiently regular
(for example,
if $f$ is piecewise continuous and the points of discontinuity
do not have any accumulation point in the real line),
then it is possible to define the operator $f(\linA)$
by means of the spectral resolution
$\{\linP_{\lambda}\}_{\lambda \xeof \R}$ of $\linA$:
\[
f(\linA) x \= \int_{\,\R} f(\lambda) \, d\linP_{\lambda}x,
\]
where $\domain{f(\linA)}$,
the domain of definiton of the operator $f(\linA)$,
is defined by
\[
\domain{f(\linA)} \eqdef
\{x \xeof \hilH: \int_{\R} \abs{f(\lambda)} ^2 \, d\linP_{\lambda}x < \infty \}.
\]
Notice that the operator $f(\linA)$ is self-adjoint iff $f$ is real-valued.

\subparagraph{}
We are especially interested in the case,
where the self-adjoint operator $\linA$ is positive and
$f(\lambda) := 0$ for $t<0$ and $e^{\lambda t}$ for $t \geq 0$
(see also section 5,
in which the family $\{e^{-t\linA}\}_{t \geq 0}$ of bounded linear operators
is shown to constitute an analytic semigroup).
But also other operators $f(\linA)$ are meaningful in applications,
which for example are given by the functions
$x \mapsto \sqrt{x}$,
$x \mapsto \abs{x} $,
$x \mapsto \cos{x}$,
$x \mapsto \sin{x}$
and are denoted as
$\sqrt{\linA}$,
$\abs{\linA} $,
$\cos{\linA}$,
$\sin{\linA}$.

\paragraph{} %- perturbation theory of self-adjoint linear operators
We now turn to perturbation theory of self-adjoint operators.
Let $\linA + \linB$ be a linear operator in Hilbert space $\hilH$,
where $\linA$ is assumed to be self-adjoint and $\linB$ to be symmetric
with domain of definition $\domain{\linB}$
being a superset of $\domain{\linA}$,
so $\linA + \linB$ is well defined
with domain of definition $\domain{\linA + \linB} \eq \domain{\linA}$.
We want to show
under which conditions on the operator $\linB$
the sum $\linA + \linB$ proves to be self-adjoint.

\subparagraph{} %- A-boundedness
Let $\linA \: \domain{\linA} \xto \hilH$ be self-adjoint.
Then each other symmetric operator $\linB \: \domain{\linB} \xto \hilH$
with $\domain{\linA} \xsubset \domain{\linB}$
is called $\linA$\textit{-bounded},
if there are constants $\theta, c \geq 0$ such that
\begin{equation}\label{A_BOUNDEDNESS}
\norm{\linB \,x} \, \leq \, \theta \norm{\linA \,x} \+ c \norm{x}
\qfa x \xeof \domain{\linA}.
\end{equation}

\subparagraph{} %- A-bound
The infimum of all numbers $\theta$,
for which there is $c \geq 0$
such that the estimate \eqref{A_BOUNDEDNESS} holds,
is called the $\linA$\textit{-bound} of $\linB$.
Also notice that
if $\linB$ is bounded,
then it is $\linA$-bounded with $\linA$-bound $\theta <1$ for any $\linA$.
In addition,
the validness of the estimate
\[
\norm{\linB \,x}^2 \, \leq \,
\theta^2 \norm{\linA \,x}^2 \+ c^2 \norm{x}^2
\qfa x \xeof \domain{\linA}.
\]
is a sufficient condition for \eqref{A_BOUNDEDNESS} to be true.

\paragraph{} %- Kato-Rellich theorem
The \textsc{Kato-Rellich} \textit{theorem} provides a criterion
on self-adjoint\-ness of the operator $\linA + \linB$,
which has been independently found by \textsc{T.\,Kato}
and \textsc{F.\,Rellich}:

\begin{theorem1}\label{KATO-RELLICH_THEOREM}
Let $\linA$ be self-adjoint and 
the linear operator $\linB$ be symmetric,
with $\domain{\linB} \supset \domain{\linA}$ 
and $\linA$-bounded with $\theta {\,<\,} 1$.
Then $\linA + \linB$ with domain of definition 
$\domain{\linA + \linB} \eq \domain{\linA}$ is self-adjoint.
\end{theorem1}

\subparagraph{} %- Kato-Rellich theorem for essentially self-adjoint operators
A similar result is obtained for the problem,
whether \textit{essentially self-adjointness} of a linear operator $\linA$
is preserved by a symmetric and $\linA$-bounded perturbation $\linB$:

\begin{theorem1}\label{KATO-RELLICH_THEOREM_II}
Let $\linA \: \domain{\linA} \xto \hilH$ be essentially self-adjoint
and assume that $\linB$ is symmetric and $\linA$-bounded,
so $\domain{\linA} \xsubset \domain{\linB}$ and \eqref{A_BOUNDEDNESS} holds
with bound $\theta < 1$.
Then the operator $\linA + \linB$
is essentially self-adjoint on $\domain{\linA}$,
where
$
\overline{\linA + \linB} \= \overline{\linA} + \overline{\linB}.
$
\end{theorem1}

\item[5.] %-- operators with pure point spectrum
%------------------------------------------------------------------------------%
Let us again remember
that the \textit{discrete spectrum} $\sigma_d(\linA)$
of a linear operator $\linA$ acting in a Banach or Hilbert space
consists of normal eigenvalues of $\linA$,
that is,
such values are isolated in $\sigma(\linA)$ and their eigen\-spaces
own finite multiplicity.
The discrete spectrum is at the same time
the relative complement of the Weyl spectrum $\sigma_{ess,5}$
in $\sigma(\linA)$.
We say that $\linA$ has \textit{pure point spectrum},
if $\sigma(\linA) \eq \sigma_d(\linA)$ and 
$\sigma_{ess,5}(\linA) \eq \emptyset$,
respectively.

\paragraph{} %- Riesz-Schauder theory for compact linear operators
The theory of linear operators having pure point spectrum
mainly relies on two fundamental theorems:

\begin{itemize}[leftmargin=9mm]
\item[(a)]
the \textsc{Riesz-Schauder} \textit{spectral theorem},
claiming that the spectrum $\sigma(\compA)$
of a compact linear operator $\compA$ consists of eigen\-values
$\lambda \neq 0$ having finite-dimensional eigenspaces
and possibly in addition of the value $\lambda \eq 0$; 

\item[(b)]
the \textit{spectral mapping theorem for the discrete spectrum},
which quotes that
the discrete spectrum of any resolvent $R(\linA,\zeta)$ depends on that
of $\linA$ according to:
\[
\sigma_d(R(\linA,\zeta)) \setminus \{0\} \=
\{(\lambda-\zeta)^{-1}: \lambda \xeof \sigma_d(\linA)\}.
\]
\end{itemize}

Both of them together imply that $\linA$ has pure point spectrum
if and only if at least one of the resolvents are compact,
hence $\linA$ is also said to \textit{have compact resolvents} in this case.
Notice that if one resolvent is compact,
then all other resolvents are compact,
too,
which is a consequence of the resolvent equation \eqref{RESOLVENT_EQUATION}.
Moreover,
a closed linear operator $\linA$ has compact resolvents
\iff the embedding of the Banach space $(\domain{\linA},\norm{\,\cdot\,}_A)$
into the space $\banX$ is compact.
%where $\norm{x}_A := \norm{x} + \norm{\linA x}$ is the \textit{operator norm}
%on the domain of definition $\domain{\linA}$.

\paragraph{} %- order of the eigenvalues
The eigenvalues of an operator $\linA$ with discrete spectrum
can be ordered by their moduli,
that is,
there is a sequence $\{ \lambda_n \}_{n \xeof \N}$ of eigen\-values
such that
\[
\abs{\lambda_1} \,\leq\, \abs{\lambda_2} \,\leq\, \abs{\lambda_3} \,\leq\,
\ldots \, .
\]

\paragraph{} %- the symmetric case
But more can be said in the case,
where $\linA$ is symmetric.
In fact,
all values $\lambda_n$ are real-valued,
hence $\linA$ proves to be self-adjoint.
If in addition $\linA$ is half-bounded,
we may improve the order of the sequence of eigenvalues to
\[
\lambda_1 \,\leq\, \lambda_2 \,\leq\, \lambda_3 \,\leq\, \ldots \, .
\]

\subparagraph{} %- discrete spectral theorem
Finally,
the collection of eigen\-vectors $\{x_n\}_{n \xeof \N}$ of $\linA$
forms an orthonormal base in case of $\hilH$ is separable.
This follows from the spectral mapping theorem for the discrete spectrum
and the celebrated \textsc{Hilbert-Schmidt} \textit{spectral theorem}
for compact-symmetric operators,
also called the \textit{discrete spectral theorem}:

\begin{theorem1}\label{DISCRETE_SPECTRAL_THEOREM}
If the linear operator $\compA$
acting in a separable Hibert space $\hilH$
is compact and symmetric,
then the collection of eigen\-vectors of $\compA$
forms an orthonormal base for $\hilH$.
\end{theorem1}

\subparagraph{} %- spectral decomposition of operators with pure point spectrum
We may conclude from this theorem that
a symmetric operator $\linA$ having pure point spectrum
owns the \textit{spectral decomposition}
\[
\linA \,x \= \sum_{n=1}^{\infty} \lambda_n \scalar{x}{x_n} \, x_n
\]
for all $x$ in the domain of definition $\domain{\linA}$ defined by
\[
\domain{\linA} \=
\{x \xeof \hilH:
\sum_{n \eq 1}^{\infty} \, \lambda_n^2 \, \abs{\scalar{x}{x_n}} ^2 < \infty \}.
\]
\end{itemize}

\paragraph{} %- equivalence of linear equations
A great portion of {Fredholm theory} deals
with the study of linear operators of type $\1+\compA$, 
where $\compA \: \banX \xto \banY$ is compact.
Such operators are indeed Fredholm with index $0$,
to which \textsc{Fredholm}'s alternative theorem can be applied.
We also emphasize the usage of Fredholm theory by the following considerations,
which are helpful in solving the \textit{Dirichlet problem}
for elliptic operators (see section 9).
Consider again the linear equation $\linA x \eq y$,
where $\banX$ is a Banach space and
$\linA \: \domain{\linA} \xsubset \banX \xto \banX$
is assumed to be a linear operator with pure point spectrum.
Then for $\zeta \xeof \rho(\linA)$ we surely have
\[
\linA \,x \= (\linA -\zeta)\,x + \zeta \,x \= y,
\]
and applying the mapping $(\linA -\zeta)^{-1}$ from the left yields
\[
x \+ \zeta (\linA - \zeta)^{-1}\,x
\=
(\linA -\zeta)^{-1}\,y \, =: \, \tilde{y}
\]
or equivalently $x + \tilde{\compA} \,x \eq \tilde{y}$,
where $\tilde{\compA} := \zeta (\linA-\zeta)^{-1}$ is compact.

\subparagraph{}
The equations $\linA \,x \eq y$ and $(\1 + \tilde{\compA})\,x \eq \tilde{y}$
are called \textit{equivalent to each other} in the sense
that $x$ is a solution of the first one
\iff
it is a solution of the second one.
By section 2,
the operator $\1 + \tilde{\compA}$ is Fredholm with vanishing index,
hence theorem \ref{FREDHOLM_ALTERNATIVE_THEOREM} is applicable
and provides assertions also for the solution of the equation $\linA x \eq y$.

\paragraph{}
Let us summarize these considerations to the following

\begin{theorem1}\label{FREDHOLM_ALTERNATIVE_THEOREM_II}
Let $\linA \: \domain{\linA} \xsubset \banX \xto \banX$ be a linear
operator with pure point spectrum.
Then the equation $\linA x \eq y$ is equivalent to an equation
of the form $(\1+\compA)x \eq y$,
where $\compA$ is a compact linear operator,
and \textsc{Fredholm}'s alternative theorem may be applied
to obtain assertions on the solutions of $\linA x \eq y$.
\end{theorem1}

%------------------------------------------------------------------------------%
\section{Generator Theorems for Closed Operators}                              %
%------------------------------------------------------------------------------%
This section serves as an outline on methods to construct solutions
of first and second-order homogeneous Cauchy problems.
The most important ones are the \textit{semigroup method} and
the \textit{functional calculus method},
where the latter one is usually applied to initial value problems
governed by positive and self-adjoint linear operators.

\paragraph{} %-- semigroups of bounded linear operators
%------------------------------------------------------------------------------%
Let us start with the semigroup method.
The notion of \textit{semigroup} has been introduced
by \textsc{J.\,A.\,S\'{e}guier} in 1904.
Later some applications to partial differential equations were detected,
but the first systematic and comprehensive treatement of semigroup theory
was published in the pioneering monograph \cite{HP}
by \textsc{E.\,Hille} and \textsc{R.\,Phillips}.
Today semigroup theory provides a powerful method
in solving evolution problems
with a multiplicity of applications to partial differential equations
governed by linear-elliptic operators.

\subparagraph{} %- infinitesimal generators and well-posedness
As it has been elaborated in section 3,
the homogeneous Cauchy problem of first order
\eqref{HOMOGENEOUS_CAUCHY_PROBLEM_I} is well-posed iff
the linear operator ${-}\linA$ is the infinitesimal generator
of a strongly continuous semigroup.
We already pointed out that
in this case
the operator $\linA$ must be \textit{densely defined} and \textit{closed}.

\subparagraph{} %- generator theorems
Reversely,
the question occurs
under which conditions a densely defined and closed linear operator ${-}\linA$
constitutes the infinitesimal generator of a $C_0$-semigroup.
This lies at the heart of semigroup theory and
can be answered by a couple of \textit{generator theorems}.

\paragraph{} %- classification of strongly continuous semigroups
Preliminary,
the totality of strongly continuous semigroups may be classified as follows:
for each C$_0$-semigroup $\{\linS(t)\}_{t \geq 0}$
there are constants $M \geq 1$ and $\omega \xeof \R$
such that for all $t \geq 0$ we have
\[
\norm{\linS(t)} \, \leq \, M e^{\omega t}.
\]

\begin{comment}
In fact,
let $M \geq 1$ such that $\norm{S(t)} \leq M$ for all $s \xeof [0,1]$.
For $t {\,>\,} 0$ there is $s \xeof [0,1)$ and $n \xeof N$
such that $t \eq n+s$.
Then $\norm{\linS(n)}\leq M \norm{S(n-1)}$ for $n {\,>\,} 1$ and
\[
\norm{\linS(t)} \eq
\norm{\linS(n+s)} \eq \norm{\linS(n) \circ \linS(s)} \leq
\norm{\linS(n)} {\cdot} \norm{\linS(s)} \leq M^{n+1} \eq M e^{n \cdot \ln{M}}.
\]
\end{comment}

\paragraph{} %-- generator theorems for strongly continuous semigroups
%------------------------------------------------------------------------------%
The \textsc{Feller-Miyadera-Hille-Yosida} \textit{theorem}
is outstanding in semigroup theory
and characterizes properties of the operator ${-}\linA$
to be the infinitesimal generator of a C$_0$-semigroup: 

\begin{theorem1}\label{MAIN_THEOREM_OF_SEMIGROUP_THEORY}
Let $\linA$ be a densely defined and closed linear operator
acting in an arbitrary Banach space $\banX$.
Then the following assertions are equivalent to each other:

\begin{itemize}[leftmargin=9mm]
\item[(a)]
The operator ${-}\linA$ generates a strongly continuous semigroup
of type $(\omega,M)$.
\item[(b)]
The resolvent set of $\linA$ fulfills $(-\infty,\omega) \xsubset \rho(\linA)$
and the resolvent function $R(\linA,\zeta)$ satisfies
\begin{equation}\label{HILLE-YOSIDA_CONDITION}
\norm{R(\linA,\zeta)^n} \, \leq \,
\frac{M}{(\omega - \zeta)^n} \quad \quad (n \in \N, \zeta < \omega).
\end{equation}
\end{itemize}
\end{theorem1}

\subparagraph{}
However,
since the number of conditions on the operator $\linA$ is infinite,
this theorem has limited worth {\wrt} applications
and even in the simplest cases.

\subparagraph{} %- Hille-Yosida theorem
But the original theorem for generators of \textit{contraction semigroups}
($\omega \eq 0$ and $M \eq 1$)
due to \textsc{Hille} and \textsc{Yosida} quotes
that in this case only the first inequality
($n \eq 1$) in \eqref{HILLE-YOSIDA_CONDITION} has to be verified:

\begin{theorem1}\label{HILLE-YOSIDA_THEOREM}
The operator ${-}\linA$ acting in a Banach space $\banX$
is the infinitesimal generator of a C$_0$-semigroup of \textit{contractions}
iff

\begin{itemize}[leftmargin=9mm]
\item[(a)]
$\linA$ is densely defined and closed,
\item[(b)]
the interval $(-\infty,0)$ belongs to the resolvent set $\rho(\linA)$,
\item[(c)]
$\norm{\zeta \, R(\linA,\zeta) } \leq 1$ for all $\zeta {\,<\,} 0$.
\end{itemize}
\end{theorem1}

\subparagraph{} %- Laplace transformation
If ${-}\linA$ owns the properties
mentioned in theorem \ref{HILLE-YOSIDA_THEOREM},
then the semigroup $\{\linS(t)\}_{t \geq 0}$ generated by ${-}\linA$
may be constructed by
\[
\linS(t) \, x \eqdef
\lim_{\,n \xto \infty} (\1 + \frac{t}{n} \linA)^{-n} \,x 
\fe t > 0 \and x \xeof \banX.
\]
In addition,
the resolvents of the operator $\linA$ can be derived
from its generated semigroup according to the \textit{Laplace transformation}
\[
R(\linA,\zeta)x \= -\int_0^{\infty} e^{-\zeta \, t} \, S(t) x \, dx
\quad (\zeta \xeof \rho(\linA), x \xeof \banX).
\]

\paragraph{} %-- Lumers-Phillips theorem
In the following we want to characterize the properties (b) and (c)
of theorem \ref{HILLE-YOSIDA_THEOREM} by equivalent conditions in the case,
where $\linA$ is acting in a Hilbert space.

\subparagraph{}
By means of the concept of \textit{m-accretivity},
the \textsc{Lumer-Phillips} \textit{theorem}
classifies the linear operators
being infinitesimal generators of contractive C$_0$-semigroups,
but without using estimates on their resolvents:

\begin{theorem1}\label{LUMER-PHILLIPS_THEOREM}
Let $\linA$ be a linear operator in Hilbert space.
Then ${-}\linA$ generates a contractive C$_0$-semigroup
iff $\linA$ is m-accretive.
\end{theorem1}

\paragraph{} %-- semigroup methods for second-order evolution equations
%------------------------------------------------------------------------------%
In the following we want to apply the semigroup method
to \textit{second order} problems.
Similar to the first-order Cauchy problem,
a function $x$ is called
a \textit{classical solution} of the second-order homogeneous Cauchy problem
\eqref{HOMOGENEOUS_CAUCHY_PROBLEM_II},
if $x \xeof C([0,\T),\domain{\linA}) \, \cap \, C^2((0,\T),\hilH)$
and satisfies in each point $t \xeof (0,\T)$ the differential equation
and the both initial conditions.

\subparagraph{}
We suppose that $\linA$ is \textit{strictly positive} and \textit{self-adjoint},
hence the operator $\sqrt{\linA}$ is well-defined
by the functional calculus for self-adjoint operators.
Letting $U := (x,x')$,
the problem \eqref{HOMOGENEOUS_CAUCHY_PROBLEM_II} can be written
as the first order problem
\[
\begin{cases}[1.2]
\quad U'(t) \= \linB \, U(t) \\
\quad U(0)  \= (x_0,x_1),\\
\end{cases}
\]
where the linear operator $\linB$ is defined by
\[
\linB \eqdef \begin{pmatrix} 0 & \1 \\ \sqrt{\linA} & 0 \end{pmatrix}.
\]

\subparagraph{}
The main idea in proving well-posedness of \eqref{HOMOGENEOUS_CAUCHY_PROBLEM_II}
is to show that the operator $\linB$
constitutes the infinitesimal generator of a semigroup of contractions,
when defined on a suitable chosen domain of definition $\domain{\linB}$.
Let us define $\domain{\linB}$ according to
\[
\domain{\linB} \eqdef
\domain{\linA} \xtimes \domain{\sqrt{\linA} }
\, \xsubset \,
\domain{\sqrt{\linA} }_{\norm{\cdot}_{\linA}} \xtimes \hilH,
\]
where $\domain{\sqrt{\linA} }_{\norm{\cdot}_{\linA}}$
is the Hilbert space $\domain{\sqrt{\linA} }$
furnished with the graph norm.
In fact,
it can be shown that
the operator $\linB$ generates a C$_0$-semigroup,
if $\sqrt{\linA}$ does so and $0 \xeof \rho(\sqrt{\linA})$,
see also theorem 7.2 in \cite{Go} (see page\,110).

\paragraph{} %-- analytic semigroups
%------------------------------------------------------------------------------%
Semigroups of bounded linear operators are usually parametrized
by the half line $[0,\infty)$.
But it is also customary to study possible extensions to a symmetric sector
\[
\Sigma_{\,\theta} \eqdef
\{z=re^{i\alpha} \xeof \C: r > 0\,, \abs{\alpha} < \theta\}
\]
of angle $2* \theta$ for some $\theta \xeof (0,\frac{\pi}{2}]$.
Indeed,
a semigroup $\{\linS(t)\}_{t \geq 0}$ is called \textit{analytic},
if there is an angle $\theta \xeof (0,\frac{\pi}{2})$
and an analytic extension $\linS' \: \Sigma_{\,\theta} \xto \BO{X}$ of $\linS$,
which is \textit{locally bounded},
that is,
\[
\sup \,
\{ \, \norm{\linS'(z)}: z \xeof \Sigma_{\,\theta}, \, \norm{z} \leq 1 \}
\, < \, \infty,
\]
and in which the semigroup property $\linS'(t+s) \eq \linS'(t) \circ \linS'(s)$
holds for all $s,t \xeof \Sigma_{\,\theta}$.

\paragraph{} %- sectorial operators generate analytic semigroups
It was already shown in section 4
that the functional calculus for sectorial operators enables
to define the semigroup $\{e^{-tA}\}_{t \geq 0}$ of linear operators
by a contour integral around the spectrum $\sigma(\linA)$.

\paragraph{}
But more is true in this case:

\begin{theorem1}\label{SEMIGROUP_CREATED_BY_SECTORIAL_OPERATOR_IS_ANALYTIC}
If $\linA$ is a sectorial operator,
then the semigroup $\{e^{-tA}\}_{t \geq 0}$ is {analytic}.
\end{theorem1}

\subparagraph{} %- sectorially contractive analytic semigroups
Reversely,
if a strongly continuous semigroup $\{\linS(t)\}_{t \geq 0}$
is analytic and the estimate $\norm{\linS(\zeta)} \leq 1$ holds
for all $\zeta \xeof \Sigma_{\,\theta}$,
then it is called \textit{sectorially contractive analytic} of angle $\theta$,
and its infinitesimal generator proves to be sectorial.

\paragraph{} %-- the generator theorem for analytic semigroups
%------------------------------------------------------------------------------%
In summary,
the generation theorem for the class of analytic semigroups reads as follows:

\begin{theorem1}\label{CRITERION_FOR_SEMIGROUP_TO_BE_SECTORIALLY_CONTRACTIVE}
Let $\linA$ be a densely defined linear operator
acting in a Banach space $\banX$
and let $\theta \xeof (0,\frac{\pi}{2}]$.
Then the following assertions are equivalent:

\begin{itemize}[leftmargin=9mm]
\item[(a)]
$\linA$ is sectorial of angle $\theta$;
\item[(b)]
${-}\linA$ generates a sectorially contractive analytic $C_0$-semigroup.
\end{itemize}
\end{theorem1}

\subparagraph{} %- properties of analytic semigroups
The advantage of working with analytic semigroups
instead of strongly continuous ones is given by the following facts:

\begin{itemize}[leftmargin=4mm]
\item[1.]
If $\linA$ generates an analytic semigroup,
then $\linA x \xeof \domain{\linA}$ for each $x \xeof \domain{\linA}$.

\item[2.]
The solution $x$ of the first-order homogeneous Cauchy problem is indefinitely
differentiable in $(0,\T)$.
By this reason,
the differential equation $x' + \linA x \eq y$
is often called \textit{parabolic},
when governed by a sectorial operator $\linA$.
\end{itemize}

\paragraph{} %-- generation of solution operators by spectral resolution
%------------------------------------------------------------------------------%
In the second part of this section we want to review semigroups
generated by \textit{positive self-adjoint} or \textit{skew-adjoint} operators.

\subparagraph{}
Let $\hilH$ be a Hilbert space and $\linA$ be self-adjoint and positive,
that is,
$\linA^* \eq \linA$ and $\Theta(\linA) \xsubset [0,\infty)$.
Such an operator is doubtless m-accretive and even sectorial,
hence ${-}\linA$ forms the generator
of a strongly continuous and analytic semigroup.
But due to the spectral theorem and the functional calculus
for self-adjoint operators,
the semigroup can defined in a second way.

\subparagraph{} %- semigroups defined by spectral resolutions
Let $\{ P_{\lambda} \}_{\lambda \xeof \R}$
be the spectral resolution of $\linA$.
We first notice that $P_{\lambda} \eq \0$ for $\lambda {\,<\,} 0$,
since $\linA$ is positive.
Then the functional calculus applied to the exponential function
$t \mapsto e^{-t}$ gives rise to the family of spectral operators
\[
\linL_t(x) \eqdef e^{-t \linA} x
\=
\int_{\,0}^{\,\infty} e^{-\lambda t} \, dP_{\lambda} \, x
\quad \quad (t \geq 0, \, x \xeof \hilH).
\]
We apparently have $\linL_0 \eq \1$,
and $\norm{x(t)} \leq \norm{x(0)}$ for $t>0$ implies
$\norm{\linL_t} \leq 1$ for all $t {\,>\,} 0$.
Hence the family $\{ \linL_t\}_{t \geq 0}$ forms a contractive semigroup
of bounded linear operators in $\hilH$,
which especially solves the homogeneous Cauchy problem
\eqref{HOMOGENEOUS_CAUCHY_PROBLEM_I}.

\subparagraph{} %- functional calculus approach
                %  for the second-order Cauchy problem
The functional calculus for self-adjoint linear operators
may also be applied to the homogeneous Cauchy problem of second order.
Preliminary,
each solution $x$ of this problem satisfies
the \textit{energy conservation law}
\[
\norm{x'(t)}^2 \+ \snorm{\sqrt{\linA} \, x(t)}^2
\=
\norm{x'(0)}^2 \+ \snorm{\sqrt{\linA} \, x(0)}^2
\quad (0 < t < \T),
\]
which implies that the problem owns mostly one solution.
Let $\{\linP_{\lambda}\}_{\lambda \xeof \R}$ be
the spectral resolution of $\linA$.
Then we define a function $x \: [0,\T) \xto \hilH$
by means of the spectral operator $\linL_t$,
induced by the functions $t \mapsto \cos(t)$ and $t \mapsto \sin(t)$.
In fact,
the operator $x \mapsto \linL_t x$
is built as the sum of two single spectral operators
for each $0 \leq t {\,<\,} \T$,
namely
\[
x(t)
:=
\linL_t(x_0,x_1)
:=
\int_{\,0}^{\,\infty} \cos{\sqrt{\lambda}t} \, dP_{\lambda} \, x_0
\+
\int_{\,0}^{\,\infty} \frac{\sin{\sqrt{\lambda}t}}{\sqrt{\lambda}}
\, dP_{\lambda} \, x_1.
\]

\subparagraph{} %- well-posedness of the second-order Cauchy problem
Having defined the family of linear operators $\{\linL_t(x_0,x_1)\}$ this way,
we obtain the following existence and uniqueness theorem
for the second order Cauchy problem:
the mapping $t \mapsto \linL_t(x_0,x_1)$ is continuous in $[0,\T)$,
twice continuously differentiable in $(0,\T)$ and
constitutes the unique solution
of problem \eqref{HOMOGENEOUS_CAUCHY_PROBLEM_II}.

\subparagraph{} %-- the generation theorem for unitary groups
%------------------------------------------------------------------------------%
We finally want to consider so-called \textit{unitary groups}.
A semigroup $\{\linU(t)\}_{t \geq 0}$ of bounded linear operators
acting in a Hilbert space $\hilH$ is called \textit{unitary},
if each single operator $U(t)$ is so.
Furthermore,
if one defines $\linU(-t) := \linU(t)$ for $t {\,>\,} 0$,
then a group $\{\linU(t)\}_{t \xeof \R}$ of unitary linear operators is defined,
which fulfills
\[
\linU(s+t) \= \linU(s) \circ \linU(t) \quad \quad (t,s \xeof \R).
\]
\subparagraph{} %- Stone's theorem
\textsc{Stone}'s theorem claims
that the infinitesimal generator of a unitary group is \textit{skew-adjoint}:

\begin{theorem1}\label{STONE_THEOREM_FOR_UNITARY_GROUPS}
Each unitary group $\{\linU(t)\}_{t \xeof \R}$
is infinitesimally generated by a linear operator of type $i \linA$,
where $\linA$ is self-adjoint.
\end{theorem1}

\subparagraph{} %- abstract Schrödinger equations
The functional calculus for self-adjoint linear operators
is also helpful to solve further classes of abstract initial value problems.
By the \textit{abstract {\SCHROEDINGER} equation}
we mean the linear operator-differential equation
\begin{equation}\label{ABSTRACT_SCHROEDINGER_EQUATION}
x'(t) \+ i \linA \, x(t) = 0\,,
\end{equation}
where $\linA$ is a self-adjoint operator in Hilbert space $\hilH$.
Each solution $x$ of \eqref{ABSTRACT_SCHROEDINGER_EQUATION}
satisfies the \textit{energy equation}
\begin{equation}\label{ENERGY_CONSERVATION_LAW}
\norm{x(t)} = \norm{x(0)}
\end{equation}
for $t \xeof \R$,
which immediately follows from
\[
\frac{d}{dt} \, \norm{x(t)}^2
\=
\scalar{-i\linA x(t)}{x(t)} + \scalar{x(t)}{-i \linA x(t)} \= 0.
\]
If an initial value $x(0) \eq x_0$ is given,
then \eqref{ENERGY_CONSERVATION_LAW} implies
that the solution of the initial value problem 
for \eqref{ABSTRACT_SCHROEDINGER_EQUATION} is unique.
Furthermore,
the existence of a solution for the abstract {\SCHROEDINGER} equation
is ensured by the presence of the unitary group
infinitesimally generated by the skew-symmetric operator $i \linA$.

%------------------------------------------------------------------------------%
\section{Coercive Sesquilinear Forms}                                          %
%------------------------------------------------------------------------------%
This section is an exposition on sesquilinear forms
in Hilbert space and their usage in the solution theory of linear problems.
Especially the possibility to derive linear operators from sesquilinear forms, 
which are \textit{coercive} in some sense,
is motivation enough to consider this special class of forms.

\paragraph{} %-- generalities on sesquilinear forms
%------------------------------------------------------------------------------%
We start with some preliminaries.
By a \textit{sesquilinear form} defined in complex linear space $\vecV$
is meant a mapping $\sql \: \vecV \xtimes \vecV \xto \C$,
which is linear in the first and conjugate linear in the second argument,
that is,
for all $x,y,z \xeof \vecV$ and $\lambda,\mu \xeof \C$ the computation rules
\[
\begin{cases}[1.2]
\quad \sql(x+y,z) \= \sql(x,z) \+ \sql(y,z) \\
\quad \sql(x,y+z) \= \sql(x,y) \+ \sql(x,z) \\
\quad \sql(\lambda x,\mu y) \= \lambda \overline{\mu} \, \sql(x,y) \\
\end{cases}
\]
are valid.
If the linear space $\vecV$ is real, 
however,
then $\sql$ is assumed to be linear also in the second argument,
so it becomes into a \textit{bilinear form} in this case.

\subparagraph{} %- representation operators
A sesquilinear form $\sql$ may also be regarded
as a linear mapping $\repA \: \vecV \xto \vecV_a^{\#}$,
where $\vecV_a^{\#}$ is the (algebraic) \textit{conjugate space} of $\vecV$,
that is,
the linear space of conjugate linear functionals $\phi \: \vecV \xto \C$.
More precisely,
given a sesquilinear form $\sql$,
each element $x \xeof \vecV$ gives rise to a conjugate linear functional
$\phi_x \: \vecV \xto \C$,
defined by $\phi_x(y) := \sql(x,y)$ for $y \xeof \vecV$.
Let $\vecV_a^{\#}$ be the linear space
of conjugate linear functionals on $\vecV$.
Then the mapping
\[
\repA: \;
\begin{cases}[1.2]
\quad \repA \: \vecV \xto \vecV_a^{\#},\\
\quad \repA \, x \eqdef \phi_x \quad (x \xeof \vecV)
\end{cases}
\]
is linear and denoted the \textit{representation} or \textit{form operator}
of $\sql$.

\subparagraph{}
Reversely,
for any linear operator $\cal{A} \: \vecV \xto \vecV_a^{\#}$,
there is a sesqui\-linear form $\sql \: \vecV \xtimes \vecV \xto \F$
such that $\sql(x,y) \eq (\cal{A} \, x) y$ for $x,y \xeof \vecV$.
Hence sesquilinear forms acting on a complex linear space $\vecV$
appear in a one-to-one correspondence to linear mappings
acting between the spaces $\vecV$ and $\vecV_a^{\#}$.
This correspondence remains also valid,
if $\vecV$ is a \textit{real} linear space,
$\sql$ is a \textit{bilinear} form
and $\cal{A}$ is the associated mapping
from $\vecV$ to the adjoint space $\vecV_a^{\#}$.

\paragraph{} %-- continuous and coercive forms
%------------------------------------------------------------------------------%
From now on we assume that $\vecV$ is a \textit{normed} linear space
with underlying field $\C$.
Then any sesquilinear form $\sql \: \vecV \xtimes \vecV \xto \C$
is called \textit{bounded},
if there is $\Theta {\,>\,} 0$ such that
\[
\abs{\sql(x,y)} \, \leq \Theta \, \norm{x} \cdot \norm{y}
\quad \quad (x,y \xeof \vecV).
\]
Due to the \textit{uniform boundedness principle},
this property is equivalent to say
that $\sql$ is continuous on the space $\vecV \xtimes \vecV$.

\subparagraph{} %- representation of continuous forms
The prototypical example for a continuous form is given by
the inner product in a complex Hilbert space $\hilH$,
where boundedness is ensured by the Cauchy-Schwarz inequality
$\abs{ \scalar{x}{y} } \leq \norm{x} \cdot \norm{y}$.

\subparagraph{}
For each continuous sesquilinear form $\sql \: \hilH \xtimes \hilH \xto \C$
there is a bounded linear operator $\linB_{\sql} \: \hilH \xto \hilH$
such that
\begin{equation}\label{FORM_REPRESENTATION}
\sql(x,y) \= \scalar{\linB_{\sql} \, x}{y} \qfa x,y \xeof \hilH,
\end{equation}
which is a representation for the form $\sql$.
In fact,
if $x \xeof \hilH$ is fixed,
then the mapping $\sql(x,\cdot) \: \hilH \xto \C$
causes a continuous antilinear functional $\phi_x \xeof \hilH_a^*$,
and by \textsc{Riesz}'s representation theorem
there is $\tilde{x} \xeof \hilH$ such that
$\phi_x(y) \eq \scalar{\tilde{x}}{y}$ for all $y \xeof \hilH$.
This suggests to define $\linB_{\sql} x := \tilde{x}$
for every $x \xeof \hilH$.
If \eqref{FORM_REPRESENTATION} holds for some operator
$\linB_{\sql} \xeof \BO{\hilH}$,
then the form $\sql$ is said to be represented by $\linB_{\sql}$.

\subparagraph{} %- coercive forms
Finally,
any continuous sesquilinear form $\sql \: \hilH \xtimes \hilH \xto \C$
is called \textit{coercive}
if there is $\theta {\,>\,} 0$ such that
\begin{equation}\label{COERCIVITY_PROPERTY}
\re \, \sql(x,x) \, \geq \theta \norm{x}^2
\qfa x \xeof \hilH.
\end{equation}
The estimate \eqref{COERCIVITY_PROPERTY} especially implies the validness of
\[
\norm{\linB_{\sql} x}\norm{x}    \, \geq \,
\abs{\scalar{\linB_{\sql} x}{x}} \, \eq  \,
\abs{\sql(x,x)}                  \, \geq \,
\re \sql(x,x)                    \, \geq \,
\theta \norm{x}^2
\]
for all $x \xeof \hilH$.
Hence the form operator $\linB_{\sql}$ is bounded below
with lower bound $\theta$,
and we conclude by theorem \ref{BOUNDED_BELOW_THEOREM}
that $\linB_{\sql}$ is injective and the range $\range{\linB_{\sql}}$
forms a closed subspace in $\hilH$.

\paragraph{} %-- variational equations and the Lax-Milgram theorem
%------------------------------------------------------------------------------%
Coercivity of a sesquilinear form $\sql$ is of great interest {\wrt}
the solvability of the so-called \textit{variational problem},
which may be formulated as follows:
given a continuous and conjugate linear functional $\phi \xeof \vecV_a^*$,
it is to find $x \xeof \vecV$ such that
\begin{equation}\label{VARIATIONAL_EQUATION}
\sql(x,y) \= \phi(y) \qfa y \xeof \vecV.
\end{equation}

\subparagraph{} %- Lax-Milgram theorem
The \textsc{Lax-Milgram} \textit{theorem} is of outstanding meaning
in applied functional analysis
and ensures well-posedness of the variational problem
for any right-hand side $\phi \xeof \vecV_a^*$:

\begin{theorem1}\label{LAX-MILGRAM_THEOREM}
If $\sql \: \vecV \xtimes \vecV \xto \C$ is bounded and coercive,
then for each $\phi \xeof \vecV_a^*$
the equation \eqref{VARIATIONAL_EQUATION} owns exactly
one solution $x_{\phi} \xeof \vecV$.
In addition,
$\norm{\,x_{\phi}\,} \leq \theta^{-1} \norm{\phi}$ holds.
\end{theorem1}

In other words,
if $\sql$ proves to be bounded and coercive,
then the form operator $\repA \: \vecV \xto \vecV_a^*$ is bjective.
But the {bounded inverse theorem} implies
that the inverse operator $\repA^{-1}$ is bounded,
too,
so the solution $x_{\phi}$ of \eqref{VARIATIONAL_EQUATION}
depends continuously on $\phi$.

\begin{comment}
\subparagraph{}
The proof of theorem \ref{LAX-MILGRAM_THEOREM} is based on the fact
that the representation operator $\linB_{\sql}$ of $\sql$ is bounded below.

\Proof:
Since $\linB_{\sql}$ is bounded below,
it is injective and $\range{\linB_{\sql}}$ is closed.
It is still to show that $\linB_{\sql}$ is surjective.
If $z \xeof \range{\linB_{\sql}}^{\bot}$,
then for $\linB_{\sql} z \xeof \range{\linB_{\sql}} $ we have
\[
0 \=
\scalar{\linB_{\sql} z}{z}
\=
\sql(z,z)
\, \geq \,
\theta \norm{z}^2
\, \geq \, 0,
\]
which implies $z \eq 0$
and therefore $\overline{\range{\linB_{\sql}} } \eq \vecV$.
But $\range{\linB_{\sql}}$ is closed,
so $\range{\linB_{\sql}} \eq \vecV$.
Finally,
the inequality $\norm{x_f} \leq \theta^{-1} \norm{f}$
follows from the fact that $\norm{\linB_{\sql}^{-1}} \leq \theta^{-1}$.
\qed
\end{comment}

\paragraph{} %- operator version of the Lax-Milgram theorem
There is also an \textit{operator version} of theorem \ref{LAX-MILGRAM_THEOREM}.
Let $\linA \: \hilH \xto \hilH$ be linear and bounded.
We shall denote $\linA$ \textit{coercive} if
\[
\re \scalar{\linA x}{x} \, > \, \theta \norm{x}^2
\]
holds for some $\theta {\,>\,} 0$ and all elements $x \xeof \hilH$.
It should be clear
that the mapping $(x,y) \mapsto \scalar{\linA x}{y}$
defines a coercive form on $\hilH$.
Then the \textsc{Lax-Milgram} theorem implies
invertibility of $\linA$ and boundedness of $\linA^{-1}$,
where $\norm{\linA^{-1}} \leq \theta^{-1}$ holds.

\paragraph{} %- Gelfand triples
%------------------------------------------------------------------------------%
Each Hilbert pair $(\vecV,\hilH)$ with embedding $\iota \: \vecV \xto \hilH$
gives rise to introduce the \textit{transposed pair}
$(\hilH_a^*,\vecV_a^*)$,
where the conjugate linear space $\hilH_a^*$
is continuously and densely embedded into
the conjugate linear space $\vecV_a^*$
by a further embedding $\iota_2 \: \hilH_a^* \xto \vecV_a^*$.
Then,
by identifying the spaces $\hilH$ and $\hilH_a^*$
via the \textsc{Riesz} map,
combining both Hilbert pairs leads to the triple
\begin{equation}\label{GELFAND_TRIPLE}
\vecV \, \overset{\iota_1}{\hookrightarrow} \,
\hilH \, \overset{\iota_2}{\hookrightarrow} \, \vecV_a^*,
\end{equation}
sometimes called a \textit{Gelfand triple},
an \textit{evolution triple} or a \textit{rigged Hilbert space}.
The space $\hilH$ is often denoted the \textit{pivot space}
of \eqref{GELFAND_TRIPLE}.

\subparagraph{}
Gelfand triples own a remarkable property.
Let $(\,\cdot\,,\,\cdot\,)$ and $\scalar{\,\cdot\,}{\,\cdot\,}$
be the inner products in $\vecV$ and $\hilH$,
respectively,
then for fixed chosen $x \in \hilH$
the action of the antilinear linear functional
$f_x := \iota_2 \, x \in \vecV_a^*$ to any element $v \in \vecV$
is represented by the inner product of $x$ and $v$ in $\hilH$,
{\ie} $f_x(v) \eq \scalar{x}{\iota_1 \, v}$ holds for all $v \in \vecV$.

\paragraph{} %-- quasi-coercive forms
%------------------------------------------------------------------------------%
The presented theory of coercive sesquilinear forms in a Hilbert space $\vecV$
is often not sufficent for applications.
Hence it is quite natural to examine sesquilinear forms
also {\wrt} pairs $(\vecV,\hilH)$ of Hilbert spaces.
In applications it is often the case
that sesquilinear forms are merely continuous but not coercive.
In the setting of a Hilbert pair $(\vecV,\hilH)$,
however,
we are able to make a form $\sql$ coercive
by adding an additional term,
which is just a multiplicity of the inner product in $\hilH$.

\subparagraph{}
More precisely,
if $\sql \: \vecV \xtimes \vecV \xto \C$ is continuous
and the slightly modified form
\[
\sql_{\lambda}(u,v) \eqdef
\sql(u,v) \+ \lambda \scalar{\iota_1 u}{\iota_1 v}_H \, \quad
(u,v \xeof \vecV)
\]
is coercive for some $\lambda \xeof \R$,
that is,
there is $\theta {\,>\,} 0$ such that
\[
\re \sql(u,u)  \+
\lambda \norm{\iota_1u}_{\hilH}^2 \, \geq \, \theta \norm{u}_{\vecV}^2
\quad (u \xeof \vecV),
\]
then $\sql$ is called
\textit{quasi-coercive},
\textit{$\hilH$-elliptic}
or simply a \textit{\GARDING} \textit{form}.
It should be clear that
if $\sql_{\lambda_0}$ is coercive for some $\lambda_0 \xeof \R$,
then $\sql_{\lambda}$ is so for $\lambda {\,>\,} \lambda_0$.

\paragraph{} %-- operators associated with H-elliptic forms
%------------------------------------------------------------------------------%
Quasi-coercive sesquilinear forms can be used
to define densely defined and closed linear operators
acting in the pivot space $\hilH$ of a Gelfand triple.

\subparagraph{}
To elaborate this fact,
let $\sql$ be a sesquilinear form defined on a Hilbert space $\vecV$,
which is the left-hand side of a Gelfand triple $(\vecV,\hilH,\vecV_a^*)$
with continuous embeddings
$\iota_1 \: \vecV \xto \hilH$ and
$\iota_2 \: \hilH \xto \vecV_a^*$.
It is our aim to introduce a linear operator $\linA_{\sql}$
with domain space $\hilH$ and range space $\hilH$.
This will be done by restricting the form operator
$\repA \: \vecV \xto \vecV_a^*$ associated with $\sql$
to those elements $u \xeof \vecV$,
for which $\repA \, u \xeof \hilH$.
Indeed,
the domain of definition $\domain{\linA_{\sql}}$ is defined
to be the subspace of those elements $u \xeof \vecV$,
for which $\repA \, u \xeof \iota_2(\hilH) \xsubset \vecV_a^*$ holds,
briefly
\[
\domain{\linA_{\sql}} \eqdef \{\, u \xeof \vecV: \repA \, u \xeof \hilH \,\}.
\]

\subparagraph{} %- variational operator
The set $\domain{\linA_{\sql}}$ can be regarded as a linear subspace
of $\hilH$ by means of the embedding $\iota_1$.
Then for $u \xeof \domain{\linA_{\sql}}$ there is $y \xeof \hilH$
such that $\iota_2 y \eq \repA \, u$.
We set $x := \iota_1 u$ and $\linA_{\sql} x := y$
for $x \xeof \domain{\linA_{\sql}}$
and denote $\linA_{\sql}$
the \textit{operator in $\hilH$ associated with $(\vecV,\sql)$},
or shortly,
the \textit{variational operator} derived from $\sql$.
Another phrase is to say that
$\linA_{\sql}$ is the \textit{part of $\repA$ in $\hilH$}.

\paragraph{} %-- main theorem on H-elliptic forms and associated operators
%------------------------------------------------------------------------------%
The next theorem collects a couple of nice properties 
of the operator $\linA_{\sql}$
associated with a continuous and quasi-coercive form $\sql$.
Besides the \textsc{Lax-Milgram} theorem,
it is our second important result in this section:

\begin{theorem1}\label{QUASI-COERCIVE_FORMS_AND_ASSOCIATED_OPERATORS}
Let $(\vecV,\hilH,\vecV_a^*)$ be a triple of Hilbert spaces
and the sesquilinear form $\sql$ be continuous and quasi-coercive
on $\vecV$ with parameter $\lambda_0 \xeof \R$.
Let $\linA_{\sql}$ be the operator associated with $\sql$.
Then the following assertions are true:

\begin{itemize}[leftmargin=9mm]
\item[(a)]
$\linA_{\sql}$ is closed.

\item[(b)]
$\linA_{\sql}$ is densely defined in $\hilH$.

\item[(c)]
If the embedding of $\vecV$ into $\hilH$ is \textit{compact},
then the operator $\linA_{\sql} + \lambda$
owns a compact inverse $(\linA_{\sql} +\lambda)^{-1}$
for each $\lambda \geq \lambda_0$.

\item[(d)]
The operator $\linA_{\sql} + \lambda$ is m-accretive
for each $\re \lambda \geq \lambda_0$.

\item[(e)]
The operator $\linA_{\sql} + \lambda$ is sectorial
for each $\lambda \geq \lambda_0$.
\end{itemize}
\end{theorem1}

\begin{comment}
\subparagraph{}
To prove that $\linA_{\sql}$ is closed,
let $\{x_n\}_{n \xeof \N}$ be a sequence in $\domain{\linA_{\sql}}$
with limit element $x \xeof \hilH$
and $(\linA_{\sql} x_n$ be the sequence mapped by $\linA_{\sql}$
with limit element $y \xeof \hilH$.
It is to show that $x \xeof \domain{\linA_{\sql}}$ and $\linA_{\sql} x \eq y$.
Let $\abs{x} $ and $\norm{x}$ be the norms in $\vecV$ and $\hilH$,
respectively.
Clearly,
we have
\[
\norm{(\linA_{\sql}+\lambda) x_n - y - \lambda x} \xto 0 
\quad \for n \xto \infty.
\]
In the following we shall use the fact
that for $\lambda {\,>\,} \lambda_0$ the operators
$\linA_{\sql}+\lambda$ and $(\linA_{\sql}+\lambda)^{-1}$
forms linear homeomorphism between the normed spaces
$(\hilH,\norm{\,\cdot\,})$
and $(\domain{\linA_{\sql}},\abs{\,\cdot\,} )$,
so the latter one forms a closed subspace
within the normed space $(\vecV,\abs{\,\cdot\,} )$.
By this reason,
\[
x_n = (\linA_{\sql}+\lambda)^{-1} (\linA_{\sql}+\lambda) x_n
\xto 
(\linA_{\sql}+\lambda)^{-1} (y + \lambda x)
\quad \for n \xto \infty
\]
{\wrt} the norm $\abs{\,\cdot\,} $ in $\domain{\linA_{\sql}}$.
Since $x$ belongs to $\hilH$ and $\domain{\linA_{\sql}}$ is closed,
the element $(\linA_{\sql}+\lambda)^{-1} x$
belongs to $\domain{\linA_{\sql}}$ with $\linA_{\sql} x \eq y$.

\subparagraph{}
To prove that $\domain{\linA_{\sql}}$ is dense in $\hilH$,
suppose that $y \xeof \hilH$ fulfills $\scalar{x}{y} \eq 0$
for all $x \xeof \domain{\linA_{\sql}}$.
It is to show that $y \eq 0$ in this case.
Since $\linA_{\sql}+\lambda$ is surjective,
there must exist $x_0 \xeof \domain{\linA_{\sql}}$ such that
$(\linA_{\sql}+\lambda) x_0 \eq y$.
Then
\[
0 \=
\scalar{\linA_{\sql} x_0}{x} \+ \scalar{\lambda x_0}{x} \=
\sql(x_0,x) \+ \lambda \scalar{x_0}{x}.
\]
Especially for $x \eq x_0$ the last equation leads to
\[
0 \=
\scalar{\linA_{\sql} x_0}{x_0} \+ \scalar{\lambda x_0}{x_0} \,\geq\,
\theta \norm{x_0}^2
\]
for some $\theta {\,>\,} 0$,
due to to quasi-coercivity of $\sql$,
hence $x_0 \eq 0$ and finally $y \eq 0$.

\subparagraph{}
Finally,
if $\iota_1 \: \vecV \xto \hilH$ is compact,
then $\iota_1 \circ (\linA_{\sql}+ \lambda)^{-1} \: \hilH \xto \hilH$
is compact,
too,
since the composition of a compact linear operator
and a continuous linear operator proves to be compact.
\qed.
\end{comment}

\paragraph{} %- operators associated with adjoint and symmetric forms
It is often helpful to consider the \textit{adjoint form} $\sql^*$
of a given sesquilinear form $\sql$ defined by
\[
\sql^* (u,v) \eqdef \overline{\sql(v,u)}
\quad \quad (u,v \xeof \vecV).
\]
Notice that if $\sql$ is continuous,
then $\sql^*$ is so.
Furthermore,
the form $\sql$ is called \textit{symmetric},
if $\sql \eq \sql^*$\footnote{
\,A well-known example of a symmetric form is the inner product
$\scalar{\,\cdot\,}{\,\cdot\,}_{\hilH}$ in a Hilbert space $\hilH$,
which differs from arbitrary symmetric sesquilinear forms
only by positive definiteness.}.
With respect to the variational operators $\linA_{\sql}$ and $\linA_{\sql^*}$,
the following theorem is valid:

\begin{theorem1}\label{QUASI-COERCIVE_FORMS_AND_ASSOCIATED_OPERATORS_II}
If $\sql^*$ is the adjoint form of $\sql$,
then ${\linA_{\sql}}^* \eq \linA_{\sql^*}$ holds.
Moreover,
if $\sql$ is symmetric,
then $\linA_{\sql}$ is self-adjoint.
\end{theorem1}

\paragraph{} %-- generating operators
%------------------------------------------------------------------------------%
Given a Hilbert pair $(\vecV,\hilH)$,
theorem \ref{QUASI-COERCIVE_FORMS_AND_ASSOCIATED_OPERATORS_II}
proves to be helpful to construct self-adjoint linear operators in $\hilH$
by starting with symmetric forms in $\vecV$.
This approach can especially be performed,
if the form is the inner product in $\vecV$.
The operator $\linA$ associated with the inner product
$\scalar{\,\cdot\,}{\,\cdot\,}_{\,\vecV}$,
called the \textit{generating operator} of $(\vecV,\hilH)$,
is coercive and even self-adjoint in $\hilH$
by theorem \ref{QUASI-COERCIVE_FORMS_AND_ASSOCIATED_OPERATORS_II}.

\paragraph{} %- construction of generating operators
In the following we derive the generating operator
of the pair $(\vecV,\hilH)$ of Hilbert spaces in a straightforward way.
First of all,
each fixed $y \in \hilH$
gives rise to a linear functional on $\vecV$ according to
\[
f_y(v) := \scalar{\iota v}{y}_{\hilH}, \quad \quad (v \in \vecV).
\]
The mapping $f_y$ is also bounded,
{\ie} $f_y \in \vecV_a^*$.
Then,
by the \textsc{Riesz} representation theorem,
there is $v_y \in \vecV$ such that
\[
\scalar{v}{v_y} \= f_y(v) \qfa v \in \vecV.
\]
We therefore obtain a mapping $\linB \: \hilH \xto \vecV$,
by $\linB y := v_y$ for $y \in \hilH$,
which proves to be linear and continuous,
and the inner product in $\hilH$ is obtained from that in $\vecV$ by
\[
\scalar{\iota x}{y}_{\hilH} \= \scalar{x}{\linB y}_{\,\vecV}
\quad \quad (x,y \in \hilH).
\]
But the operator $\linB$ is also symmetric and strictly positive.
These properties imply the existence of the inverse operator
\[
\linB^{-1} \: \domain{\linB^{-1} } = \range{\linB} \xsubset \vecV \xto \hilH,
\]
which indeed is not necessarily bounded,
but uniquely defined by the pair $(\vecV,\hilH)$.

\subparagraph{}
The operator $\linB^{-1}$ is symmetric and strictly positive,
too,
hence the square root operator $\linB^{-1/2}$ exists and owns the property
\begin{equation}\label{GENERATING_OPERATOR}
\scalar{\linB^{-1/2} x}{\linB^{-1/2} y}_{\hilH} \= \scalar{x}{y}_{\,\vecV}
\quad \quad (x,y \in \hilH).
\end{equation}

It can be shown by contradiction
that $\domain{\linB^{-1/2}} \eq \vecV$.
Especially for $y=x$ in \eqref{GENERATING_OPERATOR} we obtain the equality
\[
\snorm{\linB^{-1/2}x}_{\hilH} \= \norm{x}_{\,\vecV}
\quad \quad (x \in \vecV).
\]

%------------------------------------------------------------------------------%
\section{Sectorial Forms and m-Sectorial Operators}                            %
%------------------------------------------------------------------------------%
In this section we discuss the notations of \textit{closed sectorial form}
and \textit{m-sectorial operator} and show the relationship between them.
In opposite to the previous section,
we do not consider forms acting on the left-hand space
of a Hilbert pair $(\vecV,\hilH)$,
but assume the forms to act on linear subspaces
lying dense in a Hilbert space $\hilH$.

\subparagraph{} %-- sectorial forms
%------------------------------------------------------------------------------%
Let $\sql \: \domain{\sql} \xtimes \domain{\sql} \xto \C$
be a sesquilinear form
defined on a dense linear subspace $\domain{\sql}$ of $\hilH$,
called the \textit{domain of definition} of $\sql$.
Like for linear operators,
the \textit{numerical range} of $\sql$ is defined by
\[
\Theta(\sql) \eqdef
\{\, \sql(x,x): x \xeof \domain{\sql}, \, \norm{x}=1 \,\}.
\]
The form $\sql$ is said to be \textit{sectorial},
if $\Theta(\sql)$ is contained in a sector $\Sigma_{\,\psi,\omega}$,
that is,
\begin{equation}\label{SECTOR_CONDITION}
\Theta(\sql) \, \xsubset \,
\Sigma_{\,\psi,\omega} \eqdef
\{\zeta \xeof \C: \zeta \neq \omega, arg(\zeta-\omega) < \psi \},
\end{equation}
where $\omega \xeof \R$ is the vertex
and $0 \leq \psi {\,<\,} \pi/2$ the semi-angle of the sector.
We also say that $\linA$ is $(\psi,\omega)$-sectorial,
if \eqref{SECTOR_CONDITION} holds.
Especially for $\psi \eq 0$ let $\Sigma_{\,0,\omega} := (\omega,\infty)$,
and for $\psi \eq \pi/2$ we have $\Sigma_{\,\pi/2,0} \eq \C_+$.
Notice that sectors are symmetric {\wrt} the real line.

\subparagraph{} %- half-bounded forms
Now,
if $\sql$ is a symmetric and $\sql \eq \sql^*$,
respectively,
then $\Theta(\sql) \xsubset \R$ and $\omega \leq \inf \Theta(\sql)$, 
so in this case we denote $\sql$ to be 
\textit{half-bounded} or \textit{bounded below}
and call $\omega$ the lower bound of $\sql$.
Moreover,
if $\sql$ is half-bounded and $(\psi,\omega)$-sectorial,
then the estimates
\[
\begin{cases}[1.2]
\quad \re \sql(x,x) \, \geq \, \omega \norm{x}^2 \\
\quad \im \sql(x,x) \, \leq \,
      \tan \psi \cdot (\re \sql(x,x) - \omega \norm{x}^2)\\
\end{cases}
\]
hold for all $x \xeof \domain{\sql}$.

\paragraph{} %-- closed and closable sectorial forms
%------------------------------------------------------------------------------%
Let $\re \sql := \frac{1}{2}(\sql + \sql^*)$
be the \textit{real part} of $\sql$.
Then for every $\alpha {\,>\,} 0$ the slightly modified form
$[\cdot,\cdot]_{\alpha} := \re \sql - \omega + \alpha$
constitutes an inner product on the linear space $\domain{\sql}$.
It turns out
that the norms $\abs{\cdot} _{\alpha}$
derived from the inner products $[\cdot,\cdot]_{\alpha}$
are equivalent to each other.
We especially chose $\alpha \eq 1$,
pay attention to the inner product
\[
[x,y]_1 \= \re \sql(x,y) - \omega \scalar{x}{y} + \scalar{x}{y},
\qfa x,y \xeof \domain{\sql}
\]
and distinguish between the two cases {\wrt} completeness of the norm
\[
\abs{x} _1 \eqdef {[x,x]_1}^{1/2}
\quad \quad (x \xeof \domain{\sql}).
\]
It is in both cases our goal to obtain a Hilbert space $\hilH_{\sql}$
with inner product $[x,y]_1$ and norm $\abs{x} _1$
($x,y \xeof \hilH_{\sql}$),
together with a linear embedding $\iota \: \hilH_{\sql} \xto \hilH$,
which should be continuous if possible,
so $(\hilH_{\sql},\hilH)$ hopefully becomes into a pair of Hilbert spaces.

\begin{itemize}[leftmargin=4mm]
\item[1.]
The normed linear space $(\domain{\sql},\abs{\,\cdot\,} _1$)
is \textit{complete}.\\
We denote the form $\sql$ \textit{closed} in this case,
and $\domain{\sql}$ becomes into a Hilbert space in its own right,
hence we may define $\hilH_{\sql} := \domain{\sql}$ in this case. 
Notice that the norm $\abs{\,\cdot\,} _1$ on $\hilH_{\sql}$
is stronger than the norm $\norm{\,\cdot\,}$ in $\hilH$
restricted to $\hilH_{\sql}$,
so the natural inclusion mapping $\iota \: \hilH_{\sql} \xto \hilH$
is in fact a continuous mapping.

\subparagraph{} %- core of a closed sesquilinear form
Let us in addition mention that
a linear subspace $D \xsubset \domain{\sql}$
is denoted a \textit{core} of $\sql$,
if $\sql$ is the closure of the restriction $\sql_{|D}$.

\item[2.]
The normed linear space $(\domain{\sql},\abs{\,\cdot\,} _1)$
is \textit{not} complete rsp. $\sql$ is not closed.\\
We define the Hilbert space $\hilH_{\sql}$ as the completion
of $\domain{\sql}$ {\wrt} the norm $\abs{\,\cdot\,} _1$.
Moreover,
Cauchy sequences in $\hilH_{\sql}$ are also Cauchy sequences in $\hilH$,
hence for every $x \xeof \hilH_{\sql}$
there is $x' \xeof \hilH$ such that $x' \eq \iota x$
for a linear map $\iota \: \hilH_{\sql} \xto \hilH$.

\subparagraph{}
We want to extend the form $\sql$ from $\domain{\sql}$ to $\hilH_{\sql}$
as well,
and this will be done by defining
\[
\overline{\sql}(x,y) \eqdef \lim_{\,n \xto \infty} \sql(x_n,y_n)
\]
for $x,y \xeof \hilH_{\sql}$ and arbitrary sequences
$\{x_n\}_{n \xeof \N}, \{y_n\}_{n \xeof \N} \xsubset \domain{\sql}$
with $x_n \xto x$ and $y_n \xto y$.
It turns out that
the definition of $\overline{\sql}(x,y)$ is independent of
the sequences $\{x_n\}_{n \xeof \N}$ and $\{y_n\}_{n \xeof \N}$,
so this value is well-defined for $x,y \xeof \hilH_{\sql}$.
\end{itemize}

\paragraph{} %- closable sectorial forms
In summary,
we obtain in both cases a Hilbert space $\hilH_{\sql}$ with inner product
$[\cdot,\cdot]_1$ and norm $\abs{\,\cdot\,} _1$
together with a linear embedding $\iota \: \hilH_{\sql} \xto \hilH$,
but which in the second case is not necessarily continuous.

\subparagraph{} 
Regarding the second case,
the form $\sql$ is said to be \textit{closable},
if the mapping $\iota$ is continuous.
This definition especially means that every closed form is closable.
If $\sql$ is closable,
then the final result in both cases will be
that there is a closed form,
namely $\sql$ itself in the first and $\overline{\sql}$ in the second case,
together with a continuous embedding of $\hilH_{\sql}$ into $\hilH$.

\paragraph{} %- sectorialness
We turn back to the concept of \textit{sectorialness}
of a closed linear operator.
This class of operators has been introduced and first studied
by \textsc{T.\,Kato}.
Let again $\Sigma_{\,\psi,\omega}$ be an open sector
with vertex $\omega \xeof \R$ and semi-angle $0 {\,<\,} \psi \leq \pi$.
In his book \cite{Ka},
\textsc{Kato} denotes a linear operator $\linA$ acting in Hilbert space $\hilH$
to be \textit{sectorial},
if the numerical range $\Theta(\linA)$ is a subset of the sector
$\Sigma_{\,\psi,\omega}$
for some $\omega \xeof \R$ and $0 \leq \psi {\,<\,} \pi/2$.
Furthermore,
\textsc{Kato} calls the operator $\linA$ \textit{m-sectorial},
if it is sectorial and-quasi-m-accretive at the same time,
that is,
$\linA + \omega$ is m-accretive for some $\omega \xeof \R$.

\subparagraph{} %- relationship between closable sectorial forms and operators
We finally elaborate the correspondence
between closed sectorial forms and m-sectorial operators.
It should be clear
that if $\linA$ is a m-sectorial operator in $\hilH$
with domain of definition $\domain{\linA}$,
then the sesquilinear form
\[
\sql_{\linA}: \;
\begin{cases}[1.2]
\quad \domain{\sql_{\linA}} \eqdef \domain{\linA}, \\
\quad \sql_{\linA}(x,y)     \eqdef \scalar{\linA \,x}{y}
\; \for x,y \xeof \domain{\sql_{\linA}} \\
\end{cases}
\]
is sectorial and closed.

\subparagraph{} %-- Kato's first representation theorem
%------------------------------------------------------------------------------%
Reversely,
\textsc{Kato}\textit{'s first representation theorem} claims that
a densely defined and closed sectorial form gives rise
to a m-sectorial linear operator:

\begin{theorem1}\label{FIRST_REPRESENTATION_THEOREM}
Let $\sql$ be a densely defined and closed sectorial form in $\hilH$
with vertex $\omega$, semi-angle $\psi$ and domain of definition $\hilH_{\sql}$,
and let $\iota$ be the continuous embedding from $\hilH_{\sql}$ into $\hilH$.
Then there is a m-sectorial operator $\linA_{\sql}$ in $\hilH$,
whose domain of definition $\domain{\linA_{\sql}}$ is defined as follows:
Any element $x \xeof \hilH$ belongs to $\domain{\linA_{\sql}}$,
if $x \eq \iota x_0$ for some $x_0 \xeof \hilH_{\sql}$ and
$\sql(x_0,z) \eq \scalar{y}{\iota z}$ holds
for some $y \xeof \hilH$ and all $z \xeof \hilH_{\sql}$,
which suggests to define $\linA_{\sql} \, x := y$.

\paragraph{} %- properties of an operator associated with a sectorial form
The operator $\linA_{\sql}$ owns a couple of desirable properties:

\begin{itemize}[leftmargin=9mm]
\item[(a)]
$\domain{\linA_{\sql}}$ is dense in $\hilH$
and $\iota^{-1} \domain{\linA_{\sql}}$ is dense in $\hilH_{\sql}$.

\item[(b)]
If $x \xeof \hilH_{\sql}$, $y \xeof \hilH$
and $\sql(x,y) \eq \scalar{\iota x}{y}$ holds for all $x \xeof D$,
where $D$ is a core of $\sql$,
then $\iota x \xeof \domain{\linA_{\sql}}$ and $y \eq \linA (\iota x)$.

\item[(c)]
$\Theta(\linA_{\sql})$ is a dense subset of $\Theta(\sql)$.

\item[(d)]
The adjoint form $\sql^*$ is also closed and sectorial
with vertex $\omega$ and semi-angle $\psi$.

\item[(e)]
The adjoint ${\linA_{\sql}}^*$ of $\linA_{\sql}$
is the $m$-sectorial operator associated with the adjoint form $\sql^*$.
\end{itemize}
\end{theorem1}

\subparagraph{}
Thus m-sectorial linear operators are in one-to-one correspondence
with closed sectorial forms.

\paragraph{} %- Friedrichs extensions
We want to apply the obtained results to the class of half-bounded operators.
They distinguish from other operators by symmetry and the fact
that there is $\theta \xeof \R$ such that
\[
\scalar{\linA \,x}{x} \, \geq \, \theta \norm{x}^2
\qfa x \xeof \domain{\linA},
\]
or equivalently,
if $\Theta(\linA) \xsubset [\theta,\infty)$ for some $\theta \xeof \R$.
Recall also
that the operator $\linA$ is called \textit{positive}
(\textit{strictly positive}),
if $\theta \eq 0$ ($\theta {\,>\,} 0$) can be chosen.
In case of $\theta {\,>\,} 0$,
the {Cauchy-Schwarz} inequality implies
that $\linA$ is also bounded below,
hence $\linA$ is surely injective.
In addition,
if $\linA$ is closed,
then the range $\range{\linA}$ is closed in $\hilH$,
which follows from the bounded-below theorem,
but $\range{\linA}$ may of course be a proper subspace of $\hilH$.

\paragraph{} %-- the Friedrichs extension for half-bounded operators
%------------------------------------------------------------------------------%
The second main result in this section is given by the fact
that each densely defined and half-bounded linear operator $\linA$
owns at least one \textit{self-adjoint extension}.

\subparagraph{} %- energetic space
In fact,
we first have to notice that
if $\linA$ is half-bounded with $\theta \xeof \R$,
then it is $(0,\theta)$-sectorial,
and the sesquilinear form $\sql(x,y) := \scalar{\linA x}{y}$
is sectorial and densely defined.
The form $\sql$ is also closable,
so we may switch to the sesquilinear form $\overline{\sql}$ defined on the
Hilbert space $\hilH_{\sql}$,
called the \textit{energetic space} of the operator $\linA$.
Then by theorem \ref{FIRST_REPRESENTATION_THEOREM}
there is a m-sectorial and even self-adjoint linear operator $\linA_{\sql}$,
that forms an extension of the half-bounded operator $\linA$,
called the \textsc{Friedrichs} \textit{extension} of $\linA$.

\subparagraph{} %- properties of the Friedrichs extension
The domain of definition $\domain{\linA_{\sql}}$ of $\linA_{\sql}$
coincides with the linear subspace $\hilH_{\sql} \, \cap \, \domain{\linA^*}$.
Moreover,
$\linA_{\sql}$ owns a useful property:
if one adds to the half-bounded operator $\linA$ the dilatation 
$x \xeof \hilH \mapsto \theta x$ with $\theta \xeof \R$,
then the \textsc{Friedrichs} extensions of $\linA$ and $\linA + \theta$
are related to each other by the identity
$(\linA + \theta)_{\sql} \eq \linA_{\sql} + \theta$.

%------------------------------------------------------------------------------%
\section{Variational Evolution Equations}                                      %
%------------------------------------------------------------------------------%
In this section we want to sketch the variational approach
to first and second order evolution equations,
which has been developed mainly by \textsc{J.-L.\,Lions}.
It can be seen as the straightforward extension of
the {Lax-Milgram} theory from the stationary to the evolutional case.
But before starting with this,
we have to introduce the class of \textit{Bochner spaces},
which plays a fundamental role in this subject.

\paragraph{} %-- vector-valued functions and Bochner spaces
%------------------------------------------------------------------------------%
The {Bochner spaces} $L^p(\Omega,\banX)$ are nothing but
Lebesgue spaces for vector-valued functions $f \: \Omega \xto \banX$,
where $(\Omega,\salgB,\mu)$ is a measure space and $\banX$ is a Banach space.
To define them,
we first have to introduce concepts of measurability and integration
for vector-valued functions,
which form generalizations of the scalar-valued case.

\subparagraph{} %- measurable functions
Fist notice that
a function $f \: \Omega \xto \banX$ is called \textit{weakly measurable},
if $\phi \circ f \: \Omega \xto \F$ is measurable
for every $\phi \xeof \banX^*$.
Then every weakly measurable function $f$ is called \textit{Bochner-integrable}
if
\[
\int_{\Omega} \norm{f} \, d\mu \, < \, \infty.
\]
For any $p \geq 1$
we denote $L^p(\Omega,\banX)$ the space of functions $f \: \Omega \xto \banX$
having Bochner-integrable powers $\abs{f} ^{p}$.
The elements of this space should of course be understood
as equivalence classes of functions
being identical almost everywhere in $\Omega$.
Endowed with the norm
\[
f \mapsto \norm{f}_p \eqdef (\int_{\Omega} \, \norm{f} ^p \, d\mu)^{1/p},
\]
the Bochner spaces $L^p(\Omega,\banX)$ are complete
and thus constitute Banach spaces.
A good treatment of the theory of Bochner spaces is given in \cite{DU}
within section 2.

\subparagraph{} %- Hilbert-Bochner spaces 
If $\Omega \eq (a,b)$ is an interval and $\banX \eq \hilH$ is a Hilbert space,
then the Bochner space $L^2(a,b;\hilH)$ can be endowed with the inner product
\[
\scalar{f}{g} \eqdef \int_a^b \scalar{f(t)}{g(t)} \,dt,
\quad \quad
f,g \xeof L^2(a,b;\hilH)
\]
and the norm is derived as usual by $\norm{f}_2 \eq \scalar{f}{f}^{1/2}$.
This space is Hilbertian,
too,
so it is denoted a \textit{Hilbert-Bochner space}.
If $(a,b)$ is \textit{finite},
then the inclusion $L^2(a,b;\hilH) \xsubset L^1(a,b;\hilH)$ is valid.

\paragraph{} %-- weak derivatives of vector-valued functions
%------------------------------------------------------------------------------%
Supposed $f \xeof L^2(a,b;\hilH)$,
a further weakly measurable function $g \: (a,b) \xto \hilH$
is said to be the \textit{weak derivative} of $f$,
if
\[
\int_a^b g \phi \,dx \= - \int_a^b f \phi' \,dx
\]
holds for each $\phi \xeof C_c^{\infty}(a,b,\banX)$,
the space of indefinitely differentiable functions
with compact support in $(a,b)$.
If such a function $g$ exists,
we define $f' := g$ in this case,
which should not be confused with the classical derivative of $f$.
Note that
if $f$ is classical differentiable,
then the weak derivative of $f$ exists and coincides
with the classical derivative.

\subparagraph{} %-- Hilbert-Bochner-Sobolev spaces
%------------------------------------------------------------------------------%
Having the notation of weak derivative for vector-valued functions available,
we may introduce the \textit{Hilbert-Bochner-Sobolev space}
\[
H^1(a,b;\hilH) \eqdef
\{f \xeof L^2(a,b;\hilH): f' \xeof L^2(a,b;\hilH) \},
\]
which is a Hilbert space again with inner product given by
\[
\scalar{f}{g}_1 \eqdef \int_a^b \scalar{f}{g} \, dt
\+
\int_a^b \scalar{f'}{g'} \, dt\,,
\quad f,g \xeof H^1(a,b;\hilH).
\]

\subparagraph{} %- mixed Bochner-Sobolev spaces
Finally,
let $\hilH_1, \hilH_2$ be \textit{real} Hilbert spaces,
where $\hilH_1$ is continuously embedded into $\hilH_2$.
Then there is also a continuous embedding
\[
H^1(a,b;\hilH_1) \, \overset{\kappa}{\hookrightarrow} \, H^1(a,b;\hilH_2).
\]
Consequently,
given a Hilbert space triple 
$
\vecV \, {\overset{\iota_1}{\hookrightarrow}} \,
\hilH \, {\overset{\iota_2}{\hookrightarrow}} \vecV^*,
$
we may regard the space $H^1(a,b;\vecV)$
as a linear subspace of $H^1(a,b;\vecV^*)$,
which enables us to define the \textit{mixed Bochner-Sobolev space}
\[
H^1(a,b;\vecV,\vecV^*) \eqdef
\{ f \xeof L^2(a,b;\vecV): f' \xeof L^2(a,b;\vecV^*) \}
\]
or equivalently
\[
H^1(a,b;\vecV,\vecV^*) \eqdef L^2(a,b;\vecV) \, \cap H^1(a,b;\vecV^*).
\]
Then the mapping
\[
(f,g) \mapsto \scalar{f}{g}_1 \eqdef
\int_a^b \, \scalar{f(t)}{g(t)}_{\vecV} \, dt
\+
\int_a^b \, \scalar{f'(t)}{g'(t)}_{V^*} \, dt
\]
fulfills the properties on an inner product
and makes $H^1(a,b;\vecV,\vecV^*)$ into a Hilbert space
having the following properties:

\begin{itemize}[leftmargin=9mm]
\item[(a)]
Each $f \xeof H^1(a,b;\vecV,\vecV^*)$
is $\lambda$-nearly everywhere identical
to a continuous function $f_c \: [a,b] \xto \hilH$,
and by identifying the functions $f$ and $f_c$
the space $H^1(a,b;\vecV,\vecV^*)$ can continuously be embedded
into the space $C([a,b],\hilH)$.

\item[(b)]
The space $C^{\infty}([a,b],\vecV)$ is dense in $H^1(a,b;\vecV,\vecV^*)$.

\item[(c)]
The following product rule is valid for each $t \xeof (a,b)$:
\[
\frac{d}{dt} \norm{u(t)} ^2
\=
2 \re \, \scalar{u'(t)}{u(t)}_1
\fa u \xeof H^1(a,b;\vecV,\vecV^*).
\]
\end{itemize}

\paragraph{} %-- the variational approach to evolutional problems
%------------------------------------------------------------------------------%
We finally turn to abstract variational problems
of parabolic and hyperbolic type.

\subparagraph{}
Let $\T {\,>\,} 0$ and denote $\lambda$ the Lebesgue-Borel measure on $[0,\T)$.
By an \textit{abstract variational problem of parabolic type}
in a real Hilbert space triple $(\vecV,\hilH,\vecV^*)$
we mean a $\lambda$-nearly everywhere in $(0,\T)$ valid
linear operator-differential equation
\begin{equation}\label{ABSTRACT_PARABOLIC_VARIATIONAL_PROBLEM}
\begin{cases}[1.2]
\quad u'(t) \+ \repA(t) \, u(t) \= f(t), \\
\quad u(0) = u_0, \\
\end{cases}
\end{equation}
where $u_0 \xeof \hilH$,
$f \: [0,\T) \xto \vecV^*$ and $\{\repA(t)\}_{t \xeof [0,\T)}$
is a family of linear operators $\repA(t) \: \vecV \xto \vecV^*$.

\paragraph{} %-- well-posedness theorems for variational initial value problems
%------------------------------------------------------------------------------%
Let $\sql(t)$ be the bilinear form associated with the operator $\repA(t)$.
We call $u \xeof H^1(0,\T;\vecV,\vecV^*)$
a \textit{weak solution} of \eqref{ABSTRACT_PARABOLIC_VARIATIONAL_PROBLEM},
if $u(0) \eq u_0$ and
\[
\scalar{u'(t)\,}{v(t)} \+ \sql(t,u(t),v(t)) \= f(t)(v)
\]
holds for all $v \xeof \vecV$ and $\lambda$-nearly everywhere $t \xeof [0,\T)$,
where by $\scalar{\cdot}{\cdot}$ is meant the inner product in $\hilH$.

\subparagraph{} %- variational initial value problems of first order
Regarding the existence and uniqueness of a {weak solution}
for problem \eqref{ABSTRACT_PARABOLIC_VARIATIONAL_PROBLEM},
the following theorem is of fundamental meaning:

\begin{theorem1}\label{LIONS_THEOREM_I}
Let $\{\sql(t)\}_{t \xeof [0,\T)}$ fulfill the following conditions:

\begin{itemize}[leftmargin=9mm]
\item[(a)]
the mapping $t \xeof [0,\T) \mapsto \sql(t,u,v)$ is Lebesgue-measurable
for every fixed pair $(u,v) \xeof \vecV \xtimes \vecV$,

\item[(b)]
$\sql$ is \textit{uniformly} bounded in $[0,\T)$,
that is,
there is a constant $M \geq 0$ such that
\[
\abs{\sql(t,u,v)} \, \leq M \cdot \norm{u}_{\vecV} \norm{v}_{\vecV}
\qfa u,v \xeof \vecV, \; t \xeof [0,\T)\,,
\]
\item[(c)]
$\sql(t)$ is quasi-coercive for each $t \xeof [0,\T)$.
\end{itemize}

Then for each $u_0 \xeof \hilH$ and each $f \xeof L^2(0,\T;\vecV^*)$
there is exactly one solution $u \xeof H^1(0,\T;\vecV,\vecV^*)$
of \eqref{ABSTRACT_PARABOLIC_VARIATIONAL_PROBLEM},
which continuously depends on the data $u_0$ and $f$.
\end{theorem1}

\subparagraph{}
The proof of this theorem has been published in \cite{Li2}
as \textit{Th\'{e}or\`{e}me 1.1}, page 46.
Both terms $u'$ and $\repA(t) u$
in \eqref{ABSTRACT_PARABOLIC_VARIATIONAL_PROBLEM}
belong to the same space $L^2(0,\T;\vecV^*)$ as $f$,
which gives rise to say
that the problem has \textit{maximal regularity} in $\vecV^*$.

\paragraph{} %- variational initial value problems of second order
The concept of weak solution makes also sense
for linear evolution equations of second order.
An \textit{abstract hyperbolic variational problem}
{\wrt} the Hilbert space triple $(\vecV,\hilH,\vecV^*)$ 
and a family of bilinear forms
$\sql(t) \: [0,\T] \xtimes \vecV \xtimes \vecV \xto \R$ is given by
\begin{equation}\label{ABSTRACT_HYPERBOLIC_VARIATIONAL_PROBLEM}
\begin{cases}[1.2]
\quad \scalar{u''(t)\,}{v(t)} + \sql(t,u(t),v(t)) = \scalar{f(t)}{v} 
\quad (v \xeof \vecV), \\
\quad u(0)  = u_0 \xeof \vecV, \\
\quad u'(0) = u_1 \xeof \hilH. \\
\end{cases}
\end{equation}

\paragraph{}
Our next theorem shows that
problem \ref{ABSTRACT_HYPERBOLIC_VARIATIONAL_PROBLEM}
is \textit{weakly solvable} under a couple of conditions
on the involved bilinear forms:

\begin{theorem1}\label{LIONS_THEOREM_II}
Let $\{\sql(t)\}_{t \xeof [0,\T)}$ be a family
of bilinear forms on $\vecV$ such that

\begin{itemize}[leftmargin=9mm]
\item[(a)]
$t \mapsto \sql(t)$ is \textit{uniformly} bounded in $[0,\T)$,

\item[(b)]
$\sql(t)$ is $\hilH$-elliptic for each $t \xeof [0,\T)$,

\item[(c)]
$\sql(t)$ satisfies the symmetry condition
$\sql(t,u,v) \eq \sql(t,v,u)$
for all $u,v \xeof \vecV, \; t \xeof [0,\T)$,

\item[(d)]
the mapping $t \mapsto \sql(t,u,v)$ is continuously differentiable in $[0,\T)$
for each fixed $(u,v) \xeof \vecV \xtimes \vecV$,

\item[(e)]
there is $M_1 \geq 0$ such that
\[
\abs{ \frac{d}{dt} \sql(t,u,v)} \, \leq \, 
M_1 \cdot \norm{u}_{\vecV} \norm{v}_{\vecV}
\qfa \, t \xeof [0,\T).
\]
\end{itemize}
Then the initial value problem \eqref{ABSTRACT_HYPERBOLIC_VARIATIONAL_PROBLEM}
has a unique weak solution $u \xeof H^1(0,\T;\vecV,\vecV^*)$
for all $u_0 \xeof \vecV$,
        $u_1 \xeof \hilH$ and
        $f \xeof L^2(0,\T;\hilH)$.
\end{theorem1}

The symmetry condition (c) in \ref{LIONS_THEOREM_II} is essential
to enable the proof.

\subparagraph{}
Finally,
the solution map
\[
R\:L^2(0,\T;\hilH) \xtimes \vecV \xtimes \hilH \xto L^2(0,\T;\hilH)
\]
given by $(f,u_0,u_1) \mapsto u$ is linear in each argument and continuous,
thus problem \eqref{ABSTRACT_HYPERBOLIC_VARIATIONAL_PROBLEM} proves
to be well-posed.

%------------------------------------------------------------------------------%
\section{Applications to Linear Differential Operators}                        %
%------------------------------------------------------------------------------%
This section deals with applications of abstract functional analysis
to linear differential operators and the problems associated with them.
We shall start with the concept
of a \textit{linear partial differential operator},
which enables us to consider specific problems of type (P-1), (P-2) and (P-3)
mentioned in section 3.

\paragraph{} %-- linear partial differential operators
%------------------------------------------------------------------------------%
Remember first the notation of \textit{multi-index},
which is nothing but a $d$-tuple $\alpha \eq (\alpha_1,\ldots,\alpha_d)$
with values $\alpha_i \xeof \N_0$
and with modulus $\abs{\alpha} := \alpha_1 + \ldots + \alpha_d$.
Multi-indices are used to shorten multiple partial derivatives
$
\PA f \eqdef \partial_1^{\alpha_1}
         \partial_2^{\alpha_2} \ldots
         \partial_d^{\alpha_d} f
$
of a sufficiently often differentiable and complex-valued function $f$,
but also to denote the multiple powers
$\xi^{\alpha} := \xi_1^{\alpha_1}
\cdot \xi_2^{\alpha_2} \cdot \ldots \cdot \xi_d^{\alpha_d}$
of a vector $\,\xi \eq (\xi_1,\xi_2,\ldots,\xi_d) \xeof \C^d$.

\subparagraph{}
By the usage of multi-indices,
a \textit{linear partial differential operator} $\A$ of even order $2m$
is defined by
\begin{equation}\label{EVEN_ORDER_LPDO}
\A(x,\,\cdot\,) \eqdef
\sum_{\abs{\alpha} \leq 2m} c_{\alpha}(x) \, \PA
\end{equation}
for $x \xeof \clG$,
the closure of a \textit{domain} $\G$,
{\ie} an open and connected subset of an Euclidean space $\Rd$.
The domain $\G$ may be bounded or unbounded (and even the whole of $\Rd$).
We shall call $\A$ \textit{essentially real},
if the coefficients $c_{\alpha}$ of highest order $\abs{\alpha} \eq 2m$
are real-valued.
Moreover,
the expression $\A$ is said to be
an \textit{operator with constant coefficients},
if all functions $c_{\alpha}$ are constant in $\clG$.

\paragraph{} %-- free space and boundary value problems for partial differential equations
%------------------------------------------------------------------------------%
In case of $\G \eq \Rd$ respectively $\bndG \eq \emptyset$,
where $\bndG$ is the boundary of $\G$,
we shall denote the partial differential equation $\A \, u \eq f$
with $\A$ given by \eqref{EVEN_ORDER_LPDO} a \textit{free space problem}.

\subparagraph{} %- boundary value problems
However,
partial differential equations considered in domains with non-empty boundary
are usually furnished with \textit{boundary conditions},
which fix possible solutions and their normal derivatives up to a certain order
in the boundary $\bndG$ of the underlying domain $\G \xsubset \Rd$.
Thus,
when having in mind linear equations of the form $\A\,u \eq f$
with imposed boundary conditions,
we shall speak of \textit{linear boundary value problems}.

\subparagraph{} %- general form of boundary value problems
The most general form of a boundary value problem
posed in a domain $\G \subsetneq \Rd$ is of the form
\begin{equation}\label{LEBVP}
\begin{cases}[1.2]
\quad \A(x,u)     \= f(x)   & (x \xeof \G),\\
\quad \BDO_k(x,u) \= g_k(x) & (x \xeof \bndG \and k \eq 0,\ldots,m{-}1),\\
\end{cases}
\end{equation}
where $\A$ is a {linear} differential operator of order $2m$
and $\BDO_0,\ldots,\BDO_{m-1}$ are \textit{boundary operators} of order $m_k$,
each of them defined by
\[
\BDO_k(x,u) \=
\sum_{\abs{\alpha} \leq m_k} b_{k,\alpha}(x) \, \PA u(x).
\]

\subparagraph{}
To introduce the boundary operators $\BDO_k$,
it is sufficent
that the real or complex-valued functions $b_{k,\alpha}$
are defined in a neighbourhood of the boundary $\bndG$ only.
Then each operator $\BDO_k$ is a linear differential operator in its own right
of order $m_k {\,\leq\,} 2m{-}1$,
which maps the function $u \: \clG \xto \C$
to a function $\BDO_k u \: \bndG \xto \C$
belonging to a so-called \textit{boundary space} $B_k(\bndG)$.

\paragraph{} %- homogeneous and normal boundary operators
The collection $\{B_k\}_{k=0}^{m-1}$ of boundary operators in \eqref{LEBVP}
is said to be \textit{homogeneous}
if $\BDO_k(x,\cdot) \eq 0$ for $x \xeof \bndG$ and $k \eq 0,\ldots,m{-}1$.
Moreover,
it is denoted \textit{normal} if

\begin{itemize}[leftmargin=15mm,itemsep=2pt]
\item[(B-1)]
$m_i \neq m_j$ for $i \neq j$,

\item[(B-2)]
for all $\xi \neq 0$ being normally positioned in $x \xeof \bndG$
we have
\[
\sum_{\abs{\alpha} = m_k} b_{k,\alpha}(x) \, \xi^{\alpha} \neq 0
\quad \quad (k=0,\ldots,m{-}1).
\]
\end{itemize}

\subparagraph{} %- homogeneous Dirichlet systems
Finally,
our collection of boundary operators is called a \textit{Dirichlet system},
if it is normal and the values $m_k$ form a permutation
of the set $\{0,\ldots,m{-}1\}$,
so one may reorder the operators $\BDO_k$ to obtain $m_k \eq k$
($0 \leq k \leq m{-}1$).
The most regarded set of boundary conditions in applications
is the \textit{homogeneous Dirichlet system} defined by
\begin{equation}\label{HOMOGENEOUS_DIRICHLET_SYSTEM}
\BDO_k(x,u) \eqdef (\partial_{\nu}^k u)(x) \= 0
\end{equation}
for $x \xeof \bndG$ and $k \eq 0,\ldots,m{-}1$,
where $(\partial_{\nu}^k u)(x)$ is the $k$-th derivative of a function $u$
along the outer normal vector $\nu$ in $x \xeof \bndG$.

\subparagraph{} %- classical solutions of BVP endowed with a Dirichlet system
If the set of boundary conditions in \eqref{LEBVP} is a Dirichlet system,
then we denote a solution $u \: \clG \xto \C$
of this problem \textit{classical},
if
\begin{itemize}[leftmargin=9mm]
\item[(a)] $u \xeof C^{2m}(\G) \, \cap \, C^{m-1}(\clG)$,
\item[(b)] $u$ fullfills $\A(x,u) \eq f(x)$ in each $x \xeof G$ and
\item[(c)] the $m{-}1$ equations $\BDO_k(x,u) \eq g_k(x)$ hold
           in every $x \xeof \bndG$.
\end{itemize}

\paragraph{} %- boundary condition for second-order problems
Especially to each \textit{second-order} boundary value problem
there is assigned exactly one boundary operator $\BDO_0$,
hence problem \eqref{LEBVP} becomes into
\begin{equation}\label{SECOND_ORDER_BOUNDARY_VALUE_PROBLEM}
\begin{cases}[1.2]
\quad \A(x,u)     = f(x)   & (x \xeof \G),\\
\quad \BDO_0(x,u) = g_0(x) & (x \xeof \bndG).\\
\end{cases}
\end{equation}

\subparagraph{} %- enumeration of boundary conditions for second-order problems
Boundary value problems of the form \eqref{SECOND_ORDER_BOUNDARY_VALUE_PROBLEM}
and governed by a linear differential operator $\A$ of second order
are obtained by various boundary operators $\BDO_0$.
We denote

\begin{itemize}[leftmargin=18mm,itemsep=2pt]
\item[(BC-1)]
\;$\BDO_0(x,u) \eq u(x)$ for $x \xeof \bndG$ a \textit{Dirichlet condition},

\item[(BC-2)]
\;$\BDO_0(x,u) \eq u_{\nu}(x)$ for $x \xeof \bndG$ a \textit{Neumann condition},

\item[(BC-3)]
\;$\BDO_0(x,u) \eq \alpha(x)u(x)+\beta(x) u_{\nu}(x) \eq 1$ for $x \xeof \bndG$ 
a \textit{Robin condition},
where the functions $\alpha, \beta$ must fulfill 
\[
\alpha(x)^2 + \beta(x)^2 \neq 0 \quad \quad (x \in \bndG).
\]
\end{itemize}

\subparagraph{}
Due to physical reasons,
one often has to split the boundary set $\bndG$
into disjoint subsets $\Gamma_1,\Gamma_2,\Gamma_3$,
on which the conditions (BC-1), (BC-2) and (BC-3) are prescribed.
We then shall speak of \textit{mixed} boundary value problems
for the second-order operator $\A$.

\paragraph{} %-- operators of elliptic, parabolic and hyperbolic type
%------------------------------------------------------------------------------%
The concept of \textit{ellipticity} has been shown to be helpful in order
to classify linear partial differential expressions
by a simple algebraic criterion.
Let $\A$ be a linear differential operator of order $2m$ be given
by \eqref{EVEN_ORDER_LPDO} and let
\[
p(x,\xi) \eqdef \sum_{\abs{\alpha} = 2m} c_{\alpha}(x) \, \xi^{\alpha}
\quad \for \xi = (\xi_1,\ldots,\xi_d) \in \Rd
\]
be the \textit{characteristic function} or the \textit{main symbol} of $\A$,
which is built up by the summands of highest order in \eqref{EVEN_ORDER_LPDO}.

\begin{itemize}[leftmargin=10mm]
\item[(E-1)]
The expression $\A$ is called \textit{elliptic} in $x \xeof \clG$,
if
\[
\abs{p(x,\xi)} \, \neq \, 0 \qfa \xi \neq 0,
\]
and $\A$ is said to be \textit{uniformly elliptic} in $\G$ (in $\clG$),
if it is elliptic in \textit{each} point $x \xeof \G$ (in $x \xeof \clG$).
We then call a boundary value problem $\A \,u \eq f$ \textit{elliptic},
if $\A$ is so.

\item[(E-2)]
$\A$ is called \textit{properly elliptic} in $\clG$,
if it is elliptic and 
if for each fixed $x \xeof \bndG$ and all linear independent vectors
$\xi, \eta \xeof \Rd$ the main symbol $P_{\A}(x, \xi + \tau \eta)$ has,
{\wrt} the variable $\tau$,
exactly $m$ roots with positive imaginary part and
exactly $m$ roots with negative imaginary part.
It can be shown that
if $\A$ is \textit{essentially real} or $\G \xsubset \Rd$ with $d \geq 3$,
then any elliptic differential operator is also {properly elliptic}.

\item[(E-3)]
The concept of ellipticity is often not good enough
to obtain satisfying results.
Hence the expression $\A$ of order $2m$
is called \textit{strongly elliptic} in $x \xeof \clG$,
if there is $\theta_x> 0$ such that
\[
\re \, (-1)^m \, p(x,\xi)
\, \geq \,
\theta_x \cdot \abs{\xi} ^{2m}
\quad
\fa \xi \neq 0.
\]

\item[(E-4)]
$\A$ is called \textit{uniformly strongly elliptic} in $\clG$,
if there is $\theta {\,>\,} 0$ such that 
\[
\re \, (-1)^m \, p(x,\xi)
\, \geq \,
\theta \cdot \abs{\xi} ^{2m}
\quad
\fa \xi \neq 0 \and x \xeof \clG.
\]
\end{itemize}

\subparagraph{} %- Laplace operator
The most well-known linear partial differential operator of second order
is the \textit{Laplacian}
\[
\Delta = \partial^2 / \partial {x_1}^2 + \ldots
       + \partial^2 / \partial {x_d}^2,
\]
whose negative version $-\Delta$ constitutes the prototype
of what is meant by an {uniformly strongly elliptic operator}.
One easily computes
$\re -p_{-\Delta}(x,\xi) \eq \abs{\xi} ^2$ for $\xi \xeof \Rd$,
hence the expression $-\Delta$ often serves as a first example
when studying linear-elliptic operators.

\paragraph{} %- parabolic and hyperbolic differential operators
Besides the class of {linear-elliptic} operators
there are two further important classes
of linear partial differential operators.
They are called to be of \textit{parabolic} and \textit{hyperbolic} type.
In order to introduce these concepts,
we consider differential expressions on cylindrical subsets
$\J \xtimes \G$,
where $\J \eq (0,\T) \xsubset \R$ and $\G \xsubset \Rd$ is a domain.
Then a \textit{parabolic} differential operator $\A$ is one of the form
\[
\A(t,x,u) \= \frac{\partial u}{\partial t}(x) \+ \A_1(x,u),
\]
where $\A_1$ is \textit{uniformly elliptic} in $\G$,
and $\A$ is of \textit{hyperbolic} type,
if it is of the form
\[
\A(t,x,u) \= \frac{\partial^2 u}{\partial t^2}(x) \+ \A_1(x,u),
\]
where again $\A_1$ is uniformly elliptic in $\G$.
Simple examples of para\-bolic and hyper\-bolic operators are
the \textit{heat operator} $\partial_t    - \Delta$ and
the \textit{wave operator} $\partial_{tt} - \Delta$,
respectively.
With these operators we associate the linear evolution equations
\[
\frac{\partial u}{\partial t} + \A(\cdot,u) = f
\quad \mbox{ and } \quad
\frac{\partial^2 u}{\partial t^2} + \A(\cdot,u) = f,
\]
both of them called \textit{initial value problems},
when considered with the prescribed start value $u(0) \eq u_0$
(and in addition $u'(0) \eq u_1$ in the hyperbolic case).

\subparagraph{} %- initial value problems
When considering parabolic or hyperbolic equations,
we speak of so-called \textit{initial-boundary value problems},
if for every time $t \xeof (0,T]$ the solution fulfills
a certain set of boundary conditions in every $x \xeof \bndG$,
where we especially have in mind the homogeneous Dirichlet boundary conditions
\eqref{HOMOGENEOUS_DIRICHLET_SYSTEM}.
However,
if $\G \eq \Rd$,
we obtain free space problems or pure initial value problems
for the differential operators of parabolic or hyperbolic type.

\paragraph{} %- overview on the section
%------------------------------------------------------------------------------%
It is the aim of this section to apply
all the presented functional analytic tools to the just enumerated problems.
Let $\A$ be a linear partial differential operator of order $k \in \N$
(not necessarily even) and given by
\begin{equation}\label{ANY_ORDER_LPDO}
\A(x,u) \eqdef
\sum_{\abs{\alpha} \leq k} c_{\alpha}(x) \, \PA u.
\end{equation}
Then by a \textit{realization} of $\A$ in a Banach space $\mathbb{E}(\G)$
of functions $u \: \G \xto \C$ is meant any linear operator
\[
\linA: \domain{\linA} \subset \mathbb{E}(G) \xto \mathbb{E}(G)
\]
with domain of definition $\domain{\linA}$.
The Banach space $\mathbb{E}(\G)$ should contain the space $\TF$ 
of test functions\footnote{
\,The space $\TF$ consists of indefinitely differentiable functions
having \textit{compact} support
\[
\supp{\,\phi} \, \eqdef \, \overline{\set{x \in \G}{\phi(x) \neq 0}}\,,
\]
usually denoted the \textit{space of test functions}.
This space plays an out\-standing role in analysis,
since it lies dense in a multiplicity of other functional spaces.}
as a dense linear subspace.
We also require for each $u \in \domain{\linA}$ that
first $\A \, u$ is well-defined by
$\linA u \eq \A \, u$ and second
$\A \, u \in \mathbb{E}(\G)$ holds.
Clearly,
our first condition means that the set $\domain{\linA}$
may contain only functions $u \: \G \xto \C$ being sufficiently often
differentiable in some sense.
But $\domain{\linA}$ should also include the space $\TF$ of test functions,
so $\domain{\linA}$ lies dense in $\mathbb{E}(\G)$ as well 
and $\linA$ appears as a densely defined operator in $\mathbb{E}$.

\paragraph{} %- content of the section
In the second part of this section we shall point out that

\begin{itemize}[leftmargin=2mm,itemsep=4pt]
\item[-] 
for a linear partial differential expression $\A$
there can be defined a couple of densely defined and closable linear operators
$\linA$ acting either in the {\HOELDER} spaces $C^{k,\gamma}(\G)$
($k \in \N_0$, $0 < \gamma \leq 1$)
or in the Lebesgue spaces $\Lp$ ($1 \leq p < \infty$);

\item[-] 
using elliptic estimates,
the realization of a strongly elliptic differentiable operator
acting in Sobolev spaces
under homogeneous Dirichlet boundary conditions is closed;

\item[-] 
the study on realizations of differential operators can be supported
by that on their adjoints,
where those are determined by so-called \textit{adjointness theorems};

\item[-] 
by starting with quasi-coercive sesquilinear forms,
created by strongly elliptic differential expressions,
we can obtain realizations in $\Lh$,
which prove to be densely defined and closed;

\item[-] 
weak solutions of initial-boundary value problems exist and are unique,
supposed the governing sesquilinear form is the Dirichlet form
of a strongly elliptic differential operator;

\item[-] 
the realizations of a regularly elliptic expression $\A$ in $\Lh$ are Fredholm,
so the alternative theorem \eqref{FREDHOLM_ALTERNATIVE_THEOREM}
may be applied to the equation $\A \, u \eq f$;

\item[-] 
strongly elliptic expressions give rise to m-accretive and 
sectorial realizations in $\Lh$,
so the related parabolic and hyperbolic initial value problem
prove to be well-posed;

\item[-] 
formally self-adjoint and half-bounded differential expressions
lead to self-adjoint realizations in $\Lh$;
\end{itemize}

\subparagraph{}
This enumeration is of course only a small selection 
within the totality of applications of functional analytic methods
to linear partial differential operators and related equations,
but may give a good impression of their capabilities.

\def\TITLEA{Realizations of differential operators in $C^{\gamma}(\G)$ and $\Lp$.}
\def\TITLEB{The closeness of strongly elliptic Dirichlet operators.}
\def\TITLEC{Formally adjoint expressions and adjointness theorems.}
\def\TITLED{{\GARDING}'s inequality and quasi-coercive Dirichlet forms.}
\def\TITLEE{Weak solutions of boundary and initial-boundary value problems.}
\def\TITLEF{The Fredholm property of elliptic operators.}
\def\TITLEG{Spectral theory and sectorialness of strongly elliptic operators.}
\def\TITLEH{Self-adjoint extensions of half-bounded operators.}

%-- realizations of differential equations in C^{k,\gamma}(G) and L^2(G)
%------------------------------------------------------------------------------%
\subsection*{I. {\TITLEA}}
We basically distinguish between two classes of functional spaces,
in which realizations of a linear differential operator $\A$ are acting:

\begin{itemize}[leftmargin=4mm]
\item[1.] %- realizations in Hölder spaces
The notation of classical derivative $\PA u$,
applied to a function $u \: \G \xto \C$,
leads us to the usage
of the spaces $C^{\gamma}(\G)$ of \textit{{\HOELDER}-constinuous} functions,
which constitute Banach spaces when endowed with a suitable norm.

\subparagraph{} %- Hölder continuity
Preliminary,
a function $u \: \G \xto \F$ is called \textit{{\HOELDER}-continuous}
{\wrt} the parameter $\gamma \in (0,1]$,
if its \textit{{\HOELDER}-norm} is finite,
that is,
\[
\abs{u} _{\,\gamma} \eqdef
\sup_{x,y \, \in \, \G, \, x \neq y} \,
\frac{\abs{u(x) - u(y)} } {\abs{x-y} ^{\gamma}} \, < \, \infty.
\]

\subparagraph{} %- Hölder spaces
Let $C^{k,\gamma}(\clG)$ be the linear space
of functions $u \in C^k(\clG)$,
whose partial derivatives $\PA u$
are {\HOELDER}-continuous for all $\abs{\alpha} \eq k$.
Then the \textit{{\HOELDER} norm} $\abs{u} _{k,\gamma}$
of a function $u \in C^{k,\gamma}(\clG)$ is defined
by the sum of the norm
$\norm{u}_k := \sum_{\abs{\alpha} \leq k} \,
\sup_{x \in \clG } \; \abs{\partial^{\alpha} u(x)} $
and the {\HOELDER} norms of the partial derivatives $\PA u$
for $\abs{\alpha} \eq k$:
\[
\abs{u} _{k,\gamma}
\eqdef
\norm{u}_k \+ \sum_{\abs{\alpha} \eq k} \abs{\,\PA u} _\gamma.
\]
The spaces $C^{k,\gamma}(\clG)$ prove to be complete {\wrt} this norm,
but they are not reflexive.

\subparagraph{}
A realization of a differential operator $\A$ of order $2m$
acting in {\HOELDER} space $C^{\gamma}(\G)$ is nothing but a mapping
\[ 
\linA \: \domain{\linA} \subset C^{\gamma}(\G) \to C^{\gamma}(\G),
\]
where $\domain{\linA}$ forms a linear subspace of $C^{2m,\gamma}(\G)$
restricted by boundary conditions,
for example by a Dirichlet system.

\subparagraph{}
There is an extensive theory on the realizations
of linear partial differential operators acting in {\HOELDER} spaces,
which provide nice results especially for operators of second order.
We should mention \textsc{J.\,Schauder} as the outstanding pioneer
in the investigation of this subject.
But \textsc{Schauder} uses classical methods for the setup of his theory,
coming from potential theory and complex analysis,
so we do not want to go into the details
(a good reference for this topic is \cite{GT}).

\item[2.] % realizations in Lebesgue spaces
Instead of classical derivatives,
the modern analysis of linear partial differential operators
is based on the concept of \textit{weak derivative},
since this approach leads to larger spaces of differentiable functions
than the already mentioned {\HOELDER} spaces,
usually denoted as \textit{Sobolev spaces}.

\subparagraph{} %- Hilbert space L^2
First recall that
for any domain $\G \xsubset \Rd$
the Hilbert space $\Lh$ is a complete normed space
consisting of measurable functions $f \: \G \xto \C$ such that
\[
\int_{\,\G} \, \abs{f} ^2 \, dx \, < \, \infty,
\]
where the inner product in $\Lh$ is defined by
\[
\scalar{f}{g} \eqdef \int_{\,\G} \, f {\cdot} \overline{g} \, dx
\quad \for f,g \in \Lh.
\]
This space is a specific member in the scale of the so-called
\textit{Lebesgue spaces} $\Lp$ ($1 \leq p < \infty$),
whose norms are defined by
\[
\abs{f} _p \eqdef ( \int_{\,\G} \, \abs{f} ^p \, dx)^{\frac{1}{p}} \,<\,\infty
\]
for equivalence classes of measurable functions $f \in \Lp$.

\subparagraph{}
The space $L^{\infty}(\G)$,
however,
consists of those equivalence classes of 
measurable functions $f \: \G \xto \C$,
for which 
\[
\norm{f}_{\infty} \, \eqdef \, \inf
\{\,a \geq 0: \abs{f} \leq a \,\, \mu \mbox{-almost everywhere}\,\}
\,<\,\infty.
\]

\paragraph{} %- weak derivatives
We next turn to the notation of \textit{weak derivative}.
Preliminary,
any \textit{locally integrable function}\footnote{
\,A function $u\:\G \xto \C$ is said to be \textit{locally integrable}
and to belong to the space $\Lc$,
if $u$ is Lebesgue-integrable on each compact subset of $\G$.}
$u:\G \xto \C$ is said to be \textit{weakly differentiable} 
{\wrt} a given multi-index $\alpha$,
if there is a further locally integrable function $v$ such that
\[
\int_{\,\G} v(x) \, \phi(x) \, dx
=
(-1)^{ \abs{\alpha} } \,
\int_{\,\G} u(x) \, \PA \phi(x) \, dx
\fa \phi \in \TF.
\]
We denote $\PA u := v$ as the $\alpha$-\textit{weak derivative} of $u$
in this case.
The function $u$ is said to be \textit{$k$-times weakly differentiable},
if $u$ has $\alpha$-weak derivatives
for every index $\alpha$ with $\abs{\alpha} \,{\leq}\, k$.
Taking weak derivatives of locally integrable functions {\wrt} the multi-index
$\alpha$ arises as a linear mapping
\[
\PA \: D_{\alpha} \subset \Lc \to \Lc,
\]
where $D_{\alpha}$ forms a linear subspace of $\Lc$.
Let us also notice that
the concept of weak derivative is compliant
with that of the classical derivative:
if a function $u$ is classical differentiable in $\G$
with classical derivative $\PA u$,
then $u$ is also weakly differentiable,
and its weak derivative is equal to $\PA u$.

\subparagraph{} %- Lebesgue-Sobolev spaces
The space $W^{k,p}(\G)$ ($1 \leq p < \infty)$
is the \textit{Lebesgue-Sobolev space}
\[
W^{k,p}(\G) \eqdef
\{\,f \in \Lp: \PA f \in \Lp, \; \abs{\alpha} \leq k \, \}.
\]
If $p \eq 2$,
the space $H^k(\G) := W^{k,2}(\G)$ can be furnished with inner product
\[
\scalar{f}{g}_k \eqdef
\sum_{\abs{\alpha} \leq k} \int_{\,\G} \;
\PA f \cdot \PA \overline{g} \, dx\,,
\quad \quad f,g \xeof H^k(\G)
\]
and norm
$\abs{f} _k \eq \scalar{f}{f}^{1/2}$,
so $H^k(\G)$ forms a Hilbert space in its own right.
In other words,
the space $H^k(\G)$ consists of square-integrable functions,
which are at least $k$-times {weakly differentiable}
and all these weak derivatives belong to $\Lh$ again.

\subparagraph{} %- Hilbert-Sobolev spaces with homogeneous boundary conditions
We also want to consider the linear subspace $\WN{k,p}{\G}$ of $W^{k,p}(\G)$
fulfilling homogeneous Dirichlet boundary conditions.
This space consists of those functions $u \xeof W^{k,p}(\G)$,
which are approximated by a sequence $\{\phi_n\}_{n \xeof \N}$
of \textit{test functions} {\wrt} the norm $\abs{\,\cdot\,} _p$,
that is,
each $\phi_n$ belongs to the space $\TF$ and $\abs{u - \phi_n} _p \xto 0$. 
Clearly,
since $\WN{k,p}{\G}$ is a closed subspace of $W^{k,p}(\G)$,
it forms a Banach space again.
Furthermore,
functions belonging to $\WN{k,p}{\G}$ fulfill the $k$ boundary conditions
of a homogeneous Dirichlet system.
For $p \eq 2$,
however,
the space $\HN{k}{\G}$ owns the inner product and norm
as defined in the larger space $H^k(\G)$,
so it is a Hilbert space as well.

\subparagraph{}
Finally,
a realization of a linear differential operator $\A$ of order $k$
in Lebesgue space $\Lp$ is any linear mapping
\[
\linA \: \domain{\linA} \subset \Lp \to \Lp,
\]
for which $\domain{\linA}$ forms a linear subspace of $W^{k,p}(\G)$,
for example restricted by homogeneous Dirichlet boundary conditions.
\end{itemize}

%-- the closeness of strongly elliptic Dirichlet operators
%------------------------------------------------------------------------------%
\subsection*{II. \TITLEB}
One often has to deal with the realization $\linA_D$
of a strongly elliptic differential operator $\A$ in $\Lp$
having order $2m$ and domain of definition
\[
\domain{\linA_D} \eqdef \{ u \in \WN{m,p}{\G}: \A\,u \in \Lp \},
\]
so the functions in $\domain{\linA_D}$ satisfy all the $m$ boundary conditions
of the homogeneous Dirichlet system \eqref{HOMOGENEOUS_DIRICHLET_SYSTEM}.
By this reason it is reasonable to denote the realization $\linA_D$
a \textit{Dirichlet operator}.

\paragraph{} %- closeness of strongly elliptic differential operators
We turn to the question of closeness of the realization $\linA_D$.
Let the domain $\G$ and the expression $\A$ fulfill the following conditions:

\begin{itemize}[leftmargin=9mm]
\item[(I)\;\;]
$\G$ is bounded in $\Rd$,

\item[(I')\;\,]
$\G$ satisfies the $m$-th \textit{extension property},
that is,
the linear space $R_0(\G)$ consisting
of those functions $\phi \xeof W^{m,p}(\G)$,
which are restrictions $\phi |_{\G}$
of functions $\phi \xeof C_c^{\infty}(\Rd)$,
is dense in $W^{m,p}(\G)$;

\item[(II)\;]
$\A$ is uniformly strongly elliptic in $\G$;

\item[(III)\,]
the coefficient functions $a_{\alpha \beta}$ are essentially bounded in $\G$,
that is,
$a_{\alpha \beta} \xeof L^{\infty}(\G)$;

\item[(III')]
the coefficient functions $a_{\alpha \beta}$ of highest order
$\abs{\alpha} + \abs{\beta} \eq 2m$ are uniformly continuous in the closure
$\clG$ of $\G$.
\end{itemize}

Then,
by showing the validness of the \textit{elliptic estimate}
\[
\abs{u} _{2m} \, \leq \, c \cdot \{ \,\abs{u} _p + \abs{\linA_D u} _p \, \}
\]
for $u \xeof \domain{\linA_D}$,
where $\abs{u} _p \eq \norm{u}$ is the norm of $u$ in $\Lp$,
theorem \ref{BROWDER_THEOREM} implies that the operator $\linA_D$ is closed.
Moreover,
$\linA_D$ is injective and the range of $\linA_D$ is a closed subspace in $\Lp$.

%-- formally adjoint expressions and adjointness theorems
%------------------------------------------------------------------------------%
\subsection*{III. {\TITLEC}}
As already mentioned,
it is often helpful to combine the study of a linear operator
with that of its formally adjoint.
Supposed the linear partial differential operator $\A$ of order $k$
has coefficients $c_{\alpha} \xeof C^{\abs{\alpha} }(\G)$.
We then find for $u \xeof \TF$ and $\psi \xeof \TF$
by repeated partial integration
\[
\int_{\,\G} \sum_{\abs{\alpha} \leq k}
(c_{\alpha} \, \PA u) \cdot \overline{\psi} \, dx
\=
\int_{\,\G} u \cdot \sum_{\abs{\alpha} \leq k}
(-1)^{\abs{\alpha} } \, \overline{\PA(\overline{c_{\alpha}} \, \psi)} \, dx.
\]
Hence there is a \textit{formally adjoint} differential expression
\begin{equation}\label{FAO}
\A^{\dag}(x,u) \eqdef
\sum_{\abs{\alpha} \leq k} (-1)^{\abs{\alpha} }
\PA(\overline{c_{\alpha}(x)} \, u(x))
\end{equation}
such that the identity
\[
\int_{\,\G} \A \, \phi \, \overline{\psi} \, dx 
\=
\int_{\,\G} \phi \, \overline{\A^{\dag} \, \psi} \, dx 
\]
holds for $\phi,\psi \xeof \TF$,
briefly
$
\scalar{\A\,\phi}{\psi} \eq \scalar{\phi}{\A^{\dag}\,\psi}
$
by means of the inner product in $\TF$.
Using \textsc{Leibniz}'s rule to compute the terms
$\PA(\overline{c_{\alpha}} \, u)$ in \eqref{FAO},
there are functions $c_{\alpha}'\: \G \xto \C$
such that
\[
\A^{\dag}(x,u) \=
\sum_{\abs{\alpha} \leq k} c_{\alpha}'(x) \; \PA u(x).
\]
Thus $\A^{\dag}$ is uniquely determined by $\A$
and constitutes a differential expression of type \eqref{ANY_ORDER_LPDO},
too.
It may happen that $\A^{\dag} \eq \A$ holds,
and in this case
we shall call the expression $\A$ \textit{formally self-adjoint}.

\subparagraph{} %- a criterion for closability of differential operators
A useful criterion to ensure closability of a linear operator $\linA$
acting in Hilbert space $\Lh$
and induced by a differential expression $\A$ of order $k$
is given by the following theorem
(see also \cite{Yo}, p.\,78),
whose proof uses the presence of the formally adjoint expression $\A^{\dag}$:

\begin{theorem1}\label{L2_CLOSABILITY_THEOREM}
Let the expression $\A$ of order $n$ be given by \eqref{ANY_ORDER_LPDO}
with coefficients $c_{\alpha} \xeof C^k(\G)$ for $\abs{\alpha} \leq k$ and
the linear operator
\[
\linA: \;
\begin{cases}[1.2]
\quad \domain{\linA} \eqdef \{u \xeof C^k(\G) \cap \Lh:\A \, u \in \Lh \},\\
\quad \linA u \eqdef \A\,u \; \for u \xeof \domain{\linA}\\
\end{cases}
\]
be a realization of $\A$ in $\Lh$.
Then $\linA$ is closable.
\end{theorem1}

\subparagraph{} %- description of the domain of definition of the closure
Note that the domain of definition $\domain{\linA}$
of the operator $\linA$ considered in theorem \ref{L2_CLOSABILITY_THEOREM}
is not restricted to any boundary conditions.
It should also be mentioned that
rather the closability of the differential operator $\linA$
with domain of definition $\domain{\linA}$
like in theorem \ref{L2_CLOSABILITY_THEOREM} is ensured,
it is not an easy task to give an exact description
for the domain of definition of its closure.
But if the coefficients of $\A$ are constant,
then the set $\domain{\overline{\linA} }$ can be determined
and coincides with the Hilbert-Sobolev space $H^k(\G)$.

\subparagraph{} %- minimal and maximal operators
The closure $\linA_{min} := \overline{\linA_0}$
of a closable linear differential operator $\linA_0$
with domain of definition $\domain{\linA_0} \eq \TF$
is denoted the \textit{minimal realization}
or the \textit{minimal operator} of $\A$.
The \textit{maximal operator} $\linA_{max}$ of $\A$,
however,
is by definition the adjoint of the minimal operator of $\A^{\dag}$
with domain of definition $\domain{{\linA^{\dag}}_0} \eq \TF$,
that is,
$\linA_{max} := ({\linA^{\dag}}_{min})^*$.
One always has $\linA_{min} \, {\prec} \, \linA_{max}$,
and often any operator $\linA$ fulfilling
\[
\linA_{min} \prec \linA \prec \linA_{max}
\]
is called a \textit{realization} of $\A$,
in opposite to our selected definition of this concept.
Using this definition,
we also want to consider the linear operators $\linA_{min}$ and $\linA_{max}$
as realizations of $\A$,
too.

\paragraph{} %- adjointness theorems
To the end,
the question arises whether any realization
of the formally adjoint expression $\A^{\dag}$ is in relationship
with the adjoint of some realization of $\A$.
This problem is examined in \cite{Br2} (section\,3: Adjoints)
even in the context of $L^p$-spaces.
\textsc{Browder}'s result for Hilbert spaces $\Lh$
and the Dirichlet realization $\linA_D$
of a regularly elliptic differential expression $\A$ sounds as follows
(corresponding to theorem 5, see p.\,57):

\begin{theorem1}
Let the $2m$-order expression $\A$ be strongly elliptic
with sufficiently smooth coefficients
and defined on a smooth domain $\G \xsubset \Rd$
such that the formal adjoint $\A^{\dag}$ exists and its coefficients
are uniformly bounded in $\G$.

\subparagraph{} %- adjointness theorem in L^2
If $\linA$ and $\linA^{\dag}$ are the realizations of $\A$ and $\A^{\dag}$
in $\Lh$ with homo\-geneous Dirichlet boundary conditions
and domains of definition
\[
\domain{\linA} \= \domain{{\linA}^{\dag}} \eqdef H^{2m}(\G) \cap \HN{m}{\G},
\]
then $\linA^{\dag}$ is the Hilbert adjoint of ${\linA}$,
that is,
${\linA}^* \eq \linA^{\dag}$.
\end{theorem1}

\subparagraph{} %- adjointness theorem in L^p
A similar result holds in the setting of the spaces $\Lp$ for $1 < p < \infty$.
As it is quoted in \cite{Pa}, lemma 7.3.4,
we have 
${\linA_p}^* \eq {{\linA}^{\dag}}_q$,
where $q \eq p/(p-1)$ respectively $\frac{1}{p} + \frac{1}{q} \eq 1$
and both operators
\[
\begin{cases}[1.2]
\quad \domain{\linA_p} \eqdef W^{2m,p}(\G) \cap \WN{m,p}{\G},\\
\quad \linA_p u \eqdef \A\,u \for u \in \domain{\linA_p},\\
\end{cases}
\]
\[
\begin{cases}[1.2]
\quad \domain{{\linA^{\dag}}_q} \eqdef W^{2m,q}(\G) \cap \WN{m,q}{\G},\\
\quad {\linA^{\dag}}_q u \eqdef \A^{\dag} u
      \for u \in \domain{{\linA^{\dag}}_q}.\\
\end{cases}
\]
are densely defined and closed.

%-- Gardings's inequality and quasi-coercive Dirichlet forms
%------------------------------------------------------------------------------%
\subsection*{IV. {\TITLED}}
We next want to illustrate the usage of abstract sesquilinear
forms for the analysis of linear equations of the form $\A \, u \eq f$,
where $\A$ is a linear partial differential operator of order $2m$
and $u,f$ are real or complex-valued mappings
defined on a domain $\G \xsubset \Rd$.

\subparagraph{} %- differential expressions in divergence form
Preliminary,
a linear differential expression $\A$ of order $2m$
is said to be in \textit{divergence form},
if there are functions $a_{\alpha \beta} \xeof C^{\abs{\alpha} }(\G)$
for $\abs{\alpha} , \abs{\beta} \leq m$
such that $\A$ can be written as
\begin{equation}\label{LPDO_DIV}
\A \eqdef
\sum_{\abs{\alpha} \leq m, \abs{\beta} \leq m} (-1)^{\abs{\alpha} } \,
\PA(\,a_{\alpha \beta}(x) \, D^{\beta}).
\end{equation}

In opposite to this definition let us call an arbitary differential expression
given by \eqref{ANY_ORDER_LPDO} to be \textit{in non-divergence form}.

\subparagraph{} %- relationship between divergence and non-divergence forms
It is a matter of fact
that every expression $\A$ in divergence form
can be transformed into an expression $\A'$ in \textit{non-divergence form}.
Reversely,
a non-divergence form expression $\A'$ given by \eqref{ANY_ORDER_LPDO}
can be transformed into an expression $\A$ in divergence form \eqref{LPDO_DIV},
if $c_{\alpha} \xeof C^{2m - \abs{\alpha} }(\G)$ for $\abs{\alpha} \leq 2m$,
that is,
if the coefficients are sufficiently many differentiable in $\G$.
Comparing the expressions \eqref{ANY_ORDER_LPDO} and \eqref{LPDO_DIV},
it turns out that
$\A'$ is strongly elliptic in $\G$ \iff $\A$ is so.

\subparagraph{}
To show the purpose of a linear differential operator being in divergence form,
let the space $\TF$ of {test functions} be furnished with the inner product
\[
\scalar{\phi}{\psi} \eqdef
\int_{\,\G} \phi \, \overline{\psi} \,dx \qfa \phi,\psi \xeof \TF.
\]
Now,
applying $\A$ given in divergence form to $\phi \xeof \TF$,
then multiplying $\A \, \phi$ with $\overline{\psi} \xeof \TF$ and
finally performing $m$-times partial integration leads to
\[
\DF_{\A}(\phi,\psi) \eqdef \scalar{\A \, \phi}{\psi} \=
\sum_{\abs{\alpha} \leq m, \abs{\beta} \leq m} \,
\int_{\,\G}
a_{\alpha \beta} \; \PA \phi \; D^{\beta} \overline{\psi}
\, dx,
\]
since all the boundary terms vanish.
The mapping
\[
\DF_{\A} \eqdef \TF \xtimes \TF \xto \C
\]
is a sesquilinear form,
called the \textit{Dirichlet form} of $\A$ in $\G$.
If the coefficent functions $a_{\alpha \beta}$ of $\A$
are measurable and essentially bounded,
then the mapping $\DF_{\A}$ is continuous
and thus can continuously be extended to the product space
$\HN{m}{\G} \xtimes \HN{m}{\G}$.

\paragraph{} %- Garding's inequality
We are mainly interested in a $2m$-order linear differential operator $\A$
given in divergence form,
which together with the underlying domain $\G$ fulfills the already 
mentioned conditions (I), (I'), (II), (III) and (III').
\textsc{{\GARDING}}\textit{'s inequality} can be seen as a bridge
from the theory of linear-elliptic differential operators
to abstract functional analysis:

\begin{theorem1}\label{GARDING_INEQUALITY_THEOREM}
Let the domain $\G$ fulfill (I),\,(I') 
and let the linear differential expression $\A$ of order $2m$ fulfill
(II),\,\,(III),\,(III').\\
Then there is $\lambda_0 \xeof \R$ and $a_0 {\,>\,} 0$ such that
for all $u \xeof \HN{m}{\G}$
\begin{equation}\label{GARDING_INEQUALITY}
\DF_{\A}(u,u) + \lambda_0 \norm{u}^2 =
\scalar{\,(\A+\lambda_0)\,u}{u} \geq a_0 \, {\abs{u} _m}^2.
\end{equation}
\end{theorem1}

\subparagraph{}
In other words,
theorem \ref{GARDING_INEQUALITY_THEOREM} quotes that
the Dirichlet form associated with the operator $\A + \lambda_0$ is coercive,
hence the \textsc{Lax-Milgram} theorem may be applied.
It should be clear that this assertion also holds
for all $\lambda {\,>\,} \lambda_0$.

\paragraph{} %- Garding's inequality for second order operators
The most important class of linear partial differential operators
in divergence form {\wrt} applications is that of \textit{second order},
that is,
our operator $\A$ is of the form
\begin{equation}\label{GENERAL_SECOND_ORDER_DIFFERENTIAL_OPERATOR}
\A(x,\cdot) = 
-\sum_{i=1}^d \,
\partial_i \, (\sum_{j=1}^d a_{ij}(x)) 
+
\sum_{j=1}^d \, (b_j(x) \, \partial_j)
+
c(x),
\end{equation}
which is nothing but \eqref{LPDO_DIV} for $m \eq 1$.
The Dirichlet form $\DF_{\A}$ belonging to $\A$ is the sesquilinear form
\[
\DF_{\A}(u,v) \= \int_{\,\G} 
(\sum_{i,j=1}^d a_{ij} \, \partial_i u \, \partial_j \overline{v} \+ 
 \sum_{j=1}^d b_j \, \partial_j u \, \overline{v} \+
 c \,u \, \overline{v}
)
\, dx.
\]

\subparagraph{} %- coerciveness of the Dirichlet form of second-order operators
A short computation shows
that if the function $c$ is large enough in comparison
with the sum of partial derivatives of the functions $b_j$,
then $\DF_{\A}$ is coercive.
More precisely,
$\DF_{\A}$ proves to be coercive if
\begin{equation}\label{COERCIVITY_CONDITION}
c(x) \, \geq \, \frac{1}{2} \,
\sum_{k=1}^d \, \frac{\partial b_k}{\partial x_k}(x)
\end{equation}
holds in all points $x \xeof \G$.

\subparagraph{}
On the other hand,
if the matrix $A(x) \eq (a_{ij}(x))_{i,j=1}^d$ of second-order coefficients
in \eqref{GENERAL_SECOND_ORDER_DIFFERENTIAL_OPERATOR}
fails to be positive definite for all $x \xeof \G'$,
where $\G' \xsubset \G$ has positive measure,
then $\DF_{\A}$ cannot be coercive on $\HN{1}{\G}$.
Also,
if $A(x)$ is positive definite in all $x \xeof \G$
and condition \eqref{COERCIVITY_CONDITION} does not hold.
then the functions $b_j$ disturb the coercivity of $\DF_{\A}$.
But after all,
the form $\DF_{\A}$ proves at least to be quasi-coercive,
which is nothing but the assertion of \textsc{\GARDING}'s inequality.

\subparagraph{} %- Garding's inequality for expressions of second order
\textsc{\GARDING}'s inequality remains even true
for the Dirichlet forms of second-order differential operators,
if the basic conditions (I)' and (III') are dropped:

\begin{theorem1}\label{GARDING_INEQUALITY_THEOREM_II}
Let $\G \xsubset \Rd$ fulfill (I) and
the differential operator $\A$ of second order and given in divergence form
fulfill (II) and (III).
Then there is $\lambda_0 \xeof \R$ and $a_0 {\,>\,} 0$ such that
\[
\DF_{\A}(u,u) \+ \lambda_0 \norm{u}^2
\=
\scalar{\,(\A {+} \lambda_0)\,u}{u}
\, \geq \,
a_0 \, {\abs{u} _1}^2
\]
for all $u \xeof \HN{1}{\G}$,
where $a_0$ is the uniform lower bound of the eigenvalues
of the matrices $A(x)$ with $x \xeof \G$.
\end{theorem1}

%-- weak solutions of boundary and initial-boundary value problems
%------------------------------------------------------------------------------%
\subsection*{V. {\TITLEE}}
The variational approach to stationary and evolutional boundary value problems
provides the concept of weak solution,
and in this part of the section we want to apply the abstract theory
boundary value problems governed by strictly elliptic differential operators
by quoting the related existence and uniqueness assertions.

\subparagraph{} %- the variational Dirichlet problem
We basically assume that

\begin{itemize}[leftmargin=15mm]
\item[(DP-1)]
$\A$ is a the differential operator of order $2m$,
strongly elliptic in a \textit{bounded} domain $\G \xsubset \Rd$
and given in divergence form;

\item[(DP-2)]
the coefficients $a_{\alpha \beta}$ belong to $L^{\infty}(\G)$
and the top-order coefficients are uniformly continuous in $\clG$.
\end{itemize}

\paragraph{} %- homogeneous Dirichlet problem
It is our first aim to solve the \textit{homogeneous Dirichlet problem}
\[
\begin{cases}[1.2]
\; (\A\,u)(x)\= f(x)& (x \xeof \G),\\
\; u(x) =
\partial u/ \partial \nu (x) =
\ldots = \partial u^{m{-}1} / \partial \nu^{m{-}1}(x) = 0 & (x \xeof \bndG).\\
\end{cases}
\]
Let $\vecV \eq \HN{m}{\G}$ and $\hilH \eq \Lh$,
hence the pair $(\vecV,\hilH)$ is a Hilbert triple.
By \textsc{{\GARDING}}'s inequality,
the Dirichlet form
\[
\DF_{\A}(u,v) \eqdef \scalar{\A u}{v} \quad \quad (u,v \in \vecV)
\]
is quasi-coercive with some $\lambda_0 \xeof \R$.
Then,
by means of the \textsc{Lax-Milgram} theorem,
we obtain the following \textit{well-posedness theorem}:

\begin{theorem1}\label{DIRICHLET_EXISTENCE_UNIQUENESS_THEOREM}
For $\lambda \geq \lambda_0$, the variational equation
\[
\DF_{\A}(u,v) \+ \lambda \scalar{u}{v}
\= f(v) \quad (v \xeof \HN{m}{\G})
\]
owns for $f \xeof H^{-m}(\G) := \HN{m}{\G}^*$
a unique solution $u \xeof \HN{m}{\G}$ depending
continuously on the right-hand side $f$ and
called the \textit{weak solution} of the homogeneous Dirichlet problem.
\end{theorem1}

\paragraph{} %- parabolic variational initial value problems
If a linear differential operator $\A$ is given in divergence form
and its Dirichlet form $\DF_{\A}$ is coercive,
then the variational approach may also be applied
to the \textit{heat equation} $u_t + \A \, u \eq f$.

\subparagraph{}
Consider for $f \xeof L^2([0,\T), H^{-1}(\G))$ and $u_0 \xeof \Lh$
the \textit{parabolic} initial value problem
\[
\begin{cases}[1.2]
\quad \scalar{u_t}{v} \+ \DF_{\A}[u,v] \= \scalar{f}{v} 
\qfa v \xeof \HN{1}{\G},\\
\quad u(t,\cdot) \xeof \HN{1}{\G}, \\
\quad u(0) = u_0.\\
\end{cases}
\]
Following theorem \ref{LIONS_THEOREM_I},
this problem owns a uniquely defined weak solution
in $H^1(0,\T; \HN{1}{\G}, H^{-1}(\G))$,
which fulfills the variational equation $\lambda$-nearly everywhere in $[0,\T)$.
Notice that this assertion is especially true in case of $\A \eq -\Delta$.

\paragraph{} %- variational approach to the hyperbolic problem
A similar result holds for the variational approach
to the \textit{wave equation} $u_{tt} + \A \, u \eq f$
governed by the differential operator $\A$
having the properties (DP-1) and (DP-2).
Then the \textit{hyperbolic} initial value problem in variational form sounds as
\[
\begin{cases}[1.2]
\quad \scalar{u_{tt}}{v} \+ \DF_{\A}[u,v] \= \scalar{f}{v} 
\qfa v \xeof \HN{1}{\G}, \\
\quad u(t,\cdot) \xeof \HN{1}{\G},\\
\quad u(0)=u_0, \; u_t(0)=u_1, \\
\end{cases}
\]
where the coefficients $a_{ij}$ of the operator $\A$ are assumed
to fulfill the symmetry conditions
\[
a_{ij}(x) \= \overline{a_{ji}(x)}
\quad \quad
(1 \leq i,j \leq d, \; x \xeof \G)
\]
and the Dirichlet form $\DF_{\A}$ is assumed to be coercive.

\subparagraph{}
If $f$ belongs to $L^2([0,\T), \Lh)$,
$u_0 \xeof \HN{1}{\G}$ and $u_1 \xeof \Lh$,
then theorem \ref{LIONS_THEOREM_II} quotes
that this problem has exactly one weak solution $u \xeof L^2(0,\T; \HN{1}{\G})$
with $u(t) \xeof \HN{1}{\G}$ for $t \xeof [0,\T)$ $\lambda$-nearly everywhere.

\subparagraph{}
Since the Dirichlet $\DF_{-\Delta}$ is symmetric,
the variational initial value problem for the classical wave equation
$u'' - \Delta u \eq f$
owns a uniquely defined weak solution
$u \xeof H^1(0,\T; \HN{1}{\G}, H^{-1}(\G))$,
if the functions $f, u_0, u_1$ fulfill the prerequisites
mentioned in theorem \ref{LIONS_THEOREM_II}.

%-- the Fredholm property of elliptic operators
%------------------------------------------------------------------------------%
\subsection*{VI. {\TITLEF}}
This statement is properly not correct,
since it only proves to be true for properly elliptic differential operators
defined in a \textit{bounded} domain $\G \xsubset \Rd$
and endowed with certain boundary conditions,
for example those given by a homogeneous Dirichlet system.

\subparagraph{} %- Rellich's embedding theorem
We first want to present \textsc{Rellich}\textit{'s embedding theorem},
which holds for Sobolev spaces whose underlying domain is \textit{bounded}.
It can be seen as a key theorem in the analysis
of linear-elliptic boundary problems on such domains:

\begin{theorem1}\label{RELLICH_EMBEDDING_THEOREM}
If $\G \xsubset \Rd$ is bounded and owns smooth boundary $\bndG$,
then the continuous embeddings
$\iota_k \: \HN{k}{\G} \hookrightarrow \Lh$ are compact for all $k \xeof \N$.
\end{theorem1}

%\subparagraph{}
%A proof of this theorem can be found for example in \ldots.

\subparagraph{} %- elliptic operators as continuous mappings
The basic approach to investigate the Fredholm property
of a properly elliptic operator $\A$ in divergence form
and of order $2m$
consists of two small changes for the Dirichlet realization $\linA_D$
having domain of definition $\domain{\linA} \eq H^{2m}(\G) \cap \HN{m}{\G}$:
let us consider the linear operator $\linA_D$ as

\begin{itemize}[leftmargin=9mm]
\item[(F-1)]
a mapping $\linA_D \: \HN{m}{\G} \xto \Lh$,
which is continuous due to the essential boundedness
of the coefficients $a_{\alpha,\beta}$ in \eqref{LPDO_DIV};

\item[(F-2)]
the sum $\linA_1 + \linA_2$ of two further linear operators,
where we shall assume that
$\linA_1 \: \HN{m}{\G} \xto \Lh$ is the main part 
having weak derivatives of order $2m$ only,
whereas $\linA_2 \: \HN{m}{\G} \xto \Lh$ is the part of $\linA_D$
consisting of weak derivatives with order less than $2m$.
\end{itemize}

\subparagraph{}
Let $\lambda_0$ be the smallest parameter
for the Dirichlet form $\DF_{\linA+\lambda_0}$ to be coercive.
Since
\[
\linA_D \= (\linA_1 + \lambda_0) \+ (\linA_2 - \lambda_0)
\]
and $\linA_2- \lambda_0$ may be seen as a compact perturbation
of $\linA_1+\lambda_0$,
which maps $H^m(\G)$ one-to-one onto $\Lh$ and consequently is Fredholm 
with vanishing index,
the operator $\linA_D$ is Fredholm again with index $\chi(\linA_D) \eq 0$.

%-- spectral theory and sectorialness of strongly elliptic operators
%------------------------------------------------------------------------------%
\subsection*{VII. {\TITLEG}}
Let us consider again a strongly elliptic differential operator $\A$
given in divergence form and its Dirichlet realization $\linA_D$
in a bounded domain $\G \xsubset \Rd$.
By {\GARDING}'s inequality,
the Dirichlet form $\DF_{\A}$ is quasi-coercive
{\wrt} the parameter $\lambda_0$,
hence $\linA_D + \lambda$ proves to be bijective
for all $\re \lambda {\,>\,} \lambda_0$.
Furthermore,
the homogeneous Dirichlet problem in $\Lh$
\begin{equation}\label{HOMOGENEOUS_DIRICHLET_PROBLEM}
\linA_D \, u  \+ \lambda u \= f \quad \quad (f \in \Lh)
\end{equation}
is well-posed,
which is a direct consequence of
theorem \ref{DIRICHLET_EXISTENCE_UNIQUENESS_THEOREM}.

\subparagraph{}
In the special case,
where $f \in \Lh$,
the solutions $u$ of \eqref{HOMOGENEOUS_DIRICHLET_PROBLEM}
belong to the Hilbert-Sobolev space $H^{2m}(\G)$ and are called \textit{strong}.
It should be clear that 
each strong solution of \eqref{HOMOGENEOUS_DIRICHLET_PROBLEM}
is also a weak solution.
But in order to show that these solutions $u$ are even classical,
that is,
$u \xeof C^{2m}(\G) \cap C^m(\clG)$,
methods of \textit{regularity theory} for elliptic problems have to be applied.

\paragraph{} %- spectrum of the Dirichlet operator
Let us turn to the spectrum of $\linA_D$.
Since $-\Delta_D + \lambda$ forms a bijective and closed operator
between the spaces $H^{2m}(\G) \cap \HN{m}{\G}$ and $\Lh$
for each $\lambda \in C$ with $\re \lambda \geq \lambda_0$,
we surely have
\[
\{ \zeta \in \C: \re \zeta \geq \lambda_0\} \xsubset \rho(\linA_D)
\and
\sigma(\linA_D) \xsubset \{ \zeta \in \C: \re \zeta < \lambda_0\}.
\]
Furthermore,
if $\G$ is bounded,
then $\linA_D$ is an operator with pure point spectrum.

\paragraph{} %- sectorialness of the Dirichlet operator
We can also make assertions on the accretivity and sectorialness
of the operator $\linA_D$.
As theorem \eqref{QUASI-COERCIVE_FORMS_AND_ASSOCIATED_OPERATORS}
(d),\,(e) shows,
$\linA_D + \lambda_0$ is the generator
of a semigroup of bounded linear operators
that is strongly continuous and analytic at the same time.
Hence the initial value problems for the heat and wave equation
governed by the operator $\linA_D + \lambda_0$ is well-posed.
By perturbation theory of semigroups,
this is also true for both problems governed by the operator $\linA_D$ itself.

\paragraph{} %- Gauss-Weierstrass semigroup
Let us finally consider the heat equation
for the negative Laplacian $-\Delta_D$ in the entire space $\Rd$,
which of course is a free space problem.
Theorem \ref{SEMIGROUP_CREATED_BY_SECTORIAL_OPERATOR_IS_ANALYTIC} ensures
that the realization $-\Delta_D$ in Hilbert space $\Lh$
with domain of definition
\[
\domain{-\Delta_D} \= H^2(\Rd)
\]
is the infinitesimal generator
of a sectorially contractive analytic semigroup $\{W(t)\}_{t \geq 0}$,
denoted the \textsc{Gauss-Weierstrass} \textit{semigroup}.
Hence there is a uniquely defined solution $u \xeof C^1(\J \xtimes \Rd,\R)$
for the homogeneous Cauchy problem
\[
\begin{cases}[1.2]
\quad u_t - \Delta u \= 0,\\
\quad u(x,0)         \= u_0(x) \for x \in \Rd \\
\end{cases}
\]
with initial mapping $u_0 \in H^2(\Rd)$.
Then one may represent the semigroup $\{W(t)\}_{t \geq 0}$
explicitly by means of the \textit{heat kernel} function $\Phi_d$,
that is,
for all $t > 0$ and $x \xeof \Rd$
\[
(W(t) \, u_0) \, (x)
=
\frac{1}{(4 \pi \, t)^{d/2}}
\int_{\,\Rd} e^{- \frac{\abs{x-y} ^2}{4t}} u_0(y) \, dy
=
(\Phi_d(t) \ast u_0)\,(x).
\]

%-- self-adjoint extensions of half-bounded operators
%------------------------------------------------------------------------------%
\subsection*{VIII. {\TITLEH}}
Remember the definition of the differential expression $\DO^{\dag}$.
We then denote a linear differential expression $\DO$
\textit{formally self-adjoint} in case of $\DO^{\dag} \eq \DO$.
If $\DO$ has constant coefficients,
then $\A$ proves to be formally self-adjoint
iff
$c_{\alpha} \xeof \R$ for $\abs{\alpha} $ even
and $c_{\alpha} \xeof i \, \R$ for $\abs{\alpha} $ odd.

\subparagraph{}
In the following we want to enumerate a couple of
formally self-adjoint differential expressions
and describe some of their self-adjoint realizations.

\begin{itemize}[leftmargin=9mm]
\item[(a)] \textit{Sturm-Liouville operators.}
A first example of a formally self-adjoint linear differential operator
(given in divergence form) is the \textsc{Sturm-Liouville} \textit{operator}
\[
\A_{sl}(u,x) \eqdef -(p(x)\, u'(x))' \+ q(x)\, u(x),
\quad \quad x \xeof (a,b)
\]
with $p(x) {\,>\,} 0$ and $q(x) \geq 0$ for $x \xeof [a,b]$
with $-\infty {\,<\,} a {\,<\,} b {\,<\,} \infty$.
This operator determines the elongations of a vibrating string
being fixed at the endpoints of the interval $[a,b]$.
The corresponding Dirichlet form reads as
\[
\DF_{\A_{sl}}(u,v) \=
\int_a^b \,
(p \, u'\, \overline{v}' \+ q \, u \, \overline{v}) \, dx,
\]
which shows that the expression $\A_{sl}$ is symmetric and half-bounded
on the domain of definition $\TF$.
The energetic space $H_{\A_{sl}}$ is the Hilbert-Sobolev space $\HN{1}{a,b}$
containing the restrictions of 
\textit{absolutely continuous functions} $u\: [a,b] \xto \R$
with boundary conditions $u(a) \eq u(b) \eq 0$ to the interval $(a,b)$.

\subparagraph{}
The self-adjoint extension $\linA_{sl,F}$
of the \textsc{Sturm-Liouville} operator in $L^2(a,b)$ has domain of definition
\[
\domain{\linA_{sl,F}} = H^2(a,b) \, \cap \, \HN{1}{a,b}.
\]
The operator $\linA_{sl,F}$ has pure point spectrum,
since the interval $(a,b)$ is assumed to be bounded.
Especially in case of $p(x) \eq 1$ and $q(x) \eq 0$ for $x \xeof [a,b]$, 
the eigen\-values and eigen\-functions
of the operator $\A_{sl,F} \eq - {d^2} / {dx^2}$ can exactly be determined.
In fact,
the eigenvalue distribution of $\A_{sl,F}$ is given by
\[
\lambda_n \= \frac{n^2 \pi^2}{(b-a)^2} \quad \quad (n \xeof \N),
\]
so the asymptotical behaviour of the eigen\-values is $\lambda_n \sim n^2$.
Each eigen\-value $\lambda_n$ owns a one-dimensional eigen\-space,
and the eigen\-function $u_n$ of $\lambda_n$ is given by
\[
u_n(x) \eqdef \sin \, (\frac{n \pi x}{b-a}) \quad \for x \in [a,b].
\]
Finally,
the sequence of the eigen\-functions $u_n$
constitutes after suitable renormations an ortho\-normal base in $L^2(a,b)$.

\item[(b)] \textit{(Negative) Laplacian in various domains $\G \xsubset \Rd$.}
In order to apply the theory of self-adjoint operators
to the investigation of the negative Laplacian,
we start with the realization $\linA_{-\Delta,0}$ of the expression $-\Delta$
in the space $\TF$ of test functions.

\subparagraph{} %- Green's formula for the negative Laplacian
Using \textsc{Green}'s second formula,
and due to the boundary condition $\phi(x) \eq 0$ for $x \xeof \partial \G$
and $\phi \xeof \TF$,
we compute for $\phi,\psi \xeof \TF$
\[
\int_{\G} -\Delta \phi \, \overline{\psi} \, dx
\=
\sum_{i=1}^n \int_{\G}
\frac{\partial^2 \phi}{\partial x_i^2} \, \overline{\psi} \, dx
\]
\[
\=
\sum_{i=1}^n \int_{\G}
u \, \overline{\frac{\partial^2 \phi}{\partial x_i^2}} \, dx 
\=
\int_{\G} \phi \, \overline{\Delta \psi} \, dx,
\]
which implies symmetry of the operator $\linA_{-\Delta,0}$.
Moreover,
we obtain for $\phi \xeof \TF$ the inequality
\[
\scalar{-\Delta \phi}{\phi} \=
\sum_{i=1}^n \int_{\G}
\frac{\partial^2 \phi}{\partial x_i^2} \, \overline{\phi} \, dx
\=
\sum_{i=1}^n \int_{\G} \abs{ \frac{\partial \phi}{\partial x_i}} ^2 \, dx 
\, \geq \, 0,
\]
hence $\linA_{-\Delta,0}$ is half-bounded with lower bound $0$
and fulfills the prerequisites
to apply the \textsc{Friedrichs} extension method.
Thus we may quote
that the operator $\linA_{-\Delta,0}$ owns at least one self-adjoint extension,
which is true for any domain $\G \xsubset \Rd$.

\paragraph{} % realizations of the Laplacian
In the following enumeration we consider various self-adjoint realizations
of the negative Laplacian in a domain $\G$,
all of them being extensions of the operator $\linA_{-\Delta,0}$:

\begin{itemize}[leftmargin=5mm]
\item[(i)]
\textit{(Negative) Laplacian in $\Rd$}.
The analysis of the operator $-\Delta$ in the entire Euclidean space $\Rd$
is easy,
since the differential operator proves to be unitarily equivalent to
the multiplication operator $\linA_m$ generated by the function
\[
m \: (x_1,\ldots,x_n) \mapsto (x_1^2 + \ldots + x_n^2) \= \abs{x} ^2
\quad (x \xeof \Rd).
\]
In fact,
the operator $\linA_m$ is real-valued and self-adjoint,
which is a consequence of
the \textit{Fourier-Plancherel transformation}\footnote{
\,We introduce the \textit{Fourier-Plancherel transformation} in $L^2(\Rd)$
by 
\[
\hat{u}(\xi) \eqdef
\frac{1}{(2 \pi)^{d/2}}
\int_{\,\Rd} \, e^{-2 \pi i \scalar{\xi}{x}} \, u(x) \, dx
\]
for all $u \in C_c^{\infty}(\Rd)$,
which forms a partial isometry $\FT \: u \mapsto \hat{u} $
from the space $C_c^{\infty}(\Rd)$ to the space $L^2(\Rd)$.
Then for functions $u,v \in L^2(\Rd)$ 
the inner product between $u$ and $v$
and the inner product of their Fourier-Plancherel transforms
$\hat{u}$ and $\hat{v}$ are related by \textsc{Plancherel}\textit{'s formula}
$\scalar{u}{v} \eq \scalar{ \hat{u} }{\hat{v} }$.
Since the mapping $\FT$ is linear and continuous on $C_c^{\infty}(\Rd)$,
which is a dense subspace of $L^2(\Rd)$,
it can be extended to an \textit{unitary},
that is,
a bijective and isometric linear operator
$\FT_P \: L^2(\Rd) \xto L^2(\Rd)$,
called the \textsc{Fourier-Plancherel} \textit{operator}
or the \textit{$L^2$-Fourier transformation} in $L^2(\Rd)$.}
.
This implies
that the operator $\linA_{sa}$ with domain of definition
$\domain{\linA_{sa}} := H^2(\Rd)$ must be self-adjoint,
too.
The spectrum of $\linA_{sa}$ is that of the operator $\linA_m$,
that is,
it is purely continuous and $\sigma_c(\linA_{sa}) \eq [0,\infty)$.
Moreover,
since $\linA_{sa}$ is the closure of the operator $\linA_{-\Delta,0}$,
the latter proves to be essentially self-adjoint.

\item[(ii)]
\textit{Dirichlet-Laplacian}.
Let the domain $\G \xsubset \Rd$ be \textit{bounded}
and consider $-\Delta$ with domain of definition $\domain{-\Delta} := \TF$.
The related Dirichlet form $\DF_{-\Delta}$ is given by 
\[
\DF_{-\Delta}(\phi,\psi) \eqdef
\int_{\G} \nabla \phi \, \nabla \psi \, dx
\quad \quad
(\phi,\psi \xeof \TF),
\]
which can continuously be extended to the sesquilinear form
$\overline{\DF_{-\Delta}}$ with domain of definition
$\domain{\overline{\DF_{-\Delta}}} := \HN{1}{\G} \xtimes \HN{1}{\G}$.

\subparagraph{}
Due to theorem \ref{QUASI-COERCIVE_FORMS_AND_ASSOCIATED_OPERATORS_II},
the operator $\linA_D$ associated with the form $\DF_{-\Delta}$
and defined by
\[
\begin{cases}[1.2]
\quad \domain{\linA_D} \eqdef
\{u \xeof \HN{1}{\G}: \Delta u \xeof L^2(\G)\} \\
\quad \linA_D u \eqdef \, - \Delta u \; \for u \xeof \domain{\linA_D} \\
\end{cases}
\] 
is self-adjoint,
called the \textit{negative Laplacian with Dirichlet boundary conditions}
or briefly the \textit{Dirichlet-Laplacian}.

\subparagraph{}
We obviously have
$\HN{1}{\G} \cap H^2(\G) \xsubset \domain{\linA_D}$,
and one can prove that
if $\G$ is bounded and owns a $C^2$-boundary or $\G \eq \R_+^n$,
then the reversed inclusion also holds,
hence $\domain{\linA_D} \eq \HN{1}{\G} \cap H^2(\G)$.
This equality is even true in the case,
where $\G$ is bounded and convex.

\subparagraph{}
Furthermore,
if $\G$ has the $1$-extension property,
then the space $\HN{1}{\G}$ is compactly embedded into the space $\Lh$
and $\linA_D$ is an operator with pure point spectrum.
In this case it is an important property of the Dirichlet-Laplacian
that the number $0$ does \textit{not} belong to the spectrum $\sigma(\linA_D)$,
so $\linA_D$ is invertible and the operator $\linA_D^{-1}$ proves to be bounded.

\item[(iii)]
\textit{Neumann-Laplacian}.
Let $\G \xsubset \Rd$ be a bounded domain with $C^1$-boundary $\partial \G$
and let the realization $\linA_N$ of the negative Laplacian in $\Lh$
be defined by
\[
\begin{cases}[1.2]
\quad \domain{\linA_N} \eqdef
\{ u \xeof H^2(\G): u_v(x)= 0, \; x \xeof \partial \G \} \\
\quad \linA_N u \eqdef - \Delta u \; \for u \xeof \domain{\linA_N}. \\
\end{cases}
\]
The realization $\linA_N$ is said to be the
\textit{(negative) Laplacian with Neumann boundary condition}
or briefly the \textit{Neumann-Laplacian}.
It owns $\lambda \eq 0$ as the minimal eigen\-value,
and all constant functions in $\domain{\linA_N}$ are eigen\-functions
of $\lambda \eq 0$.
\end{itemize}

\item[(c)] \textit{{\SCHROEDINGER} operators.}
By the \textit{{\SCHROEDINGER} operator} is meant
the linear partial differential expression given by
\[
\HAM_q(u,x) \eqdef (-\Delta u)(x) \+ q(x),
\]
defined in the entire Euclidean space $\Rd$.
In the special case,
where $q \eq 0$,
self-adjointness of the operator $-\Delta$ acting in $L^2(\Rd)$
and with domain of definition $H^2(\Rd)$ is well-known.

\subparagraph{} % potential functions and multiplication operators
The so-called \textit{potential function} $q \: \Rd \xto \R$
gives rise to a multiplication operator $M_q$ acting in $L^2(\Rd)$,
that is,
\[
\begin{cases}[1.2]
\quad \domain{M_q} \eqdef \{u \in L^2(\Rd): q {\cdot} u \in L^2(\Rd)\},\\
\quad M_q u        \eqdef q {\cdot} u \; \for u \in \domain{M_q}   ,\\
\end{cases}
\]
so the operator $\HAM_q$ can be seen as a perturbation
of the negative Laplacian by the operator $M_q$.

\subparagraph{} %- application of the Kato-Rellich theorem
As an application of the \textsc{Kato-Rellich} theorem
\ref{KATO-RELLICH_THEOREM_II} let us consider the case,
where $d \eq 3$ and the function $q$ belongs to $L^2(\R^3) + L^{\infty}(\R^3)$,
{\ie}
$q \eq q_1 + q_2$ for functions $q_1, q_2$
with $q_1 \xeof L^2(\R^3)$ and $q_2 \xeof L^{\infty}(\R^3)$.
Since of the negative Laplacian $-\Delta$ with domain of definition $\TF$
is essentially self-adjoint in $L^2(\R^3)$
and $M_q$ is $-\Delta$-bounded with bound $\theta {\,<\,} 1$,
the operator $\HAM_q := -\Delta + M_q$
with domain of definition $C_c^{\infty}(\R^3)$
is essentially self-adjoint in $L^2(\R^3)$,
too.
In conclusion,
the operator $\HAM_q$ with domain of definition $\domain{\HAM_q} \eq H^2(\R^3)$
proves to be self-adjoint.

\paragraph{} %- Schroedinger equation
By the time-dependent \textit{{\SCHROEDINGER} equation} is meant
the initial value problem
\[
i \, u' \= -\Delta u \+ M_q u
\]
with prescribed initial value $u(0) \eq u_0 \in H^2(\Rd)$,
where it is supposed that $\HAM_q \eq -\Delta + M_q$
is a self-adjoint operator in $L^2(\Rd)$
with domain of definition $\domain{\HAM_q} \eq H^2(\Rd)$.
The existence and uniqueness of a solution for this problem is ensured
by the functional calculus for self-adjoint operators
using the spectral resolution of the operator $\HAM_q$
and the function $f:\lambda \in \R \mapsto e^{it \lambda}$.
The solution $u$ of the problem is then provided
by a group of unitary operators $\{U(t)\}_{t \in \R}$,
which implies $u(t) \eq U(t) u_0$ for every $t \in \R$
(see also section 5 for the treatment of the abstract {\SCHROEDINGER} equation).

\paragraph{} %- spectral analysis of Schroedinger operators
The study of \textit{time-independent solutions}
of the {\SCHROEDINGER} equation leads to the eigenvalue problem
\begin{equation}\label{SCHROEDINGER_EIGENVALUE_PROBLEM}
\HAM_q u \eqdef -\Delta u + M_q u \= \lambda u,
\end{equation}
which means to find those real numbers $\lambda$,
for which there are non-trivial solutions
of \eqref{SCHROEDINGER_EIGENVALUE_PROBLEM}.
The discrete spectrum $\sigma_d(\HAM_q)$ of the operator $\HAM_q$
is in fact of great interest,
since the eigen\-functions of eigen\-values $\lambda \xeof \sigma_d(\HAM_q)$
describe the stationary states of a quantum mechanical system.

\subparagraph{} %- spectrum of Schroedinger operators
The presence of a potential function $q$
can fairly change the spectrum of a {\SCHROEDINGER} operator 
like the example of the \textit{harmonic oscillator} shows:
the spectrum of the negative Laplacian is the connected set $[0,\infty)$,
but the spectrum of the operator $-\Delta {\,+\,} M_q$ 
with $q \: x \mapsto \frac{1}{2}x^2$ consists of discrete points only.

\subparagraph{} %- lower and upper bounds for the discrete spectrum
It is a further and interesting problem in quantum mechanics
to determine lower and upper bounds for the discrete spectrum
of the operator $\HAM_q$.
In the following we want to present three assertions,
which contain information on the bounds of the set $\sigma_d(\HAM_q)$,
namely
the \textsc{Birman-Schwinger} \textit{bound}, 
the \textsc{Lieb-Cwikel-Rozenblum} \textit{bound} and
    \textsc{Kato}\textit{'s theorem}:

\begin{itemize}[leftmargin=4mm]
\item[1.]
The \textsc{Birman-Schwinger} \textit{bound.}
This assertion presents an upper bound for the set $\sigma_d(\HAM_q)$,
which is valid for certain potentials $q$ in case of $d \eq 3$.
Let $q \in L^{\infty}(\R^3)$ and 
\[
B_q \eqdef
\int_{\,\R^3} \int_{\,\R^3} \frac{q(x) \, q(y)}{\abs{x-y} ^2} \ dx \, dy
\, < \, + \infty.
\]
Then the total number $N(\HAM_q)$ of eigen\-values of the operator $\HAM_q$
can be estimated according to
\[
N(\HAM_q) \, \leq \, \frac{1}{16 \pi^2} \, B_q.
\]

\item[2.]
The \textsc{Lieb-Cwikel-Rozenblum} \textit{theorem}
quotes an estimate on the number $N_(\HAM_q)$
of negative eigen\-values of the operator $\HAM_q$:
if $d > 3$ and $q \xeof L_{loc}^{\infty}(\Rd)$,
then
\[
N(\HAM_q) \, \leq \, c_d \int_{\,\Rd} \min \{q(x),0\}^{n/2} \, dx,
\]
where $c_d$ is a constant only depending on the dimension $d$.

\item[3.]
Our last result is often denoted as \textit{Kato}\textit{'s theorem}.
It does not make assertions on bounds for the discrete spectrum,
but quotes the absence of positive eigen\-values
under mild assumptions on the function $q$:
if $q \in L_{loc}^{\infty}(\Rd)$ and
\[
\lim_{\abs{x} \to \infty} \, \abs{x} \, q(x) \= 0,
\]
then the operator $\HAM_q$ does not own any positive eigen\-value.
\end{itemize}
\end{itemize}

%------------------------------------------------------------------------------%
\section{Distribution Theory and Differential Equations}                       %
%------------------------------------------------------------------------------%
A \textit{distribution} my be regarded as a generalization of a function.
The corresponding subject in functional analysis,
called the \textit{theory of distributions},
was developed by \textsc{L.\,Schwartz} in the midth of the 20th century
to describe ideas and concepts of theoretical physics
by an exact mathematical formalism.
Distribution theory can be seen as a direct application
of the previously sketched theory including
locally convex spaces, continuous linear functionals and
transposed linear operators.

\paragraph{} %-- families of seminorms and locally convex topologies
%------------------------------------------------------------------------------%
The class of \textit{locally convex spaces} can be seen as the most important
one within all kinds of topological linear spaces.
The topology $\topT_{\,\vecV}$ of a locally convex space $\vecV$ distinguishes
by the fact
that its origin $\0$
(and thus any other element in $\vecV$) owns a neighbourhood base
consisting of \textit{convex}, \textit{balanced} and \textit{absorbing} sets.

\subparagraph{} %- seminorms
It turns out
that each such pair $(\vecV,\topT_{\,\vecV})$ gives rise to
a family of \textit{seminorms} on the space $\vecV$.
By the concept of seminorm
we mean a mapping $\abs{\,\cdot\,} \:\vecV \xto \R$ having the properties

\begin{itemize}[leftmargin=55pt,itemsep=2pt]
\item[(SN-1)] \quad
$\abs{\lambda x} \= \abs{\lambda} \cdot \abs{x} \for \lambda \in \F$
\textit{(homogenity)},
\item[(SN-2)] \quad
$\abs{x+y} \, \leq \, \abs{x} + \abs{x} \for x,y \in \vecV$
\textit{(subadditivity)},
\end{itemize}

which especially implies $\abs{x} \leq 0$ for $x \in \vecV$ and $\abs{\0} \eq 0$.

\subparagraph{} %- Minkowski functional
Moreover,
if $\U \xsubset \vecV$ forms a convex, balanced and absorbing neighbourhood
of the origin $\0$,
then the \textit{Minkowski functional}
\[
\begin{cases}[1.3]
\quad \abs{\,\cdot\,} _{\,\U}\:\vecV \xto \R_0^+,\\
\quad x \mapsto \abs{x} _{\,\U} \, \eqdef \,
\mbox{inf} \, \set{\lambda \geq 0}{ x \in \lambda \, \U} \\
\end{cases}
\]
fulfills all the properties of a seminorm.
On the other hand,
a family $\famF$ of seminorms defined on a linear space $\vecV$
leads in a natural way to a locally convex topology $\topT_{\,\vecV}$:
it is the coarsest topology on $\vecV$
such that all seminorms in $\famF$ remain continuous.

\subparagraph{}
In summary,
there is a one-to-one correspondence
between linear spaces endowed with locally convex topologies
and those endowed with families of seminorms.
Notice that a locally convex topology is Hausdorff
iff
the family of seminorms is \textit{separating},
{\ie} for each element $x \neq 0$
there is at least one seminorm $\abs{\,\cdot\,} _{\alpha}$
such that $\abs{x} _{\alpha} \neq 0$.

\paragraph{} %- F-spaces and Frechet spaces
There are a two subclasses of locally convex spaces,
we want to review in the following:

\begin{itemize}[leftmargin=4mm]
\item[1.]
A topological linear space $\vecV$ is called \textit{F-space},
if it is also a \textit{complete} metric space with translation-invariant
metric $\d$,
{\ie}
\[
\d(x,y) \= \d(x+z,y+z) \; \fa x,y,z \in \vecV,
\]
and the metric topology coincides with the given topology on $\vecV$.

\item[2.]
If the topology of the F-space $\vecV$ is locally convex,
then $\vecV$ is said to be a \textit{{\FRECHET} space}\footnote{
A metrizable locally convex space that is not necessarily complete,
is called a \textit{pre-{\FRECHET} space}.}.
Alternatively,
a locally convex space is {\FRECHET}
\iff
it is complete and its family of seminorms is mostly countable.
In this case
we may define a translation-invariant metric on a {\FRECHET} space $\vecV$
according to
\[
\d(x,y) \, \eqdef \,
\sum_{n=1}^{\infty} p_n \cdot \frac{\abs{x-y} _n}{1+ \abs{x-y} _n}
\quad (x,y \in \topV)
\]
where $\{p_n\}_{n=1}^{\infty}$ is any sequence of positive real numbers,
whose infinite sum converges.
\end{itemize}

\paragraph{} %- continuous linear functionals and adjoint spaces
Given a topological $\F$-linear space $\banX$,
let $\banX^*$ be the set of all continuous linear mappings
$\phi\:\banX \xto \F$,
called the \textit{adjoint space} $\banX^*$ of $\banX$,
which forms a linear space again.
If the topology of $\banX$ is at least locally convex,
then a specific version of the celebrated \textsc{Hahn-Banach} \textit{theorem}
ensures that
there are non-trivial continuous linear mappings on $\banX$,
thus we may assume $\banX^* \neq \{\0\}$.
We denote the elements of the adjoint space $\banX^*$ as
\textit{continuous linear functionals}.

\paragraph{} %- bi-adjoint spaces and weak*-topologies
%------------------------------------------------------------------------------%
It is known from linear algebra
that an ordinary $\F$-linear space $\vecV$ may be regarded
as a linear subspace of its (algebraic) bi-adjoint space
$\vecV^{\#\#} \eqdef (\vecV^{\#})^{\#}$,
where $\vecV^{\#}$ is the space of linear functionals $\phi\:\vecV \xto \F$.
More precisely,
every element $x \in \vecV$ gives rise to a linear functional
$F_x\:\vecV^{\#} \xto \F$ according to
$F_x(\phi) \eqdef \phi(x)$ for all $\phi \in \vecV^{\#}$.
The natural mapping $\iota\:\vecV \xto \vecV^{\#\#}$,
defined by $\iota(x) \eqdef F_x$ for $x \in \vecV$,
is linear and injective but in general not surjective.

\subparagraph{} %- weak*-topologies
Having these facts in mind,
we may introduce the \textit{weak*}-topology on the adjoint $\banX^*$
of a locally convex topological linear space $\banX$:
it is the smallest topology in $\banX^*$
such that
for each $x \in \banX$ the linear functional $F_x$ remains continuous.

\paragraph{} %-- continuous linear operators in topological linear spaces
%------------------------------------------------------------------------------%
Let us again have a look at continuous linear mappings in the context
of topological linear spaces.
As for normed linear spaces,
the notion of \textit{continuity} is of fundamental interest
for linear operators between topological linear spaces $\vecV$ and $\vecW$.
It turns out that a linear operator $\linA \: \vecV \xto \vecW$ is continuous
\iff
it is continuous in $0 \xeof \vecV$ or in any other point $x \neq 0$.
But continuity is not equivalent to boundedness in this general setting.

\subparagraph{} %- continuous linear maps between topological linear spaces
If $\vecV$ and $\vecW$ are topological linear spaces,
then $\vecV$ is called \textit{continuously embedded} into $\vecW$,
if there is an injective and continuous linear operator
$\iota \: \vecV \xto \vecW$.
In addition,
if this mapping is surjective and both spaces are {\FRECHET},
then the \textit{open mapping theorem} claims that $\iota^{-1}$ is continuous,
too,
and therefore forms a \textit{linear homeomorphism}
between these spaces (notation: $\vecV \cong \vecW$).

\paragraph{} %-- transposed operators in topological linear spaces
%------------------------------------------------------------------------------%
The concept of \textit{transposed operator} is useful to generalize the domains
of linear partial differential operators and Fourier transformation
to the very extensive spaces of distributions.
Let $\vecV$ and $\vecW$ be locally convex and $\vecV^*, \vecW^*$
the dual spaces of them.
Given a continuous linear operator $\linA \: \vecV \xto \vecW$,
each linear functional $f \xeof \vecW^*$ gives rise to
a linear functional $g \xeof \vecV^*$ by the composition $g:= f \circ \linA$.
The assignment $f \mapsto g$ induces a continuous mapping
$\linA^t \, :\vecW^* \xto \vecV^*$,
called the \textit{trans\-posed operator} of $\linA$.
One says that
the continuous linear functional $g$ is the \textit{pullback} of $f$
along the operator $\linA$.

\subparagraph{}
The transposed operator $\linA^t$ is also linear.
If $\vecV, \vecW$ are locally convex linear spaces
and $\linA \: \vecV \xto \vecW$ is continuous,
then $\linA^t$ is injective
iff $\range{\linA} $ is dense in $\vecW$.
Moreover,
if $\linA$ is bijective,
then $\linA^t$ is bijective, too,
and its inverse $(\linA^t)^{-1}$ is the transposed of $\linA^{-1}$.

\subparagraph{}
Finally,
if $\vecV$ and $\vecW$ are {\FRECHET} spaces,
then the solvability of the linear equation $\linA x \eq y$ 
is related to certain properties of the transposed operator $\linA^t$:
the linear operator $\linA \: \vecV \xto \vecW$ is surjective
iff
$\linA^t$ is injective and $\range{\linA^t} $ is closed in $\vecV^*$.

\paragraph{} %-- the space of test functions
%------------------------------------------------------------------------------%
The space $\TF$ of \textit{test functions} is of fundamental meaning
in distribution theory and its starting point.
Remember that this space forms a linear subspace of $\SF$
consisting of smooth functions with \textit{compact} support
\[
\supp{\,\phi} \, := \, \overline{\{x \xeof \G: \phi(x) \neq 0\,\} }.
\]
There is neither a countable family of seminorms
making $\TF$ into a {\FRECHET} space nor there is a metric making 
$\TF$ into a F-space.

\subparagraph{}
But it is still possible to define a locally convex topology $\topT$ on $\TF$
such that this space is complete in the sense of a topological linear space.
Instead of saying how the open sets of this topology $\topT$ look like,
we define $\topT$ by means of a suitable concept of convergence,
known as $\D$-convergence in $\TF$.

\subparagraph{}
A sequence $(\phi_n)$ in $\TF$ is said to be convergent to $\phi \xeof \TF$,
if there is a compact set $K \xsubset \G$
such that

\begin{itemize}[leftmargin=15mm]
\item[(D-1)]
$\supp{\phi_n} \xsubset K$ for all $n \xeof \N$,

\item[(D-2)]
for each multi-index $\alpha$ the sequence $(\PA \phi_k)$
converges to the function $\PA \phi$ \textit{uniformly} in $K$,
\end{itemize}
the linear space $\TF$ equipped with the topology
induced by $\D$-convergence is usually denoted $\D(\G)$.
We have $\D(\G) \hookrightarrow \E(\G)$,
that is,
the space $\D(\G)$ is continuously embedded into the space $\E(\G)$
{\wrt} their locally convex topologies.

\paragraph{} %- linear functionals on the space of test functions
%------------------------------------------------------------------------------%
The benefit of the function space $\TF$ relies on the fact that,
due to the smoothness and compact support property of its elements,
the integral
\begin{equation}\label{REGULAR_DISTRIBUTION}
\disF_f(\phi) \, := \, \int_{\,\G} f \, \phi \, dx, \quad \quad \phi \xeof \TF
\end{equation}
is well defined for a multiplicity of real or complex-valued functions $f$.
The largest class of functions
for which this integral has a reasonable meaning,
is the class $\Lc$ of \textit{locally integrable functions},
that is,
$f \xeof \Lc$ is Lebesgue-integrable on each compact subset $K \xsubset \G$.
This space incorporates a lot of other well-known function spaces,
for example all the spaces $L^p(\G)$ with $p \geq 1$,
the space $L^{\infty}(\G)$ of essentially bounded functions
and the space $C(\G)$ of continuous functions.
Moreover, 
if $\mu$ is a positive measure on $\G$,
the integral
\begin{equation}\label{RADON_DISTRIBUTION}
\disF_{\mu}(\phi) \, := \, \int_{\,\G} \phi \, d\mu, \quad \, \phi \xeof \TF
\end{equation}
especially makes sense for all the Radon measures $\mu$ on $\G$,
since such measures are finite on compact sets by definition.
Both mappings $\phi \mapsto \disF_f$ and $\phi \mapsto \disF_{\mu}$
as defined in \eqref{REGULAR_DISTRIBUTION} and \eqref{RADON_DISTRIBUTION}
act linearly on $\TF$ and hence give rise to linear functionals.
But both of them are also continuous on the topological linear space $\D(\G)$.

\paragraph{} %- Schwartz distributions
%------------------------------------------------------------------------------%
Let us denote each continuous linear functional $\disF \: \D(\G) \xto \C$ 
a \textit{Schwartz distribution}.
We call $\D'(\G)$ the dual space of $\D(\G)$,
so $\D'(\G)$ is the linear space of Schwartz distributions on $\G$.
One may also define Schwartz distributions {\wrt} $\D$-convergence,
when applied to test functions:
a linear functional $\disF$ on the space $\TF$
is denoted a Schwartz distribution,
if $\disF(\phi_n) \to 0$ for any $\D$-convergent sequence $(\phi_n)$
converges to the zero function.

\subparagraph{}
As already mentioned,
all linear functionals derived from locally integrable functions
and Radon measures are continuous and therefore Schwartz distributions.
Indeed,
the space $\D'(\G)$ is extremely extensive
and contains nearly all functions being important in applications,
amongst others the functions
belonging to the well-known spaces $\Lp$, $\Li$, $C(\G)$,
$C(\clG)$ and $C_0(\G)$.

\subparagraph{}
A distribution is said to be \textit{regular},
if it is the continuous linear functional $\disF_f$
induced by a locally integrable function $f$,
otherwise it is called \textit{singular}.
Well-known examples of singular distributions are
the \textit{Dirac distributions}
\[
\delta_a(\phi) \, := \, \phi(a)
\qfa \phi \xeof \TF
\quad \quad  (a \xeof \G).
\]

\paragraph{} %- space of distributions with compact support
%------------------------------------------------------------------------------%
Regarding the space $\E(\G)$,
which is the space $\SF$ in the locally convex topology
induced by the family of seminorms
\[
\abs{\phi} _{N,n}
\, := \,
\max_{x \xeof K_n} \;
\{ \, \abs{\, (\PA \phi)(x) \,} : \abs{\alpha} \leq N \}
\quad \quad (N,n \xeof \N)
\]
for each $\phi \xeof \SF$,
where $(K_n)$ is a sequence of compact subsets of $\G$ with
\[
K_1 \xsubset K_2 \ldots \xsubset \G,
\quad \quad
\bigcup_{n \xeof \N} \, K_n \= \G
\]
each element $\disF \xeof \E'(\G)$ is called
a \textit{distribution with compact support}.
This notion is founded by the fact
that for such $\disF$ there is a compact set $K_F \xsubset \G$,
called the \textit{support} of $\disF$
such that
$\disF \phi \eq 0$ for all smooth functions $\phi$
having support in $\G \setminus K_F$.

\subparagraph{}
The denotation for this class of distributions is motivated as follows:
if $\disF_f \xeof \E'(\G)$ is regular,
then the function $f \xeof \Lc$ related to $\disF_f$
has compact support, too.
On the other hand,
each measurable function $f$ with compact support of course
gives rise to a compactly supported distribution $\disF_f \xeof \E'(\G)$.
Since it is always possible to increase a dual space
by decreasing the underlying linear space,
there is a continuous inclusion $\E'(\G) \, \hookrightarrow \, \D'(\G)$,
which is dual to the mentioned embedding
$\D(\G) \, \hookrightarrow \, \E(\G)$.
So we basically may regard
each compactly supported distribution as a Schwartz distribution.

%------------------------------------------------------------------------------%
\paragraph{} %-- convolution and tensor products                               %
%------------------------------------------------------------------------------%
The usage of distributions for the solution theory of linear partial
differential equations can only be exhausted in full strength
if the notion of \textit{convolution} between elements
of the spaces $\D'(\Rd)$ and $\E'(\Rd)$ has been introduced
as far as possible.

\subparagraph{}
The convolution $\phi \ast \psi$ of Schwartz functions $\phi, \psi \xeof \SRd$
is defined according to
\[
(\phi \ast \psi) (x) \, := \, \int_{\,\Rd} \phi(y) \, \psi(x-y) \, dy
\]
and always exists,
hence the space of Schwartz functions forms an algebra
{\wrt} convolution as multiplicative operation.
Furthermore,
convolution between a regular distribution
$\disF_f \xeof \D'(\Rd)$ and a test function $\phi \xeof \D(\Rd)$
is similarly defined:
\[
(\disF_f \ast \phi) (x) \, := \, \int_{\,\Rd} f(y) \, \phi(x-y) \, dy.
\]

\subparagraph{}
More general,
the convolution $\disF * \disG$ of distributions
$\disF, \disG \xeof \D'(\Rd)$ exists,
if the \textit{support condition} is fulfilled,
meaning that for each compact subset $K \xsubset \Rd$ the set
\[
K_{\disF,\disG} :=
\{(x,y): x \xeof \supp{F}, y \xeof \supp{G}, x+y \xeof K \}
\xsubset \Rd \xtimes \Rd
\]
is compact, too.
In this case we may define the convolution $\disF \ast \disG$ according to
\[
(\disF \ast \disG) \, \phi \, := \,
\iint \limits_{\compK_{\disF,\disG}}
\disF(x)\,\disG(y)\,\phi(x+y) \, dx \, dy, \quad \phi \xeof C_c^{\infty}(\Rd).
\]
The support condition is sufficient
but not necessary for the existence of the convolutional distribution
$\disF \ast \disG$.
Nevertheless,
the condition is always fulfilled
if at least one of the distributions $\disF, \disG$
is compactly supported.
Notice also that
convolution on the space $\D'(\Rd)$ is commutative, if exists.

\subparagraph{}
Especially the convolution between any distribution $\disF \xeof \D'(\G)$
and the Dirac distribution $\delta_0$ exists and
fulfills $\disF \ast \delta_0 \eq \disF$,
thus $\delta_0$ serves as identity element of this operation.
But convolution is also associative,
hence for distributions $\disF, \disG, \disH \xeof \D'(\Rd)$ the equality
\[
(\disF \ast \disG) \ast \disH \= \disF \ast (\disG \ast \disH)
\]
holds provided that
the involved convolutions involved exist.

\subparagraph{}
Besides the notion of convolution one also requires
that of \textit{tensor product} between distributions.
Tensor products are useful to transform an inital value
of a distributional differential equation
into an inhomogeneous part of that equation.
Based on the fact
that for $\phi \xeof \D(\R^{d_1})$ and $\psi \xeof \D(\R^{d_2})$ 
the tensor product $\phi \otimes \psi$ defined by
\[
(\phi \otimes \psi)(x,y) \= \phi(x) \cdot \psi(y)\,,
\quad (x,y) \xeof \R^{d_1+d_2}
\]
belongs to the space $\D(\R^{d_1+d_2})$,
there is
a uniquely defined distribution
$
\disF \otimes \disG \xeof \D'(\R^{d_1+d_2})
$
for $\disF \xeof \D'(\R^{d_1})$ and $\disG \xeof \D'(\R^{d_2})$,
which for all test functions
$\phi \xeof \D(\R^{d_1})$ and $\psi \xeof \D(\R^{d_2})$
fulfills the property
\[
\disF \otimes \disG \, (\phi \otimes \psi)
\= \disF(\phi) \cdot \disG(\psi)\,,
\]
also denoted the \textit{tensor product} of $\disF$ and $\disG$.
Performing the tensor product is a commutative and associative operation,
that is
\[
\disF \otimes \disG = \disG \otimes \disF,
\quad \quad
(\disF \otimes \disG) \otimes \disH =
 \disF \otimes (\disG \otimes \disH).
\]

\subparagraph{}
Let $\PA$ be the elementary differential expression
given by the multi-index $\alpha$.
Then $\PA$ is a continuous linear operator $\PA \:\D(\G) \xto \D(\G)$,
due to the fact
that the space $\D(\G)$ is endowed with a very fine topology.

\subparagraph{}
The action $\PA_{\D}$ of the expression $\PA$
on the distributional space $\D'(\G)$
is defined by transposition of the \textit{formally adjoint} expression
$(\PA)^* := (-1)^{\alpha} \PA$,
or equivalently
\begin{equation}\label{DISTRIBUTIONAL_DERIVATIVE}
(\PA_{\D} \disF) \, \phi \, := \, \disF((\PA)^* \phi)
\= (-1)^{\abs{\alpha} } \, \disF(\PA \phi)
\end{equation}
for all $\disF \xeof \D'(\G)$ and $\phi \xeof \D(\G)$.

\subparagraph{}
It should be emphasized
that this definition is chosen in order
to obtain compatibility with possible classical derivatives
of \textit{regular} distributions.
In fact,
if $F_f \xeof \D'(\G)$ is the distribution generated
by a differentiable function $f \:\G \xto \C$
as defined in \eqref{REGULAR_DISTRIBUTION},
then $\PA_{\D} F_f \eq F_{\PA f}$ holds.

\subparagraph{}
The mapping $\disF \mapsto \PA_{\D} \disF$ proves to be continuous
on the Schwartz space $\D'(\G)$.
We also have the following compatibility rules
between convolutions and tensor products:
\[
\PA_{\D} (\disF \ast \disG)  \=
(\PA_{\D} \disF) \ast \disG \=
\disF \ast (\PA_{\D} \disG) \,,
\]
\[
 \PA_{\D} (\disF \otimes \disG) =
(\PA_{\D} \disF) \otimes \disG  =
 \disF           \otimes (\PA_{\D} \disG).
\]

\paragraph{} %- Schwartz space
%------------------------------------------------------------------------------%
We continue with the introduction of the \textit{Schwartz space} $\SRd$,
the linear space of all smooth functions $\phi \xeof C^{\infty}(\Rd)$
having the additional property
\[
\abs{ \PA \phi(x) \cdot x^{\beta}} \to 0 \,
\for \abs{x} \to \infty
\mbox{ and all multi-indices } \alpha,\beta,
\]
where $\abs{x} $ denotes the Euclidean norm of any vector $x \xeof \Rd$.
Then a function $\phi \xeof \SRd$ is called \textit{rapidly decreasing}
or briefly a \textit{Schwartz function}.
Hence Schwartz functions and all their derivatives vanish at infinity
faster then the reciprocal of any polynomial.
The most famous example of a Schwartz function is the Hermite function
\[
\cal{H}_0(x) \, := \, e^{- \frac{\abs{x} ^2}{2}}, \quad x \xeof \Rd.
\]

\subparagraph{}
Rather similar to the space $\SF$,
one can find for the Schwartz space $\SRd$
a family of seminorms as follows:
define for $\phi \xeof C^{\infty}(\Rd)$ and multi-indices $\alpha, \beta$
\begin{equation}\label{X_SRN_FAMILY_OF_SEMINORMS}
p_{\alpha,\beta}(\phi) \, := \,
\sup_{x \xeof \Rd} \abs{ \, x^{\alpha} \partial^{\beta} \phi(x) \, } .
\end{equation}

\subparagraph{}
Then $\phi$ is a Schwartz function
iff $p_{\alpha,\beta}(\phi) < \infty$
holds for all multi-indices $\alpha, \beta$.
The collection of mappings $p_{\alpha,\beta} \:\SRd \xto \R$
defines a countable family of seminorms,
which gives rise to a locally convex topology $\topT_{\mathscr{S}}$ on $\SRd$.
Moreover,
since the Schwartz space is complete {\wrt} $\topT_{\mathscr{S}}$,
it actually forms a \textit{{\FRECHET} space}.

\subparagraph{}
Comparing the Schwartz space $\SRd$ with the function spaces
$\D(\Rd)$ and $\E(\Rd)$,
it is easy to see that there is the chain of inclusions
\[
\D(\Rd) \, \hookrightarrow \, \SRd \, \hookrightarrow \, \cal{E}(\Rd).
\]
Moreover,
by comparing the locally convex topologies of these spaces,
it turns out that these inclusions are actually continuous.

\paragraph{} %- the space of tempered distributions
%------------------------------------------------------------------------------%
Similar to the spaces $\D(\G)$,
we introduce the concept of \textit{tempered distribution}
as a continuous linear functional $T \:\SRd \xto \C$.
The space of tempered distributions is usually denoted $\TD$.
A tempered distribution generalizes the class of bounded and slow-growing
locally integrable functions.

\subparagraph{}
The following criterion determines
under which conditions a linear functional $f$ on $\SRd$
forms a tempered distribution:
this is the case 
\iff
there is a constant $C > 0$ and integers $m,l \xeof \N$ such that
\[
\abs{ f(\phi) } \, \leq \,
C \cdot \sum_{\abs{\alpha} \leq l, \abs{\beta} \leq m} \, p_{\alpha,\beta}(\phi)
\qfa \phi \xeof \SRd\,,
\]
where the functions $p_{\alpha,\beta}$ are the seminorms
defined by \eqref{X_SRN_FAMILY_OF_SEMINORMS}.

\paragraph{} %-- comparison of distributional spaces
%------------------------------------------------------------------------------%
Comparing the known kinds of distributional spaces,
we obtain a chain of inclusions,
which is dual to that of the base spaces:
\[
\cal{E}'(\Rd) \, \xsubset \, \TD \, \xsubset \, \D'(\Rd).
\]
Endowing each of these spaces with the so-called \textit{weak*-topology},
these embeddings are continuous.
Hence a distribution having compact support can also be regarded
as a tempered distribution,
which in turn gives rise to a Schwartz distribution.

\paragraph{} %-- Fourier transforms in distributional spaces
%------------------------------------------------------------------------------%
Ee next shall elaborate that Fourier transformation can be
extended to the space $\TD$ of tempered distributions.
First notice 
that the Fourier transform $\hat{\phi}$
of a Schwartz function $\phi \xeof \SRd$ is well defined,
since $\SRd$ is a linear subspace of the convolution algebra $L^1(\Rd,+,\ast)$.
Moreover, 
Fourier transformation restricted to the Schwartz space $\SRd$
forms a linear mapping $\FT \:\SRd \xto \SRd$,
which sends each function $\phi \xeof \SRd$ to a further function
$\hat{\phi} \xeof \SRd$ according to the well-known formula
\begin{equation}\label{FOURIER_TRANSFORM_OF_SCHWARTZ_FUNCTION}
\hat{\phi}(x) \= \FT(\phi)(x) :=
\frac{1}{(2\pi)^{d/2}}
\int_{\,\Rd} e^{-i <\xi,x>} \phi(\xi) \, d \xi
\quad (x \xeof \Rd).
\end{equation}

\subparagraph{}
The extraordinary meaning of Schwartz functions is founded by the fact that
the Fourier transformation has excellent properties on the space $\SRd$.
More precisely,
besides the well-known features of Fourier transformation in $L^1(\Rd,+,\ast)$,
the mapping $\FT$ is \textit{bijective} on $\SRd$.
Hence the \textit{inverse Fourier transform} exists for all $\phi \xeof \SRd$
and is described by the \textit{Fourier inversion formula}
\[
\check{\phi}(\xi) \= 
\FT^{-1}(\phi)(\xi) \eqdef
\frac{1}{(2\pi)^{d/2}}
\int_{\,\Rd} e^{i \scalar{x}{\xi}} \, \phi(x) \, dx
\quad \quad (\xi \xeof \Rd).
\]

\subparagraph{}
We next extend Fourier transformation in $\SRd$ to the space $\TD$
of tempered distributions by transposition,
that is
\[
(\FT \disT) \, \phi \, := \, \disT(\FT \phi),
\quad \quad \; \phi \xeof \SRd, \; \disT \xeof \TD.
\]
This mapping is bijective in the space $\TD$
and even continuous {\wrt} the weak*-topology on $\TD$.
The inverse Fourier transform of a tempered distribution $\disT$
is represented by the formula
\[
\FT^{-1}(\disT)(\phi) \, := \, \disT(\FT^{-1} \phi)
\quad \quad (\phi \xeof \SRd).
\]
For example,
the Fourier transform $\FT(\delta_0)$ of the Dirac distribution
$\delta_0 \xeof \cal{S}'(\R)$ is the continuous linear functional
\[
\FT(\delta_0)(\phi) \, := \,
\delta_0(\FT(\phi)) \= \hat{\phi}(0) \=
\frac{1}{\sqrt{2\pi}} \, \int_{\,\R} \phi(x) \, dx,
\quad \phi \xeof \SR1\,,
\]
so $\FT(\delta_0)$ is nothing but the distribution induced by
the constant function $x \mapsto \frac{1}{\sqrt{2\pi}}$.

\paragraph{} %-- fundamental solutions
%------------------------------------------------------------------------------%
We now study the action of linear partial differential operators 
in distributional spaces.
First recall
that if $\psi \xeof \SF$ and $\disF \xeof \D'(\G)$,
then $\psi \disF$ defined by
\[
(\psi \disF) \phi \, := \, \disF (\psi \phi), \quad \quad \phi \xeof \TF
\]
also belongs to $\D'(\G)$. 
This implies that the linear differential operators $c_{\alpha} \, \PA$ 
with $c_{\alpha} \xeof \SF$ are acting on the space $\D'(\G)$,
since for $\disF \xeof \D'(\G)$
the mapping $c_{\alpha} \PA \disF$ is a distribution again.
By taking linear combinations of such operators,
we conclude that the action
of a linear partial differential operator
\[
\DO(x,u) \, := \,
\sum_{\abs{\alpha} \leq 2m} c_{\alpha}(x) \; \PA u(x)
\]
having smooth coefficients and acting on the space $\D'(\G)$ is well-defined,
which of course is especially true in case of $\G \eq \Rd$.

\paragraph{} %-- definition of fundamental solution
%------------------------------------------------------------------------------%
In the following we introduce the concept of \textit{fundamental solution}
for a linear partial differential operators $\DO$
having smooth coefficients $c_{\alpha}$.
It is nothing but a distribution $\Phi_{\DO} \xeof \D'(\Rd)$
that forms a solution of the differential equation
$\DO \Phi_{\DO} \eq \delta_0$,
where $\delta_0$ is the Dirac distribution {\wrt} the origin $x \eq 0$.

\subparagraph{}
It is useful to know a fundamental solution of $\DO$
as the following computation shows:
let $\disF \xeof \D'(\Rd)$ be any Schwartz distribution and consider the
inhomogeneous differential equation $\DO \, \disU \eq \disF$.
If convolution between $\Phi_{\DO}$ and $\disF$ exists,
then the equalities
\[
\DO(\Phi_{\DO} \ast \disF)  \=
(\DO \Phi_{\DO}) \ast \disF \, = \;
\delta_0 \ast \disF         \= \disF
\]
are valid,
so the convolutional product $\Phi_{\DO} \ast \disF$
solves the distributional differential equation $\DO \, \disU \eq \disF$. 
Notice that
there is a distributional solution $\disU$ of this equation
especially for all distributions $\disF$ having compact support.

\subparagraph{}
The fundamental solution of a linear differential operator
need not be unique,
which can be showed by the following example:
let $\disU \xeof \D'(\Rd)$ be a solution of the homogeneous equation
$\DO \, \disU \eq 0$,
then the distribution $\Phi_{\DO} + \disU\,$
is a fundamental solution of $\DO$, too.

\paragraph{} %- Malgrange-Ehrenpreis theorem
%------------------------------------------------------------------------------%
Regarding the existence of fundamental solutions for linear 
differential operators,
the \textsc{Malgrange-Ehrenpreis} \textit{theorem}
has been shown to be a milestone in PDE theory:

\begin{theorem1}\label{MALGRANGE-EHRENPREIS_THEOREM}
Each linear differential operator $\DO \neq 0$
with constant coefficients
has a fundamental solution in the distributional space $\D'(\Rd)$.
\end{theorem1}

\subparagraph{}
It should be mentioned 
that there are counterexamples which demonstrate
that this theorem is not valid for differential operators
whose coefficients are non-constant.

\paragraph{} %-- examples of fundamental solutions
%------------------------------------------------------------------------------%
In the following we describe the fundamental solutions
of three linear differential operators,
namely
the \textit{Laplace operator},
the \textit{heat operator} and
the \textit{wave operator}.

\begin{itemize}[leftmargin=4mm]
\item[1.]
The fundamental solution $\Phi_{\Delta}^d$ of the Laplacian $\Delta$
depends on the number of spatial coordinates.
Since $\Delta$ is invariant {\wrt} translations and rotations,
it is nearby to conjecture $\Phi_{\Delta}^d(x) \eq \phi_d(\abs{x} )$.
Indeed,
the functions 
\[
\Phi_{\Delta}(x) \=
\begin{cases}[1.5]
\quad \frac{1}{2} \abs{x}                         & (d=1)      \\
\quad - \frac{1}{2 \pi} \ln \abs{x}               & (d=2)      \\
\quad \frac{1}{(d-2) \, \omega_d } \abs{x} ^{2-d} & (d \geq 3) \\
\end{cases}
\]
are solutions of the differential equations
$\Delta \Phi_{\Delta}^d \eq \delta_0$ for $d \xeof \N$,
where $\omega_d$ is the Lebesgue-Borel measure of
the $d$-dimensional unit sphere.
Each other fundamental solution of the Laplacian is of the form
$\Phi_{\Delta}^d + P_1$,
where $P_1$ is any polynomial of degree less than two.

\item[2.]
The fundamental solution of the \textit{heat operator} 
$\DO := \partial_t - \Delta u$ in $\Rd$ is given 
by the so-called \textit{heat kernel} function
\[
\Phi_d(x,t) \=
\begin{cases}[1.5]
\quad \frac{1}{(4 \pi t)^{d/2}} \cdot e^{- \frac{\abs{x} ^2}{4t}},
         & \; x \xeof \Rd, \; t > 0,   \\
\quad 0, & \; x \xeof \Rd, \; t \leq 0.\\
\end{cases}
\]
Each other fundamental solution of the heat operator
adopts the form $\Phi_d + c$ ($c \xeof \C$),
so the kernel of the heat operator 
purely consists of constant distributions.

\item[3.]
When investigating the fundamental solutions
of the hyperbolic \textit{wave operator} $\Box := \partial_{tt} - \Delta u$,
one has to distinguish between even and odd spatial dimensions.

\subparagraph{}
Let $H$ be the \textit{Heaviside function},
defined by $H(t) := 0$ for $t<0$ and $H(t) := 1$ for $t \geq 0$.
Then we obtain for $d \leq 3$ 
\[
\Phi_{\Box}^d(x,t) \= \;
\begin{cases}[1.5]
\quad \frac{1}{2} H(t-\abs{x} )                                      & (d=1)\\
\quad \frac{1}{2 \pi} \, (t^2 - \abs{x} ^2)^{-1/2} \, H(t- \abs{x} ) & (d=2)\\
\quad \frac{1}{2 \pi} \, \delta(t^2 - \abs{x} ^2) \, H(t)            & (d=3).\\
\end{cases}
\]

\subparagraph{}
It is a notable fact
that all the higher fundamental solutions of the wave operator
can be derived from the functions $\Phi_{\Box}^1$ and $\Phi_{\Box}^2$.
Using the variable $r := \abs{x} $,
the functions $\Phi_{\Box}^{d+2}$ ($d \xeof \N$) are recursively given by
\[
\Phi_{\Box}^{d+2}(t,r) \=
\frac{-1}{2\pi r} \cdot \frac{\partial}{\partial_r} \Phi_{\Box}^d(t,r),
\quad d \geq 1.
\]
To determine the support of a fundamental solution $\Phi_{\Box}^d$,
we first define for fixed $x_0 \xeof \Rd$ the \textit{light cone}
\[
C^+(x_0) \, := \, \{(t,x): t > 0, \abs{x-x_0} < t \}
\]
and obtain the support of the fundamental solution according to
\[
\supp{\Phi_{\Box}^d} \=
\begin{cases}[1.5]
\quad \partial C^+(0),   \quad     d \mbox{ odd  or } d \geq 3,\\
\quad \overline{C^+(0)}, \quad\;\; d \mbox{ even or } d=1.\\
\end{cases}
\]
\end{itemize}

\paragraph{} %-- hypoelliptic differential operators
%------------------------------------------------------------------------------%
A linear differential operator $\DO$ is called \textit{hypoelliptic},
if {\wrt} the linear differential equation $\DO u \eq f$
the condition $f \xeof \SF$ implies $u \xeof \SF$.
Restricting to differential operators with constant coefficients,
hypoellipticity can be characterized
by means of a feature of the fundamental solution $\Phi_{\DO_c}$:
$\DO_c$ is hypoelliptic
iff
$\Phi_{\DO_c} \eq \disF_f$,
where $\disF_f$ is a regular distribution 
with generating function $f$ being indefinitely differentiable
in every $x \neq 0$.

\subparagraph{}
We next introduce the concept of \textit{singular support} $\Sigma(\disF)$
for $\disF \xeof \D'(\Rd)$:
it is the smallest subset $S \xeof \Rd$
such that $\disF$ remains regular on $\Rd \setminus S$.
Then hypoellipticity of linear differential operators
with constant coefficients are characterized as follows:

\begin{theorem1}\label{THEOREM_ON_HYPOELLIPTIC_OPERATORS}
A linear partial differential operator $\DO_c$ with constant coefficients
is hypoelliptic
\iff
$\,\Sigma(\DO_c \, \disF) \eq \Sigma(\disF)$ for all $\disF \xeof \D'(\Rd)$.
\end{theorem1}

\subparagraph{}
Knowing the fundamental solutions of the elliptic operator $\Delta$ and
the parabolic heat operator $\partial_t-\Delta$,
it turns out that both operators are hypoelliptic.
But theorem \ref{THEOREM_ON_HYPOELLIPTIC_OPERATORS} implies
that the hyperbolic wave operator $\Box \eq \partial_{tt}-\Delta$
is not hypoelliptic.

\paragraph{} %-- free space problems
%------------------------------------------------------------------------------%
To show the benefit of knowing the fundamental solution
of a linear partial differential operator with constant coefficients,
we consider free space problems ($\G \eq \Rd$)
related to the Laplacian, the heat and the wave operator.

\begin{itemize}[leftmargin=4mm]
\item[1.]
The \textit{Poisson equation} in the entire space $\Rd$ is given by
\begin{equation}\label{CLASSICAL_POISSON_EQUATION}
\Delta \, u \= - 4\pi \,f, \quad f \xeof C(\Rd).
\end{equation}

\subparagraph{} %- Poisson equation
A function $u \:\Rd \xto \R$ is called a \textit{classical solution}
of \eqref{CLASSICAL_POISSON_EQUATION},
if it is twice continuously differentiable
and fulfills the differential equation in each $x \xeof \Rd$.
However,
$\disU \xeof \D'(\Rd)$ is said to be a \textit{distributional solution}
of the Poisson equation,
if
\begin{equation}\label{DISTRIBUTIONAL_POISSON_EQUATION}
\Delta \, \disU \= - 4\pi \,\disF, \quad \disF \xeof \D'(\Rd)
\end{equation}
holds in the sense of distributions.

\subparagraph{}
Let $\Phi_{\Delta}$ be the fundamental solution of the Laplacian.
If the convolutional product $\Phi_{\Delta} \ast \disF$ exists,
which is especially true for each regular distribution $\disF_f$
with generating function $f \xeof C_c(\Rd)$,
then the solution of \eqref{DISTRIBUTIONAL_POISSON_EQUATION}
is obtained by convolution,
that is,
$\disU \eq \Phi_{\Delta} \ast \disF$.

\subparagraph{}
For the physically interesting case $d=3$
and for $f \xeof C_c(\R^3)$ we get a classical solution $u_f$,
called the \textit{Newton potential} for the function $f$:
\[
u_f(x) \=
\int_{\,\R^3} \, \frac{f(y)}{\abs{x-y } }  \, dy, \quad x \xeof \R^3.
\]

\item[2.]
By the \textit{heat equation} equation in $\Rd$ is meant
the inital value problem
\begin{equation}\label{CLASSICAL_HEAT_EQUATION}
\begin{cases}[1.2]
\quad u_t(t,x) - \Delta u(t,x) = f(t,x) & (x \xeof \Rd, \, t > 0) \\
\quad u(0,x) \= 0,                      & (x \xeof \Rd). \\
\end{cases}
\end{equation}

\paragraph{} %- heat equation
We first define the concept of distributional solution
for problem \eqref{CLASSICAL_HEAT_EQUATION}:
let $\disU_0 \xeof \D'(\Rd)$ and $\disF \xeof \D'(\R^{d+1})$ are
Schwartz distributions satisfying 
$\supp{\disF} \xsubset [0,\infty) \xtimes \Rd$,
then $\disU \xeof \D'(\R^{d+1})$ is called a \textit{distributional solution},
if $\supp{\disU} \xsubset [0,\infty) \xtimes \Rd$ and
\begin{equation}\label{DISTRIBUTIONAL_HEAT_EQUATION}
\disU_t \, - \Delta \, \disU \= \disF + \disU_0 \otimes \delta_0,
\end{equation}
where $\delta_0 \xeof \D'(\R)$ is the Dirac distribution.
Here the initial condition of the problem has been transformed
into an inhomogeneous part of the differential equation.

\subparagraph{}
The existence of a solution for \eqref{DISTRIBUTIONAL_HEAT_EQUATION}
is ensured,
if $\disU_0 \xeof \E'(\Rd)$ and $\disF \xeof \E'(\R^{d+1})$.
In this case the solution $\disU \xeof \D'(\Rd)$ is given by
convolution of the fundamental equation $\Phi_d$ 
with the inhomogeneous parts of \eqref{DISTRIBUTIONAL_HEAT_EQUATION},
\[
\disU \=
\Phi_d \ast (\disF + \disU_0 \otimes \delta_0) \=
\Phi_d \ast \disF \, + \, \Phi_d \ast \disU_0.
\]

\subparagraph{}
Given a classical solution $u$ of \eqref{CLASSICAL_HEAT_EQUATION},
the related Schwartz distribution $\disF_u$ is a solution
in the distributional sense.
Conversely,
let $\disF_u$ be regular and solve \eqref{DISTRIBUTIONAL_HEAT_EQUATION},
then $u$ forms a classical solution,
if some regularity conditions are valid:

\begin{itemize}[leftmargin=10mm]
\item[a.]
$\disU, \disU_0$ and $\disF$ are regular distributions
with associated functions $u, u_0, f \xeof \Lc$,

\item[b.]
$u_0$ is continuous on $\Rd$ and
$f$ is continuous on $[0,\infty) \xtimes \Rd$,

\item[c.]
$u$ is twice continuously differentiable on $[0,\infty) \xtimes \Rd$.
\end{itemize}

\item[3.]
The initial-value problem for the \textit{wave equation} in $\Rd$
is given by

\begin{equation}\label{CLASSICAL_WAVE_EQUATION}
\begin{cases}[1.2]
\quad u_{tt}(t,x) - \Delta u(t,x) \= f(t,x) & (x \xeof \Rd, \, t > 0) \\
\quad u(0,x) = u_0(x)                       & (x \xeof \Rd) \\
\quad u_t(0,x) = u_1(x)                     & (x \xeof \Rd. \\
\end{cases}
\end{equation}

\paragraph{} %- wave equation
As for the heat equation,
we define for \eqref{CLASSICAL_WAVE_EQUATION}
the concept of distributional solution,
where again both initial value conditions are transformed
into inhomogeneous parts of the equation by means of tensor products.

\subparagraph{}
Let $\disU_0 \xeof \D'(\Rd)$, $\disU_1 \xeof \D'(\Rd)$
and $\disF \xeof \D'(\R^{d+1})$ be distributions,
then $\disU$ is called a distributional solution
of \eqref{CLASSICAL_WAVE_EQUATION},
if
\begin{equation}\label{DISTRIBUTIONAL_WAVE_EQUATION}
\disU_{tt} - \Delta \disU \=
\disF + \disU_0 \otimes \delta_0' \, + \, \disU_1 \otimes \delta_0 \,,
\end{equation}
where $\delta_0$ is the Dirac distribution for $t \eq 0$ 
and $\delta_0'$ is the distributional derivative of $\delta_0$
{\wrt} the variable $t$.

\subparagraph{}
Since the support of the function
$\disU_0 \, \otimes \, \delta_0' \, + \, \disU_1 \, \otimes \, \delta_0$
belongs to the set $\R^{d+1}$,
there is to each distribution $\disF \xeof \D'(\R^{d+1})$ with support
\[
\supp{\disF} \= \{\, (x,t) \xeof \Rd \xtimes \R: t \geq 0 \,\}
\]
a further distribution
$
\disU_F :=
\Phi_{\Box}^d \ast (\disF + \disU_0 \otimes \delta'_0
+ \disU_1 \otimes \delta_0)
\in \D'(\R^{d+1}),
$
which in fact solves the distributional differential equation
\eqref{DISTRIBUTIONAL_WAVE_EQUATION}.

\subparagraph{}
Each classical solution of \eqref{CLASSICAL_WAVE_EQUATION}
is also a distributional solution.
However,
the reverse statement only holds under the conditions 

\begin{itemize}[leftmargin=9mm]
\item[a.]
$\disU, \disU_0, \disU_1$ and $\disF$ are regular
and generated by ordinary functions $u, u_0, u_1, f \xeof \Lc$,

\item[b.]
$u_0 \xeof C^1(\Rd)$ and $u_1 \xeof C(\Rd)$,

\item[c.]
$f \xeof C([0,\infty) \xtimes \Rd)$ and

\item[d.]
$u \xeof C^2((0,\infty) \xtimes \Rd) \cap C([0,\infty) \xtimes \Rd)$.
\end{itemize}
\end{itemize}

\paragraph{} %-- applications of the Fourier transform
%------------------------------------------------------------------------------%
The usability of Fourier transformation mainly relies on the fact
that elementary differential operators $\partial_k$,
applied to Schwartz functions $\phi$,
appear as multiplication operators to the Fourier transform $\FT(\phi)$
and via verce,
that is, 
we always have
\[
\partial_k \, \phi(x) \, \overset{\FT}{\longmapsto}
\, i \, \xi_k \cdot \hat{\phi}(\xi)\,,
\quad \quad
i \, x_k \phi(x) \, \overset{\FT}{\longmapsto}
\, \partial_k \, \hat{\phi}(\xi)\,,
\]
for $k=1,\ldots,d$,
as one can verify by partial integration.
For any linear differential operator $\DO_c$
with constant coefficients $c_{\alpha}$
and symbol $S_{\DO_c}$ we accordingly obtain 
\[
(\DO_c \, \phi)(x) \, \overset{\FT}{\longmapsto} \,
S_{\DO_c}(i \xi) \cdot \hat{\phi}(\xi).
\]
This is nothing but the fact
that all the differential operators $\DO_c$ are Fourier multipliers.
They operate on functions $\phi \xeof \SRd$ like multi\-plications
with the symbols $S_{\DO_c}$ on the Fourier trans\-form $\FT \phi$.
Hence a linear differential operator $\DO_c$ with constant coefficients
is conjugated to its symbol function $S_{\DO_c}$ in the sense of
\[
\DO_c \phi \=
(\FT^{-1} \circ S_{\DO_c} \circ \FT) \, \phi \=
\FT^{-1}(S_{\DO_c} \hat{\phi})
\qfa \phi \xeof \SRd.
\]

\paragraph{} %-- the general usage of the Fourier transform in PDE theory
%------------------------------------------------------------------------------%
We want incident
how Fourier transformation can be used
to solve linear differential equations in the whole space $\Rd$
and governed by the Laplacian.
The basic idea is to transform the differential equation
either into a purely algebraic equation
or into an ordinary differential equation.
Once these are solved,
inverse Fourier transformation is applied
to get the solution of the original problem.

\subparagraph{}
Two examples for these methods will be presented.
In the first one,
the bijectivity of Fourier transformation on the Schwartz space
is used to solve a specific differential equation,
while in the second one the simplified computation of convolution
products after performing Fourier transformation is helpful
to find fundamental solutions of the underlying differential operators.

\begin{itemize}[leftmargin=4mm]
\item[1.]
The equation $\DO_c \, u \eq f$ with $f \xeof \SRd$
and governed by a linear differential operator $\DO_c$
with constant coefficents
can be converted by applying Fourier transformation
into the algebraic equation
\begin{equation}\label{ALGEBRAIC_EQUATION_OF_LDE}
S_{\DO_c}(i \xi) \, \hat{u}(\xi) \= \hat{f}(\xi)
\quad \quad (\xi \xeof \Rd).
\end{equation}
Now,
if the symbol function $S_{\DO_c}$ of $\DO_c$ owns the property
\begin{equation}\label{SYMBOL_CONDITION}
S_{\DO_c}(i \xi) \neq 0
\qfa \xi = (\xi_1,\ldots,\xi_d) \xeof \Rd,
\end{equation}
equation \eqref{ALGEBRAIC_EQUATION_OF_LDE} can be resolved for $\hat{u}(\xi)$,
and the solution $u$ of $\DO_c \,u \eq f$
is obtained by applying inverse Fourier transformation:
\[
\hat{u}(\xi) = \hat{f}(\xi) \, / \, S_{\DO_c}(i \xi)
\]
and
\[
u(x) \= \frac{1}{(2 \pi)^{d/2}} \,
\int_{\,\Rd}
e^{i \xi x} \cdot S_{\DO_c}^{-1}(i \xi) \cdot \hat{f}(\xi) \, d \xi.
\]

\item[2.]
If \eqref{SYMBOL_CONDITION} holds,
the Fourier transform method can also be applied
to solve the distributional differential equation
$\DO_c \, \disU \eq \disF$ with $\disF \xeof \TD$.
Since $\DO_c$ is bijective on $\SRd$ and 
\[
P_{\DO_c}(i \xi) \= P_{\DO_c^*}(-\xi)
\qfa \xi \xeof \Rd
\]
holds,
the formally adjoint operator $\DO_c^*$
is bijective on $\SRd$ and so its transposed operator
$\DO_c \eq (\DO_c^*)^t$ is bijective on the dual space $\TD$.
Thus,
if \eqref{SYMBOL_CONDITION} is valid,
the equation $\DO_c \, \disU \eq \disF$
is uniquely solvable for each tempered distribution $\disF$.

\subparagraph{}
Especially for the differential operator 
$\Delta{_\lambda} := -\Delta + \lambda$ and $\lambda > 0$ we obtain
\[
S_{\Delta_{\lambda}}(i \xi) = \abs{\xi} ^2 + \lambda \neq 0
\]
for $\xi \xeof \Rd$,
hence the solution $u$ of $\Delta_{\lambda} u \eq f$ with $f \xeof \SRd$ is
\[
u(x) \=
\frac{1}{(2 \pi)^{d/2}} \,
\int_{\,\Rd} e^{i \xi x}
\cdot \frac{\hat{f}(\xi)}{\abs{\xi} ^2 + \lambda} \, d \xi, \quad (x \xeof \Rd).
\]
\end{itemize}

\paragraph{} %-- partial Fourier transformation
%------------------------------------------------------------------------------%
Partial Fourier transformation is a method
where the operator $\FT$ is applied only to the spatial variables
of the partial differential equation.
The result is an ordinary differential equation
{\wrt} the time variable $t$.
Then solving this equation and applying inverse Fourier transformation
leads to the solution of the original problem.

\paragraph{} %-- examples of partial Fourier tranformation
%------------------------------------------------------------------------------%
In the rest of this section we shall sketch
how {partial Fourier transformation} is useful
to solve the initial-value problem of the heat and wave equation
in the entire space $\Rd$.

\begin{itemize}[leftmargin=4mm]
\item[1.]
The initial-value problem of the heat equation is given by
\[
\begin{cases}[1.2]
\quad u_t - \Delta u = 0   & (t > 0)     \\
\quad u(0,x) \= f(x)       & (x \xeof \R). \\
\end{cases}
\]
Applying Fourier transformation to both sides of the differential equation
{\wrt} the spatial variables $x_1,\ldots,x_d$,
we obtain
\begin{equation}\label{PARTIAL_FOURIER_TRANSFORMATION_I}
\frac{\partial}{\partial t} \hat{u}(t,\xi) \, + \, \xi^2 \, \hat{u}(t,\xi)
\= 0,
\quad \quad \quad
\hat{u}(0,\xi) = \hat{f}(\xi),
\end{equation}
where $\xi \eq (\xi_1,\ldots,\xi_d)$ are the variables
in the \textit{phase space}.

\subparagraph{}
Then multiplying both sides of the first equation
in \eqref{PARTIAL_FOURIER_TRANSFORMATION_I} by $e^{\abs{\xi} ^2 t}$ leads to
\[
\frac{\partial}{\partial t} \, e^{\abs{\xi} ^2 t} \, \hat{u}(t,\xi) \= 0
\]
with solution
$\hat{u}(t,\xi) \eq \phi(\xi) \, e^{- \abs{\xi} ^2 t}$.
The mapping $\xi \mapsto \phi(\xi)$ is
the integration constant for each fixed point $\xi$.
The second equation in \eqref{PARTIAL_FOURIER_TRANSFORMATION_I}
implies $\phi(\xi)=\hat{f}(\xi)$,
hence the Fourier trans\-formed solution $\hat{u}$ 
of the initial-value problem is nothing but
\begin{equation}\label{PARTIAL_FOURIER_TRANSFORMATION_II}
\hat{u}(t,\xi) \eq \hat{f}(\xi) \cdot e^{- \abs{\xi} ^2 t}.
\end{equation}

\subparagraph{}
The solution of the original equation is obtained
by applying inverse Fourier transformation and the rule
$
\FT^{-1}(f \cdot g) \eq \FT^{-1}(f) \ast \FT^{-1}(g)
$
to each side of \eqref{PARTIAL_FOURIER_TRANSFORMATION_II}:
\[
u(t,x) \=
\FT^{-1}(\hat{f}(\xi) \cdot e^{- \abs{\xi} ^2 t}) \=
\FT^{-1}(\hat{f}(\xi)) * \FT^{-1}(e^{- \abs{\xi} ^2 t})
\]
\[
\= f(x) * \FT^{-1}(e^{- \abs{\xi} ^2 t})
\=
\frac{1}{\sqrt{4\pi t}}
\int_{\,\R} e^{-\abs{x-y} ^2/(4t)} \,f(y)\,dy.
\]
The function
\[
H_t(x,y) \, := \,
\frac{1}{\sqrt{4\pi t}} \, e^{- \abs{x-y} ^2/(4t)}
\=
\Phi_{\partial_t - \Delta}(t,x-y)
\]
is called a \textit{heat kernel},
since it constitutes the kernel of the integral operator
$\cal{H}_t:f \mapsto H_t \ast f$.

\item[2.]
The method also enables to solve the initial-value problem
for the wave equation
\[
\begin{cases}[1.2]
\quad u_{tt} - \, \Delta u \= 0  & (t > 0),    \\
\quad u(0,x) \= f(x)             & (x \xeof \Rd),\\
\quad u_t(0,x) \= g(x)           & (x \xeof \Rd).\\
\end{cases}
\]

\subparagraph{}
Partial Fourier transformation let become this equation into
\[
\hat{u}_{tt}(t,\xi) \= \abs{\xi} ^2 \hat{u}(t,\xi),
\quad \quad
\hat{u}(0,\xi) \= \hat{f}(\xi),
\quad \quad
\hat{u_t}(0,\xi) \= \hat{g}(\xi),
\]
and the general solution of the ordinary differential equation
in phase space is 
\[
\hat{u}(t,\xi)
\=
\hat{f}(\xi) \, \cos( t \abs{\xi} )
\, + \,
\hat{g}(\xi) \: \frac{\sin(t \abs{\xi} )}{\abs{\xi} }.
\]

\subparagraph{}
Unfortunately,
performing inverse Fourier transformation to $\hat{u}$
for any $n >1$ is not an easy task.
But for $n=1$ we obtain
\[
u_1(x,t) \=
\FT^{-1}\left(\hat{f}(\xi) \cdot \cos(\abs{\xi} ) \right)
\=
\frac{1}{2}\,[f(x+t) \, + \, f(x-t)],
\]
\[
u_2(x,t) \=
\FT^{-1}\left(\hat{g}(\xi) \cdot
\frac{\sin(\abs{\xi} )}{\abs{\xi} } \right)
\=
\frac{1}{2}\, \int_{-t}^{+t} g(x+s) \, ds,
\]
thus the solution of the wave equation is given by $u \eq u_1+u_2$,
known as the \textsc{d'Alembert} \textit{solution}
\[
u(t,x) \=
\frac{1}{2} \, [f(x+t) \, + \, f(x-t)]
\, + \,
\frac{1}{2t}\, \int_{-t}^{+t} g(x+s) \, ds.
\]
\end{itemize}

\newpage
\thispagestyle{empty}
\fancyhead[C] {}
\fancyhead[OL]{}
\fancyhead[ER]{}
%------------------------------------------------------------------------------%
\section*{Appendix A}                                                          %
%------------------------------------------------------------------------------%
The following figure shows the inheritance tree
of topological and normed linear spaces,
where the Banach and Hilbert spaces are positioned
at the very end of this hierarchy:

\begin{center}
\scalebox{0.695}{

\begin{tikzpicture}[node distance=4.0cm,auto]

\node[align=center,rectangle,draw](TLS) at (0,-1.0)
{\;topological linear space\; \\ $(\topV,\topT_V)$};

\node[align=center,rectangle,draw](LCTLS) at (-5,-3.5)
{locally convex \\ \;topological linear space\; \\ $(\topV,\{ p_{\alpha} \}_{\alpha \in \cal{I}})$};

\node[align=center,rectangle,draw](SCTLS) at (5,-3.5)
{\textbf{sequentially complete}\\ \; \textbf{topological linear space} \; \\ $(V,\topT_V)$};

\node[align=center,rectangle,draw](MS) at (-3,-5.9)
{metric space \\ $(\metrS,\topT_d)$};

\node[align=center,thick,rectangle,draw](CMS) at (1,-5.9)
{\textbf{complete} \\ \; \textbf{metric Space} \; \\ $(\metrS,\topT_d)$};

\node[align=center,thick,rectangle,draw](F-S) at (5,-5.9)
{\; \textbf{F-space}\; \\ $(F,\topT_V=\topT_d)$};

\node[align=center,rectangle,draw](PFS) at (-5,-8.0)
{\;pre-{\FRECHET} space\; \\ $(F,\{p_n\}_{n \in \N})$};

\node[align=center,thick,rectangle,draw](FS) at (5,-8.0)
{\;\textbf{{\FRECHET} space}\; \\ $(F,\{p_n\}_{n \in \N})$};

\node[align=center,rectangle,draw](NS) at (-5,-10.5)
{\;{normed linear space}\; \\$(\banX,\norm{\cdot})$};

\node[align=center,thick,rectangle,draw](BS) at (5,-10.5)
{\;\textbf{Banach space}\; \\$(\banX,\norm{\cdot})$};

\node[align=center,thick,rectangle,draw](RBS) at (5,-12.0)
{\;\textbf{reflexive Banach space}\; \\$(\banX,\norm{\cdot})$};

\node[align=center,thick,rectangle,draw](UCS) at (5,-13.5)
{\;\textbf{uniformly convex space}\; \\$(\banX,\norm{\cdot})$};

\node[align=center,rectangle,draw](IPS) at (-5,-15)
{\;{inner product space}\; \\$(\vecH,\scalar{\cdot}{\cdot})$};

\node[align=center,thick,rectangle,draw](HS) at (5,-15)
{\;\textbf{Hilbert space}\; \\$(\hilH,\scalar{\cdot}{\cdot})$};

\draw[->,>={Stealth}, thick] (MS) to (CMS);
\draw[->,>={Stealth}, thick] (MS) to (PFS);
\draw[->,>={Stealth}, thick] (CMS) to (F-S);
\draw[->,>={Stealth}, thick] (TLS) to (SCTLS);
\draw[->,>={Stealth}, thick] (TLS) to (LCTLS);
\draw[->,>={Stealth}, thick] (SCTLS) to (F-S);
\draw[->,>={Stealth}, thick] (F-S) to (FS);
\draw[->,>={Stealth}, thick] (LCTLS) to (PFS);
\draw[->,>={Stealth}, thick] (PFS) to (NS);
\draw[->,>={Stealth}, thick] (NS) to (IPS);
\draw[->,>={Stealth}, thick] (PFS) to (FS);
\draw[->,>={Stealth}, thick] (NS) to (BS);
\draw[->,>={Stealth}, thick] (IPS) to (HS);
\draw[->,>={Stealth}, thick] (FS) to (BS);
\draw[->,>={Stealth}, thick] (BS) to (RBS);
\draw[->,>={Stealth}, thick] (RBS) to (UCS);
\draw[->,>={Stealth}, thick] (UCS) to (HS);

\node[align=center](I) at (0,0.4) {\textit{metric and topological linear spaces}};
\node[align=center](II) at (   0,-16.4) {\textit{normed linear spaces}};
\draw[dashed,thick]        (-7.8,  0.0) rectangle (8.0,-9.0);
\draw[dashed,thick]        (-7.8, -9.5) rectangle (8.0,-16.0);
\end{tikzpicture}
}
\end{center}

\begin{center}
\footnotesize{Figure A.1:
              The hierarchy of metric and topological linear spaces}
\end{center}
%------------------------------------------------------------------------------%
%                                 BIBLIOGRAPHY                                 %
%------------------------------------------------------------------------------%
\newpage
\thispagestyle{empty}
\fancyhead[C] {}
\fancyhead[OL]{}
\fancyhead[ER]{}

\titleformat{\section}
{\normalfont \bfseries \filcenter}
{\thesection.\;}{0mm}{}

\newlength{\bibitemsep}\setlength{\bibitemsep}{.5\baselineskip plus .05\baselineskip minus .05\baselineskip}
\newlength{\bibparskip}\setlength{\bibparskip}{0pt}
\let\oldthebibliography\thebibliography
\renewcommand\thebibliography[1]{%
  \oldthebibliography{#1}%
  \setlength{\parskip}{\bibitemsep}%
  \setlength{\itemsep}{\bibparskip}%
}

\vspace{8mm}
\begin{thebibliography}{99}
%\setlength\bibitemsep{3\itemsep}
\begin{small}

\bibitem{AT} W.\,Arendt, A.F.M. ter Elst:
\textit{From Forms to Semigroups}.
Operator Theory: Advances and Applications, Vol.\,221, pp.\,47-69.

\bibitem{Br0} F.\,Browder:
\textit{Estimates and Existence Theorems for Elliptic Boundary Value Problems}.
Proc NatI Acad Sci USA, Vol.\,45, pp.\,365-372, 1959.

\bibitem{Br1} F.\,Browder:
\textit{Functional Analysis and Partial Differential Equations I}.
Math. Annalen 138, pp.\,55-79, 1959.

\bibitem{Br2} F.\,Browder:
\textit{On the Spectral Theory of Elliptic Differential Operators I}.
Math. Annalen 142, pp.\,22-130, 1961.

\bibitem{Br3} F.\,Browder:
\textit{Functional Analysis and Partial Differential Equations II}.
Math. Annalen 145, pp.\,81-226, 1962.

\bibitem{DU}  J.\,Diestel, J.\,J.\,Uhl:
\textit{Vector Measures}.
American Mathematical Soc., 1977.

\bibitem{EE}  D.\,E.\,Edmunds, W.\,D.\,Evans:
\textit{Spectral theory and differential operators}.
Oxford University Press, 1987.

\bibitem{EN}  K.-J.\,Engel, R.\,Nagel:
\textit{A Short Course on Operator Semigroups}.
Universitext, Springer, 2006.

\bibitem{Ev}  L.\,Evans: \textit{Partial Differential Equations}.
AMS, Graduate Studies in Mathematics 13, 1988.

\bibitem{GT}  D.\,Gilbarg, N.S.\,Trudinger:
\textit{Elliptic Partial Differential Operators of Second Order}.
Springer, 2001.

\bibitem{Go}  J.\,Goldstein:
\textit{Semigroups of Linear Operators and Applications}.
Oxford University Press, 1985.

\bibitem{Ha}  M.\,Haase:
\textit{The Functional Calculus for Sectorial Operators}.
Birk\-h\"{a}user, 2006.

\bibitem{He}  D.\,Henry:
\textit{Geometric Theory of Semilinear Parabolic Equations}.
Lecture Notes in Mathematics Vol.\,840, Springer, 1981.

\bibitem{HP}  E.\,Hille, R.\,S.\,Phillips:
\textit{Functional Analysis and Semigroups}.
AMS Vol.\,31, 1948.

\bibitem{Ka}  T.\,Kato:
\textit{{Perturbation theory of linear operators}}.
Springer, Grund\-lehren der mathematischen Wissenschaften, Band 132, 1966.

\bibitem{Kr}  S.\,Krein:
\textit{Linear Equations in Banach Spaces}.
Birk\-h\"{a}user, 1982.

\bibitem{Li2} J.-L.\,Lions:
\textit{\'{E}quations Diff\'{e}rentielles Op\'{e}rationelles et Probl\`{e}mes aux Limites}.
Springer,
Grund\-lehren der mathematischen Wissenschaften, Band 111, 1961.

\bibitem{LM1} J.-L.\,Lions, E.\,Magenes:
\textit{Non-Homogeneous Boundary Value Problems with Applications I}.
Springer, Grund\-lehren der mathematischen Wissenschaften, Band 181, 1972.

\bibitem{Lu}  A.\,Lunardi:
\textit{Analytic Semigroups and Optimal Regularity in Parabolic Problems}.
Birk\-h\"{a}user, 1995.

\bibitem{Mi}  M.\,Miclav\^{c}i\^{c}:
\textit{Applied Functional Analysis and Partial Differential Equations}.
World Scientific Publishing, 1998.

\bibitem{Pa}  A.\,Pazy:
\textit{Semigroups of Linear Operators
        and Applications to Partial Differential Equations}.
Applied Mathematical Sciences Vol.\,44, Springer, 1983.

\bibitem{RS1} M.\,Reed, B.\,Simon:
\textit{Modern Methods of Mathematical Physics I: Functional Analysis}.
Academic Press Inc., 1980.

\bibitem{RS4} M.\,Reed, B.\,Simon:
\textit{Modern Methods of Mathematical Physics VI: Spectral Analysis}.
Academic Press Inc., 1980.

\bibitem{RR}  M.\,Renardy, R.\,Rogers:
\textit{An Introduction to Partial Differential Equations}.
Texts in Applied Mathematics 13, Springer, 2004.

\bibitem{Wl}  J.\,Wloka:
\textit{Partial Differential Equations}.
Cambridge University Press, 1987.

\bibitem{Yo}  K.\,Yosida:
\textit{Functional Analysis}.
Springer, 1980.
\end{small}
\end{thebibliography}

%------------------------------------------------------------------------------%
%                               END OF DOCUMENT                                %
%------------------------------------------------------------------------------%
\end{document}
